% Piles électrochimiques et corrosion — Chimie Tle TI — Cours signé OnBuch+
% Programme officiel MINESEC, Terminale TI. Compiler avec : tectonic L10-piles-electrochimiques-corrosion.tex
\newcommand{\DOCMATIERE}{Chimie}
\newcommand{\DOCNIVEAU}{Tle TI}
\newcommand{\DOCMODULE}{Module 3 · Oxydoréduction}
\newcommand{\DOCLECON}{10}
\newcommand{\DOCTITRE}{Piles électrochimiques et corrosion}
% OnBuch+ — préambule « cours signés » ENRICHI (compilé avec tectonic / XeLaTeX)
% Système visuel « Pop » d'OnBuch+ : fond crème (jamais blanc), bordures encre
% épaisses, ombres « dures » décalées, titres Archivo Black en capitales.
% Version enrichie pour les cours de chimie : chimie (mhchem, chemfig),
% graphes (pgfplots), boîtes pédagogiques supplémentaires (définition,
% propriété, méthode, attention, le savais-tu, exercice résolu, exercice,
% TP, bilan), page de garde refaite et signature OnBuch+ ancrée en bas.
%
% Macros attendues dans main.tex AVANT \input{preamble} :
%   \DOCMATIERE  \DOCNIVEAU  \DOCMODULE  \DOCLECON (numéro)  \DOCTITRE
\documentclass[11pt]{article}

\usepackage{fontspec}
\usepackage[french]{babel}
\usepackage[a4paper,margin=2cm,top=3.4cm,bottom=2.8cm,headheight=1.4cm,footskip=1.3cm]{geometry}
\usepackage{amsmath,amssymb,amsfonts}
\usepackage[table]{xcolor} % [table] : fournit aussi \rowcolors (alternance de lignes)
\usepackage{pagecolor}
\usepackage{tikz}
\usetikzlibrary{arrows.meta,positioning,calc,shapes.geometric,shapes.misc,shapes.arrows,decorations.pathmorphing,decorations.markings,patterns,shadows,mindmap,trees,fit,backgrounds,matrix,intersections}
\usepackage{pgfplots}
\pgfplotsset{compat=1.18}
\usepackage[version=4]{mhchem}
\usepackage{chemfig}
\usepackage{enumitem}
\usepackage{fancyhdr}
% [most] charge toutes les bibliothèques tcolorbox (listings, minted, raster,
% documentation...) et gonfle inutilement la mémoire TeX sur un document long
% (~300 boîtes) : on ne charge que ce qui est réellement utilisé.
\usepackage[skins,breakable]{tcolorbox}
\usepackage{graphicx}
\usepackage{colortbl}
\usepackage{booktabs}
\usepackage{tabularx}
\usepackage{array}
\usepackage{multicol}
\usepackage{siunitx}
% parse-numbers=false : accepte aussi \SI{2,5 \times 10^{-3}}{...}
\sisetup{locale=FR,output-decimal-marker={,},group-minimum-digits=4,parse-numbers=false}
\DeclareSIUnit\ohm{\ensuremath{\symup{\Omega}}} % U+2126 absent de la police
\usepackage{qrcode}
% \degree : commande fréquemment utilisée par les modèles pour un angle ou une
% température en dehors de \SI{}{} (non fournie par les paquets chargés).
\providecommand{\degree}{^{\circ}}
\usepackage{lastpage}
\usepackage{hyperref}
\hypersetup{hidelinks}

% ============================================================
% Palette « Pop » OnBuch+ — mêmes hex que la classe P (plus_ui.dart)
% ============================================================
\definecolor{poporange}{HTML}{F59321}
\definecolor{popdark}{HTML}{E07A0C}
\definecolor{popcream}{HTML}{FFF6E8}
\definecolor{popink}{HTML}{1C1714}
\definecolor{poppurple}{HTML}{7A5AE0}
\definecolor{popgreen}{HTML}{2E9E6A}
\definecolor{poppink}{HTML}{E0457B}
\definecolor{popblue}{HTML}{2F7DE1}
\definecolor{popgold}{HTML}{C98A12}
\definecolor{popmuted}{HTML}{8A7F76}
\definecolor{popline}{HTML}{EADFD2}
\definecolor{popgray}{HTML}{8A7F76}
\colorlet{popgrey}{popgray}
% Alias tolérants (le modèle nomme parfois les couleurs différemment)
\colorlet{popwhite}{white}
\colorlet{popblack}{popink}
\colorlet{popred}{poppink}
\colorlet{popyellow}{popgold}
% Teintes claires pour fonds de tableaux / zones
\colorlet{popblueL}{popblue!10!white}
\colorlet{poporangeL}{poporange!14!white}
\colorlet{popgreenL}{popgreen!12!white}
\colorlet{poppurpleL}{poppurple!12!white}
\colorlet{poppinkL}{poppink!12!white}
\colorlet{popgoldL}{popgold!14!white}

\pagecolor{popcream}

% ============================================================
% Typographie
% ============================================================
\defaultfontfeatures{Ligatures=TeX}
\setmainfont{PlusJakartaSans-Medium.ttf}[Path=../../../fonts/,BoldFont=PlusJakartaSans-Bold.ttf]
% Maths en sans-serif (Fira Math) avec chiffres et lettres droites de Plus
% Jakarta, pour que les formules se fondent dans le texte.
\usepackage[math-style=ISO,bold-style=ISO]{unicode-math}
\setmathfont{FiraMath-Regular.otf}[Path=../../../fonts/]
\setmathfont{PlusJakartaSans-Medium.ttf}[Path=../../../fonts/,range={up/num,up/{Latin,latin},\mathup}]
\usepackage{icomma} % virgule décimale sans espace en mode math
% Caractères absents de la police de texte (sinon carré vide) : rendus par
% leur équivalent mathématique, en texte comme en formule.
\usepackage{newunicodechar}
\newunicodechar{α}{\ensuremath{\alpha}}
\newunicodechar{Α}{\ensuremath{\alpha}}
\newunicodechar{β}{\ensuremath{\beta}}
\newunicodechar{δ}{\ensuremath{\delta}}
\newunicodechar{✓}{\popcheck{popgreen}}
\newfontfamily\popdisplay{ArchivoBlack-Regular.ttf}[Path=../../../fonts/]
\newfontfamily\poplogo{SpaceGrotesk-Bold.ttf}[Path=../../../fonts/]
\newfontfamily\popsemi{PlusJakartaSans-SemiBold.ttf}[Path=../../../fonts/]
\newfontfamily\popxbold{PlusJakartaSans-ExtraBold.ttf}[Path=../../../fonts/]
\newfontfamily\popxboldls{PlusJakartaSans-ExtraBold.ttf}[Path=../../../fonts/,LetterSpace=3.0]

\setlist[enumerate]{leftmargin=1.7em,itemsep=5pt,topsep=4pt}
\setlist[itemize]{leftmargin=1.7em,itemsep=4pt,topsep=4pt,label={\color{poporange}\textbullet}}
\setlength{\parindent}{0pt}
\setlength{\parskip}{5pt}
\renewcommand{\arraystretch}{1.25}

% chemfig : liaisons un peu plus épaisses, encre Pop
\setchemfig{atom sep=2em,bond style={line width=0.9pt,popink},cram width=3pt}

% pgfplots : style Pop par défaut
\pgfplotsset{
  every axis/.append style={axis line style={popink,line width=1pt},tick style={popink},
    grid style={popline,line width=0.5pt},label style={font=\small},tick label style={font=\footnotesize}},
  popcurve/.style={poporange,line width=1.8pt,no markers,smooth},
}

% ============================================================
% Pill / squiggle
% ============================================================
\newcommand{\poppill}[3]{%
  \tikz[baseline=-0.55ex]\node[draw=popink,line width=1.2pt,rounded corners=2.6pt,%
    fill=#2,inner xsep=6.5pt,inner ysep=3pt]{%
    \popxboldls\fontsize{7}{7}\selectfont\color{#3}\MakeUppercase{#1}};%
}
\newcommand{\popsquiggle}[1]{%
  \tikz[baseline=-0.4ex]{
    \draw[#1,line width=2.6pt,line cap=round]
      plot [smooth] coordinates {(0,0.09) (0.34,0.01) (0.68,0.21) (1.02,0.01) (1.36,0.10)};
  }%
}
% Mot-clé surligné (marqueur orange sous le texte)
\newcommand{\cle}[1]{{\popxbold\color{popdark}#1}}

% ============================================================
% En-tête / pied
% ============================================================
\pagestyle{fancy}
\fancyhf{}
\renewcommand{\headrulewidth}{0pt}
\renewcommand{\footrulewidth}{0pt}
\lhead{\raisebox{-0.15ex}{\poplogo\fontsize{13}{13}\selectfont\color{popink}OnBuch\color{poporange}+}}
\rhead{\poppill{\DOCMATIERE\ \textbullet\ \DOCNIVEAU\ \textbullet\ Leçon \DOCLECON}{white}{popink}}
\lfoot{\popxbold\fontsize{7.6}{7.6}\selectfont\color{popmuted}\MakeUppercase{Cours signé OnBuch+}\ \textbullet\ {\popsemi\DOCTITRE}}
\rfoot{\tikz[baseline=-0.6ex]\node[rounded corners=8.5pt,fill=popink,minimum height=17pt,minimum width=17pt,inner xsep=5pt,inner sep=0]{\popxbold\fontsize{8}{8}\selectfont\color{popcream}\thepage\,/\,\pageref*{LastPage}};}

% ============================================================
% Titres
% ============================================================
\newcommand{\chaptitre}[2]{%
  \begin{tcolorbox}[enhanced,colback=poporange,colframe=popink,boxrule=2.6pt,arc=11pt,
      drop shadow=popink,left=15pt,right=15pt,top=12pt,bottom=11pt,before skip=3pt,after skip=18pt]
    \poppill{#2}{popink}{popcream}\par\vspace{9pt}
    {\raggedright\hyphenpenalty=10000\exhyphenpenalty=10000\popdisplay\fontsize{22}{24}\selectfont\color{popcream}\MakeUppercase{#1}\par}
  \end{tcolorbox}
}
% Section numérotée : pastille orange avec le numéro + titre Archivo
\newcounter{popsec}
\newcommand{\coursec}[1]{%
  \stepcounter{popsec}\setcounter{popsub}{0}%
  \par\vspace{16pt}\noindent
  \begin{minipage}{\linewidth}
  \tikz[baseline=-0.7ex]\node[draw=popink,line width=1.6pt,fill=poporange,rounded corners=4pt,minimum size=22pt,inner sep=0]{\popdisplay\fontsize{12}{12}\selectfont\color{popcream}\thepopsec};%
  \hspace{8pt}{\popdisplay\fontsize{15}{17}\selectfont\color{popink}\MakeUppercase{#1}}\par
  \vspace{4pt}\noindent{\color{poporange}\rule{36pt}{3.6pt}}
  \end{minipage}\par\nopagebreak\vspace{8pt}
}
\newcounter{popsub}
\newcommand{\courssub}[1]{%
  \stepcounter{popsub}%
  \par\vspace{9pt}\noindent{\popxbold\fontsize{11.5}{13}\selectfont\color{popink}{\color{poporange}\thepopsec.\thepopsub}\hspace{6pt}#1}\par\nopagebreak\vspace{4pt}
}

% ============================================================
% Boîtes pédagogiques — même recette : fond blanc, bordure épaisse colorée,
% ombre dure encre, badge plein. Titre optionnel [#1].
% ============================================================
\tcbset{popbase/.style={breakable,enhanced,colback=white,boxrule=2.3pt,arc=9pt,
  shadow={3pt}{-3pt}{0pt}{popink},left=14pt,right=14pt,top=11pt,bottom=11pt,before skip=11pt,after skip=15pt,
  fontupper=\popsemi\fontsize{10.6}{14.5}\selectfont\color{popink},
  fontlower=\popsemi\fontsize{10.6}{14.5}\selectfont\color{popink}}}
% Coche dessinée (le glyphe ✓ n'existe pas dans les polices du document)
\newcommand{\popcheck}[1]{\tikz[baseline=-0.5ex]\draw[#1,line width=1.6pt,line cap=round,line join=round](-0.13,0.0)--(-0.04,-0.09)--(0.13,0.1);}
\newcommand{\popboxhead}[3]{\poppill{#1}{#2}{white}\ifx\relax#3\relax\else\hspace{6pt}{\popxbold\fontsize{10.6}{12}\selectfont\color{popink}#3}\fi\par\vspace{7pt}}

\newtcolorbox{aretenir}[1][]{popbase,colframe=popgreen,before upper={\popboxhead{À retenir}{popgreen}{#1}}}
\newtcolorbox{exemplebox}[1][]{popbase,colframe=poporange,before upper={\popboxhead{Exemple}{poporange}{#1}}}
\newtcolorbox{definition}[1][]{popbase,colframe=popblue,before upper={\popboxhead{Définition}{popblue}{#1}}}
\newtcolorbox{propriete}[1][]{popbase,colframe=poppurple,before upper={\popboxhead{Propriété}{poppurple}{#1}}}
\newtcolorbox{methode}[1][]{popbase,colframe=popgold,before upper={\popboxhead{Méthode}{popgold}{#1}}}
\newtcolorbox{attention}[1][]{popbase,colframe=poppink,before upper={\popboxhead{Attention}{poppink}{#1}}}
\newtcolorbox{experience}[1][]{popbase,colframe=popblue,colback=popblueL,before upper={\popboxhead{Expérience}{popblue}{#1}}}
% « Le savais-tu ? » : carte violette pleine (contexte, culture, Cameroun)
\newtcolorbox{savaistu}[1][]{popbase,colback=poppurple,colframe=popink,
  fontupper=\popsemi\fontsize{10.4}{14.2}\selectfont\color{white},
  before upper={\poppill{Le savais-tu\,?}{white}{poppurple}\ifx\relax#1\relax\else\hspace{6pt}{\popxbold\color{white}#1}\fi\par\vspace{7pt}}}
% Exercice résolu : énoncé en haut, \tcblower puis la solution sur fond crème
\newcounter{popexo}\newcounter{popexr}
\newtcolorbox{exoresolu}[1][]{popbase,colframe=poporange,colbacklower=popcream,
  segmentation style={popink,line width=1.2pt,dashed},
  before upper={\stepcounter{popexr}\popboxhead{Exercice résolu \thepopexr}{poporange}{#1}},
  before lower={\poppill{Solution}{popink}{popcream}\par\vspace{6pt}}}
% Exercice (non résolu en place) — la correction est dans la partie Corrigés
\newtcolorbox{exercice}[1][]{popbase,colframe=popink,
  before upper={\stepcounter{popexo}\popboxhead{Exercice \thepopexo}{popink}{#1}}}
% Corrigé
\newtcolorbox{corrige}[1][]{popbase,colframe=popgreen,colback=popgreenL,
  before upper={\popboxhead{Corrigé}{popgreen}{#1}}}
% Objectifs / prérequis / bilan
\newtcolorbox{objectifs}{popbase,colframe=popink,colback=poporangeL,
  before upper={\popboxhead{Objectifs de la leçon}{popink}{}}}
\newtcolorbox{prerequis}{popbase,colframe=popink,colback=popblueL,
  before upper={\popboxhead{Prérequis}{popblue}{}}}
\newtcolorbox{bilan}{popbase,colframe=popink,colback=white,boxrule=2.8pt,
  before upper={\popboxhead{Fiche bilan}{poporange}{L'essentiel à connaître pour l'examen}}}
% Figure encadrée avec légende
\newtcolorbox{popfigure}[1][]{enhanced,breakable=false,colback=white,colframe=popline,boxrule=1.4pt,arc=8pt,
  left=10pt,right=10pt,top=10pt,bottom=8pt,before skip=10pt,after skip=12pt,halign=center,
  lower separated=false,halign lower=center,
  fontlower=\popsemi\fontsize{9}{11.5}\selectfont\color{popmuted},
  #1}
\newcommand{\legende}[1]{\par\vspace{6pt}{\popsemi\fontsize{9}{11.5}\selectfont\color{popmuted}#1}}

% ============================================================
% Signature OnBuch+
% ============================================================
\newcommand{\poplogoinline}[1]{{\poplogo\fontsize{#1}{#1}\selectfont\color{popink}OnBuch\color{poporange}+}}
\newcommand{\signatureonbuch}{%
  \begin{tcolorbox}[enhanced,colback=popink,colframe=popink,boxrule=2.2pt,arc=10pt,
      drop shadow=poporange,left=14pt,right=14pt,top=10pt,bottom=10pt,before skip=0pt,after skip=0pt]
    \tikz[baseline=-0.7ex]\node[circle,fill=popgreen,draw=popcream,line width=1.3pt,minimum size=22pt,inner sep=0]{\popcheck{white}};%
    \hspace{9pt}\begin{minipage}[c]{0.62\linewidth}
      {\popxbold\fontsize{12}{13}\selectfont\color{popcream}Signé\ \textbullet\ {\poplogo OnBuch\color{poporange}+}}\par\vspace{2pt}
      {\raggedright\popsemi\fontsize{8}{10.5}\selectfont\color{popline}\MakeUppercase{Équipe pédagogique OnBuch+}\\\MakeUppercase{Conforme au programme officiel MINESEC}\par}
    \end{minipage}\hfill
    {\poplogo\fontsize{22}{22}\selectfont\color{popcream}OnBuch\color{poporange}+}
  \end{tcolorbox}%
}

% ============================================================
% Page de garde
% #1 titre · #2 matière · #3 niveau · #4 module · #5 n° leçon · #6 durée ·
% #7 description / objectifs (texte ou itemize)
% ============================================================
\newcommand{\pagegarde}[7]{%
  \thispagestyle{empty}
  \newgeometry{a4paper,margin=1.7cm,top=1.6cm,bottom=1.4cm}%
  \begin{tikzpicture}[remember picture,overlay]
    \fill[poporange!18] ([xshift=-3.2cm,yshift=-3.6cm]current page.north east) circle (4.6cm);
    \fill[poppurple!14] ([xshift=2.4cm,yshift=7.6cm]current page.south west) circle (3.8cm);
  \end{tikzpicture}%
  {\poplogo\fontsize{30}{30}\selectfont\color{popink}OnBuch\color{poporange}+}\hfill
  \raisebox{0.6ex}{\poppill{Cours signé \textbullet\ Premium}{popink}{popcream}}\par
  \vspace{1pt}\popsquiggle{poppurple}\par
  \vspace{8pt}
  \begin{tcolorbox}[enhanced,colback=poporange,colframe=popink,boxrule=2.8pt,arc=13pt,
      drop shadow=popink,left=18pt,right=18pt,top=12pt,bottom=13pt,after skip=10pt]
    \poppill{#2 \textbullet\ #3}{popink}{popcream}\hspace{5pt}\poppill{#4}{white}{popink}\par\vspace{8pt}
    {\popdisplay\fontsize{14}{14}\selectfont\color{popink}LEÇON #5}\par\vspace{4pt}
    {\raggedright\hyphenpenalty=10000\exhyphenpenalty=10000\popdisplay\fontsize{25}{27}\selectfont\color{popcream}\MakeUppercase{#1}\par}
    \vspace{8pt}
    {\popxbold\fontsize{9.5}{9.5}\selectfont\color{popink}\MakeUppercase{Durée conseillée : #6}}
  \end{tcolorbox}
  \begin{tcolorbox}[enhanced,colback=white,colframe=popink,boxrule=2.2pt,arc=10pt,
      drop shadow=popink,left=16pt,right=16pt,top=10pt,bottom=10pt,after skip=8pt]
    \poppill{Dans cette leçon}{poppurple}{white}\par\vspace{5pt}
    \popsemi\fontsize{9.3}{12.6}\selectfont\color{popink}#7
  \end{tcolorbox}
  \noindent\begin{minipage}[t]{0.487\linewidth}
    \begin{tcolorbox}[enhanced,colback=popgreen,colframe=popink,boxrule=2.2pt,arc=10pt,
        drop shadow=popink,left=11pt,right=11pt,top=10pt,bottom=10pt,height=3.1cm,valign=center]
      \tcbox[colback=white,colframe=popink,boxrule=1.3pt,arc=5pt,boxsep=0pt,nobeforeafter,
        left=3pt,right=3pt,top=3pt,bottom=3pt,valign=center]{\qrcode[height=1.9cm]{https://play.google.com/store/apps/details?id=com.luvvix.onbuch&hl=fr}}%
      \hspace{8pt}\begin{minipage}[c]{0.5\linewidth}
        {\popdisplay\fontsize{11}{12.5}\selectfont\color{white}\MakeUppercase{Télécharge OnBuch}}\par\vspace{4pt}
        {\raggedright\popsemi\fontsize{8.4}{11}\selectfont\color{white}Gratuit \textbullet\ Google Play\\Tuteur IA, annales, concours, résultats\par}
      \end{minipage}
    \end{tcolorbox}
  \end{minipage}\hfill
  \begin{minipage}[t]{0.487\linewidth}
    \begin{tcolorbox}[enhanced,colback=poppurple,colframe=popink,boxrule=2.2pt,arc=10pt,
        drop shadow=popink,left=11pt,right=11pt,top=10pt,bottom=10pt,height=3.1cm,valign=center]
      \tikz[baseline=-0.7ex]\node[circle,fill=white,draw=popink,line width=1.3pt,minimum size=1.6cm,inner sep=0]{\popdisplay\fontsize{24}{24}\selectfont\color{poppurple}+};%
      \hspace{8pt}\begin{minipage}[c]{0.6\linewidth}
        {\popdisplay\fontsize{11}{12.5}\selectfont\color{white}\MakeUppercase{Passe à OnBuch+}}\par\vspace{4pt}
        {\raggedright\popsemi\fontsize{8.4}{11}\selectfont\color{white}Tous les cours signés, Tuteur IA illimité, fiches PDF — onglet OnBuch+ de l'app\par}
      \end{minipage}
    \end{tcolorbox}
  \end{minipage}
  \vspace{8pt}
  \popcontenu
  \vfill
  \signatureonbuch
  \restoregeometry
  \newpage
}

% Rangée « Ce cours contient » (page de garde)
\newcommand{\poptile}[3]{% #1 couleur #2 gros texte #3 légende
  \begin{tcolorbox}[enhanced,colback=#1,colframe=popink,boxrule=2pt,arc=9pt,shadow={2.5pt}{-2.5pt}{0pt}{popink},
      left=6pt,right=6pt,top=9pt,bottom=9pt,halign=center,valign=center,height=1.75cm,width=\linewidth,nobeforeafter]
    {\popdisplay\fontsize{12.5}{13}\selectfont\color{white}\MakeUppercase{#2}}\par\vspace{3pt}
    {\centering\popsemi\fontsize{7.8}{9.5}\selectfont\color{white}#3\par}
  \end{tcolorbox}}
\newcommand{\popcontenu}{%
  \par\noindent{\popxboldls\fontsize{7.5}{7.5}\selectfont\color{popmuted}CE COURS SIGNÉ CONTIENT}\par\vspace{7pt}
  \noindent\begin{minipage}[t]{0.235\linewidth}\poptile{popblue}{Cours}{complet, illustré, pas à pas}\end{minipage}\hfill
  \begin{minipage}[t]{0.235\linewidth}\poptile{poporange}{Méthodes}{et exercices résolus}\end{minipage}\hfill
  \begin{minipage}[t]{0.235\linewidth}\poptile{poppink}{Exercices}{type examen + corrigés}\end{minipage}\hfill
  \begin{minipage}[t]{0.235\linewidth}\poptile{popgreen}{Fiche bilan}{l'essentiel à retenir}\end{minipage}\par}

% Sommaire
\newtcolorbox{sommaire}{enhanced,colback=white,colframe=popink,boxrule=2.2pt,arc=10pt,
  drop shadow=popink,left=18pt,right=18pt,top=13pt,bottom=13pt,before skip=6pt,after skip=20pt,
  before upper={\poppill{Sommaire}{poppurple}{white}\par\vspace{9pt}\popsemi\fontsize{10.6}{14.5}\selectfont\color{popink}}}

% Signature de fin de cours — poussée tout en bas de la dernière page
\newcommand{\findecours}{%
  \par\vfill\nopagebreak
  \begin{center}\popsquiggle{poporange}\end{center}\vspace{4pt}
  \signatureonbuch
}
% Glyphes absents de Fira Math : redéfinis à partir de glyphes disponibles
\AtBeginDocument{%
  \def\setminus{\mathbin{\backslash}}\let\smallsetminus\setminus
  \def\vdots{\mathord{\vbox{\baselineskip3.2pt\lineskiplimit0pt\kern4pt\hbox{.}\hbox{.}\hbox{.}}}}%
  \def\ddots{\mathinner{\mkern1mu\raise6.5pt\hbox{.}\mkern2mu\raise3.6pt\hbox{.}\mkern2mu\raise0.7pt\hbox{.}\mkern1mu}}%
  \def\triangle{\mathord{\scalebox{0.9}{$\symup{\Delta}$}}}%
  \def\nabla{\mathord{\raisebox{\depth}{\scalebox{1}[-1]{$\symup{\Delta}$}}}}%
  \def\checkmark{\popcheck{popdark}}%
  \def\star{\mathbin{\ast}}\def\bigstar{\mathord{\ast}}%
  \def\diamond{\mathbin{\scalebox{0.6}{$\Diamond$}}}%
  \def\bigcup{\mathop{\mathchoice{\raisebox{-0.3em}{\scalebox{1.7}{$\cup$}}}{\scalebox{1.25}{$\cup$}}{\cup}{\cup}}\displaylimits}%
  \def\bigcap{\mathop{\mathchoice{\raisebox{-0.3em}{\scalebox{1.7}{$\cap$}}}{\scalebox{1.25}{$\cap$}}{\cap}{\cap}}\displaylimits}%
  \def\lesseqqgtr{\mathrel{\lessgtr}}\def\bigoplus{\mathop{\oplus}\displaylimits}%
}
\providecommand{\ssi}{si et seulement si} % abréviation fréquente

\begin{document}

\pagegarde{Piles électrochimiques et corrosion}{Chimie}{Tle TI}{Module 3 · Oxydoréduction}{10}{6 h}{Cette leçon exploite la classification électrochimique pour prévoir les réactions d'oxydoréduction spontanées (règle du gamma), construire et analyser la pile Daniell (polarité, f.é.m, équations aux électrodes), puis expliquer le mécanisme de la corrosion du fer par micropiles et les méthodes de protection. Elle constitue l'aboutissement pratique du module d'oxydoréduction, très fréquemment évalué au Baccalauréat TI.
\vspace{4pt}
\begin{multicols}{2}\small
\begin{itemize}
\item À la fin de cette leçon, je sais exploiter le tableau de classification électrochimique des couples oxydant/réducteur
\item prévoir une réaction d'oxydoréduction naturelle entre deux couples redox en appliquant la règle du gamma
\item décrire la constitution d'une demi-pile Mⁿ⁺/M et de l'électrode standard à hydrogène (ESH) avec leurs notations conventionnelles
\item schématiser une pile Daniell, déterminer ses pôles et représenter la pile par sa notation conventionnelle
\item écrire les équations aux électrodes et l'équation-bilan de fonctionnement d'une pile
\item calculer la f.é.m d'une pile à partir des potentiels standard (complément Bac)
\end{itemize}\end{multicols}}

\begin{sommaire}
\begin{multicols}{2}
\begin{enumerate}
\item Classification électrochimique et règle du gamma
\item Demi-piles et électrode de référence
\item La pile Daniell : constitution et fonctionnement
\item Généralisation et prévision des piles
\item Corrosion des métaux : manifestations et mécanisme
\item Protection contre la corrosion
\item Travaux pratiques
\item Méthodes et exercices résolus
\item Exercices
\item Corrigés des exercices
\item Fiche bilan
\end{enumerate}
\end{multicols}
\end{sommaire}

\begin{objectifs}
\begin{itemize}
\item À la fin de cette leçon, je sais exploiter le tableau de classification électrochimique des couples oxydant/réducteur
\item prévoir une réaction d'oxydoréduction naturelle entre deux couples redox en appliquant la règle du gamma
\item décrire la constitution d'une demi-pile Mⁿ⁺/M et de l'électrode standard à hydrogène (ESH) avec leurs notations conventionnelles
\item schématiser une pile Daniell, déterminer ses pôles et représenter la pile par sa notation conventionnelle
\item écrire les équations aux électrodes et l'équation-bilan de fonctionnement d'une pile
\item calculer la f.é.m d'une pile à partir des potentiels standard (complément Bac)
\item réaliser expérimentalement une pile de type Daniell et mesurer sa f.é.m au voltmètre
\item expliquer la corrosion du fer par la formation de micropiles et citer ses facteurs favorisants
\item décrire et justifier les méthodes de protection contre la corrosion (revêtement, galvanisation, protection cathodique)
\end{itemize}
\end{objectifs}

\begin{prerequis}
\begin{itemize}
\item Couple oxydant/réducteur, oxydation, réduction, demi-équations électroniques (leçons 8 et 9)
\item Écriture et équilibrage d'une équation-bilan d'oxydoréduction
\item Nombre d'oxydation et variation du n.o. lors d'une réaction redox
\item Notions d'électricité : courant, tension, circuit, utilisation du voltmètre et de l'ampèremètre (Physique, Première)
\item Tests d'identification des ions métalliques en solution (Cu²⁺, Zn²⁺, Fe²⁺)
\end{itemize}
\end{prerequis}

\newpage

\coursec{Classification électrochimique et règle du gamma}

\textbf{Situation d'accroche.} À Douala, sur les chantiers du littoral, les barres de fer laissées à l'air libre se couvrent de rouille en quelques mois sous l'effet de l'humidité et de l'air salin. Pourtant, une simple couche de zinc déposée sur ces barres les protège pendant des années. Pourquoi le zinc protège-t-il le fer ? Par ailleurs, la batterie qui alimente ton téléphone et la pile d'une torche fonctionnent grâce aux mêmes principes que la pile inventée par John Daniell en 1836. Comment deux métaux plongés dans des solutions peuvent-ils produire du courant électrique ? Et surtout, comment \emph{prévoir}, sans réaliser la moindre expérience, quelle réaction va se produire spontanément entre deux métaux et leurs ions ? Cette leçon répond à ces questions grâce à deux outils puissants : la \cle{classification électrochimique} et la \cle{règle du gamma}.

\courssub{Rappels : couples oxydant/réducteur}

Avant de classer les couples, rappelons le vocabulaire fondamental de l'oxydoréduction.

\begin{definition}[Oxydant, réducteur, couple redox]
Un \cle{oxydant} est une espèce chimique (atome, ion ou molécule) capable de \textbf{capter} un ou plusieurs électrons. Un \cle{réducteur} est une espèce capable de \textbf{céder} un ou plusieurs électrons. Un oxydant \ce{Ox} et un réducteur \ce{Red} qui se correspondent par transfert d'électrons forment un \cle{couple oxydant/réducteur}, noté \ce{Ox}/\ce{Red}, reliés par la demi-équation électronique :
\[ \ce{Ox} + n\,e^- \ce{<=>} \ce{Red} \]
\end{definition}

\textbf{Exemples.} Le couple \ce{Cu^2+}/\ce{Cu} : $\ce{Cu^2+} + 2\,e^- \ce{<=>} \ce{Cu}$. Le couple \ce{Zn^2+}/\ce{Zn} : $\ce{Zn^2+} + 2\,e^- \ce{<=>} \ce{Zn}$. Le couple \ce{Ag+}/\ce{Ag} : $\ce{Ag+} + e^- \ce{<=>} \ce{Ag}$.

\begin{attention}[Sens de lecture d'une demi-équation]
La demi-équation $\ce{Ox} + n\,e^- \ce{<=>} \ce{Red}$ se lit dans \textbf{deux sens} : de gauche à droite, c'est une \textbf{réduction} (l'oxydant gagne des électrons) ; de droite à gauche, c'est une \textbf{oxydation} (le réducteur perd des électrons). Dans une réaction réelle, un seul sens est emprunté par chaque couple : un couple \emph{réduit} pendant que l'autre \emph{oxyde}. Il n'y a jamais d'électrons dans une équation-bilan !
\end{attention}

\courssub{Expériences qualitatives : qui réagit avec qui ?}

Pour comparer la « force » des oxydants et des réducteurs, réalisons trois expériences simples, réalisables dans n'importe quel laboratoire de lycée.

\begin{experience}[Métaux et ions métalliques : trois tests décisifs]
\textbf{Expérience 1 : lame de cuivre dans une solution de nitrate d'argent \ce{Ag+}.} On plonge une lame de cuivre rougeâtre dans une solution incolore contenant des ions \ce{Ag+}. \emph{Observation :} la lame se recouvre d'un dépôt gris brillant d'argent métallique (aspect « arbre de Diane ») et la solution bleuit progressivement : des ions \ce{Cu^2+} apparaissent. Il s'est produit :
\[ \ce{Cu(s) + 2 Ag+(aq) -> Cu^2+(aq) + 2 Ag(s)} \]
\textbf{Expérience 2 : lame de zinc dans une solution de sulfate de cuivre \ce{Cu^2+}.} On plonge une lame de zinc grise dans une solution bleue de sulfate de cuivre. \emph{Observation :} la lame se recouvre d'un dépôt rouge de cuivre et la solution se décolore lentement :
\[ \ce{Zn(s) + Cu^2+(aq) -> Zn^2+(aq) + Cu(s)} \]
\textbf{Expérience 3 : lame de cuivre dans une solution de sulfate de zinc \ce{Zn^2+}.} \emph{Observation :} rien ne se passe, même après plusieurs jours. \textbf{Aucune réaction.}
\end{experience}

\begin{popfigure}
\begin{center}
\begin{tikzpicture}[scale=0.95]
% Bécher
\draw[thick, popdark] (-2.2,0) -- (-2.2,3.4) (2.2,0) -- (2.2,3.4) (-2.2,0) -- (2.2,0);
% Solution bleue
\fill[popblue!35] (-2.15,0.05) rectangle (2.15,2.6);
\draw[popblue, thick] (-2.15,2.6) -- (2.15,2.6);
\node[popblue!70!black] at (1.15,1.1) {\small solution de \ce{CuSO4}};
\node[popblue!70!black] at (1.15,0.7) {\small (ions \ce{Cu^2+} : bleue)};
% Lame de zinc
\fill[popmuted!60] (-0.9,-0.1) rectangle (-0.2,4.0);
\draw[popdark] (-0.9,-0.1) rectangle (-0.2,4.0);
\node[rotate=90, popdark] at (-0.55,3.4) {\small lame de zinc \ce{Zn}};
% Dépôt de cuivre
\fill[poporange!80!black, decorate, decoration={random steps, segment length=2pt, amplitude=1pt}] (-0.95,0.0) rectangle (-0.15,2.2);
\node[poporange!60!black, align=center] at (-3.9,1.6) {\small dépôt rouge\\ \small de cuivre \ce{Cu}};
\draw[->, poporange!60!black, thick] (-2.9,1.5) -- (-1.05,1.3);
% Ions
\node[popblue!80!black] at (1.6,2.1) {\scriptsize \ce{Cu^2+}};
\node[popblue!80!black] at (0.6,1.9) {\scriptsize \ce{Cu^2+}};
\node[popgreen!60!black] at (0.5,0.5) {\scriptsize \ce{Zn^2+}};
\node[popgreen!60!black] at (1.7,0.35) {\scriptsize \ce{Zn^2+}};
\end{tikzpicture}
\end{center}
\legende{Lame de zinc plongée dans une solution bleue de sulfate de cuivre : le zinc s'oxyde en ions \ce{Zn^2+} (incolores) tandis que les ions \ce{Cu^2+} se réduisent en cuivre métallique rouge qui se dépose sur la lame.}
\end{popfigure}

\textbf{Interprétation.} Les expériences 1 et 2 montrent que les transferts d'électrons se font toujours dans le \emph{même sens} : le zinc donne ses électrons aux ions \ce{Cu^2+}, le cuivre donne ses électrons aux ions \ce{Ag+}. On dit que le zinc est un réducteur \emph{plus fort} que le cuivre, lui-même plus fort que l'argent. Inversement, \ce{Ag+} est un oxydant plus fort que \ce{Cu^2+}, lui-même plus fort que \ce{Zn^2+}. L'expérience 3 confirme que la réaction inverse ($\ce{Cu} + \ce{Zn^2+}$) est impossible spontanément.

\begin{definition}[Réaction d'oxydoréduction naturelle]
Une \cle{réaction d'oxydoréduction naturelle} (ou spontanée) est la réaction qui se produit d'elle-même, sans apport d'énergie extérieur, entre l'\textbf{oxydant d'un couple} et le \textbf{réducteur d'un autre couple}. Elle a toujours lieu entre l'oxydant le \emph{plus fort} et le réducteur le \emph{plus fort} mis en présence.
\end{definition}

\courssub{Pouvoir oxydant et pouvoir réducteur : construction de la classification}

\textbf{Intuition.} Imagine un concours de « tir à la corde » sur les électrons. Certains ions (comme \ce{Ag+}) attirent les électrons avec avidité : ce sont des oxydants \emph{forts}. Certains métaux (comme le sodium ou le zinc) lâchent leurs électrons facilement : ce sont des réducteurs \emph{forts}. En multipliant les expériences du type précédent entre tous les couples métalliques, on peut ranger les couples sur une échelle unique.

\begin{definition}[Classification électrochimique]
La \cle{classification électrochimique} est un tableau vertical dans lequel les couples \ce{Ox}/\ce{Red} sont rangés de telle sorte que :
\begin{itemize}
\item le \textbf{pouvoir oxydant} de l'oxydant \textbf{croît vers le haut} ;
\item le \textbf{pouvoir réducteur} du réducteur \textbf{croît vers le bas}.
\end{itemize}
Ainsi, l'oxydant le plus fort (\ce{Au^3+}, \ce{Ag+}\ldots) est en haut à gauche, et le réducteur le plus fort (\ce{K}, \ce{Ca}, \ce{Na}\ldots) est en bas à droite.
\end{definition}

\begin{popfigure}
\begin{center}
\begin{tikzpicture}[scale=1.0]
% Colonnes
\foreach \y/\ox/\red in {
  6.0/\ce{Au^3+}/\ce{Au},
  5.4/\ce{Ag+}/\ce{Ag},
  4.8/\ce{Cu^2+}/\ce{Cu},
  4.2/\ce{Pb^2+}/\ce{Pb},
  3.6/\ce{H+}/\ce{H2},
  3.0/\ce{Fe^2+}/\ce{Fe},
  2.4/\ce{Zn^2+}/\ce{Zn},
  1.8/\ce{Al^3+}/\ce{Al},
  1.2/\ce{Mg^2+}/\ce{Mg},
  0.6/\ce{Na+}/\ce{Na},
  0.0/\ce{Ca^2+}/\ce{Ca},
 -0.6/\ce{K+}/\ce{K}}{
  \node[anchor=east] at (-0.25,\y) {\ox};
  \node at (0.35,\y) {/};
  \node[anchor=west] at (0.6,\y) {\red};
  \draw[popline!60] (-1.6,\y-0.28) -- (2.0,\y-0.28);
}
% Cadres colonnes
\draw[thick, popdark] (-1.6,6.35) rectangle (2.0,-0.95);
\draw[thick, popdark] (0.2,6.35) -- (0.2,-0.95);
\node[popblue!70!black] at (-0.7,6.65) {\small \textbf{Oxydants}};
\node[poporange!80!black] at (1.3,6.65) {\small \textbf{Réducteurs}};
% Flèches de pouvoir
\draw[->, very thick, popblue] (-2.6,0.0) -- (-2.6,6.0) node[midway, left, rotate=90, align=center] {\small pouvoir oxydant croissant};
\draw[->, very thick, poporange] (3.0,6.0) -- (3.0,0.0) node[midway, right, rotate=90, align=center] {\small pouvoir réducteur croissant};
% Encadré H+/H2
\draw[popgreen!70!black, thick, rounded corners] (-1.55,3.32) rectangle (1.95,3.9);
\node[popgreen!50!black, anchor=west] at (3.6,3.6) {\small couple de référence};
\draw[popgreen!50!black] (3.55,3.6) -- (2.0,3.6);
\end{tikzpicture}
\end{center}
\legende{Classification électrochimique des couples usuels. Les oxydants sont de plus en plus forts vers le haut, les réducteurs de plus en plus forts vers le bas. Le couple \ce{H+}/\ce{H2} sert de repère central (électrode de référence, voir section 2).}
\end{popfigure}

\begin{aretenir}[Repères dans la classification]
\begin{itemize}
\item Les métaux alcalins (\ce{K}, \ce{Ca}, \ce{Na}) sont les réducteurs les plus \textbf{forts} : leurs couples sont tout en \textbf{bas}.
\item Les métaux nobles (\ce{Cu}, \ce{Ag}, \ce{Au}) sont des réducteurs \textbf{faibles} : leurs couples sont en \textbf{haut}. C'est pourquoi l'or et l'argent ne s'altèrent pas à l'air.
\item Le couple \ce{H+}/\ce{H2} occupe une position \textbf{intermédiaire} : il sépare les métaux « attaquables » par les acides de ceux qui ne le sont pas.
\end{itemize}
\end{aretenir}

\begin{savaistu}[Pourquoi conserve-t-on le sodium dans l'huile ?]
Le sodium est un réducteur \emph{très} fort, placé tout en bas de la classification. Il réduit violemment l'eau ($\ce{2 Na} + \ce{2 H2O} \rightarrow \ce{2 Na+} + \ce{2 OH-} + \ce{H2}$, réaction explosive qui enflamme le dihydrogène dégagé) et s'oxyde rapidement au contact du dioxygène de l'air. C'est pourquoi, au laboratoire, on le conserve \textbf{dans l'huile} (ou le pétrole), à l'abri de l'air et de l'humidité. À l'inverse, l'or, réducteur très faible, se trouve à l'état natif dans la nature : c'est un métal « noble ».
\end{savaistu}

\courssub{La règle du gamma : prévoir la réaction spontanée}

\textbf{Intuition.} Plaçons deux couples sur l'échelle verticale. La réaction naturelle met en jeu l'oxydant le plus fort (celui du couple le plus \emph{haut}) et le réducteur le plus fort (celui du couple le plus \emph{bas}). Si l'on trace le chemin qui va de cet oxydant vers ce réducteur en passant par les deux cases des couples, on dessine une lettre grecque : le \textbf{gamma} ($\gamma$).

\begin{propriete}[Règle du gamma]
Lorsqu'on met en présence les espèces de deux couples redox, la réaction \textbf{spontanée} possible a lieu entre l'\textbf{oxydant le plus fort} et le \textbf{réducteur le plus fort}. Sur la classification verticale, le chemin allant de l'oxydant du couple supérieur au réducteur du couple inférieur dessine un $\gamma$ :
\[ \ce{Ox_1} + \ce{Red_2} \ce{->} \ce{Red_1} + \ce{Ox_2} \]
où le couple 1 est placé \textbf{au-dessus} du couple 2 dans la classification. Si les espèces présentes ne permettent pas de tracer un gamma (oxydant faible face à réducteur faible), \textbf{aucune réaction} ne se produit.
\end{propriete}

\begin{popfigure}
\begin{center}
\begin{tikzpicture}[scale=1.05]
% Couple 1 (haut)
\draw[thick, popblue, rounded corners, fill=popblueL] (-0.2,2.6) rectangle (3.4,3.6);
\node at (1.6,3.1) {$\ce{Cu^2+} + 2\,e^- \ce{<=>} \ce{Cu}$};
\node[popblue!70!black, anchor=east] at (-0.4,3.1) {\small couple 1};
% Couple 2 (bas)
\draw[thick, poporange, rounded corners, fill=poporangeL] (-0.2,0.4) rectangle (3.4,1.4);
\node at (1.6,0.9) {$\ce{Zn^2+} + 2\,e^- \ce{<=>} \ce{Zn}$};
\node[poporange!80!black, anchor=east] at (-0.4,0.9) {\small couple 2};
% Gamma : de Ox1 (gauche haut) vers Red2 (droite bas)
\draw[->, very thick, poppink, rounded corners=6pt] (0.35,3.1) -- (-1.3,3.1) -- (-1.3,0.9) -- (2.6,0.9);
\node[poppink, rotate=90] at (-1.75,2.0) {\large $\gamma$};
% Étiquettes
\node[popblue!70!black, align=center] at (5.6,3.1) {\small oxydant le plus fort\\ \small \ce{Cu^2+} capte les électrons};
\draw[->, popblue!70!black] (4.35,3.1) -- (3.5,3.1);
\node[poporange!80!black, align=center] at (5.6,0.9) {\small réducteur le plus fort\\ \small \ce{Zn} cède les électrons};
\draw[->, poporange!80!black] (4.35,0.9) -- (3.5,0.9);
% Bilan
\node[align=center, popdark] at (1.6,-0.6) {\small \textbf{Réaction naturelle :} $\ce{Cu^2+} + \ce{Zn} \ce{->} \ce{Cu} + \ce{Zn^2+}$};
\end{tikzpicture}
\end{center}
\legende{La règle du gamma : le chemin $\gamma$ relie l'oxydant du couple supérieur (\ce{Cu^2+}) au réducteur du couple inférieur (\ce{Zn}). Ce sont ces deux espèces qui réagissent spontanément.}
\end{popfigure}

\begin{methode}[Appliquer la règle du gamma en 3 étapes]
\begin{enumerate}
\item \textbf{Repérer les couples} auxquels appartiennent les espèces présentes et les placer l'un au-dessus de l'autre selon la classification électrochimique.
\item \textbf{Tracer le gamma} : partir de l'oxydant du couple le plus haut, descendre et rejoindre le réducteur du couple le plus bas. Si les espèces présentes sont sur ce chemin, la réaction a lieu ; sinon, il n'y a \textbf{aucune réaction}.
\item \textbf{Écrire les demi-équations} dans le bon sens (réduction pour le couple du haut, oxydation pour le couple du bas), multiplier pour égaliser les électrons, puis \textbf{additionner} : les électrons disparaissent du bilan.
\end{enumerate}
\end{methode}

\begin{exemplebox}[Le zinc réduit les ions cuivre(II) ; l'inverse est impossible]
\textbf{Cas 1 : lame de zinc dans une solution de \ce{Cu^2+}.} Couples présents : \ce{Cu^2+}/\ce{Cu} (au-dessus) et \ce{Zn^2+}/\ce{Zn} (en dessous). Le gamma relie \ce{Cu^2+} à \ce{Zn} : la réaction a lieu.
\begin{align*}
\text{Réduction : } & \ce{Cu^2+} + 2\,e^- \ce{->} \ce{Cu}\\
\text{Oxydation : } & \ce{Zn} \ce{->} \ce{Zn^2+} + 2\,e^-\\[2pt]
\text{Bilan : } & \ce{Zn(s) + Cu^2+(aq) -> Zn^2+(aq) + Cu(s)}
\end{align*}
Réaction spontanée, totale, exothermique, observable par le dépôt rouge de cuivre et la décoloration de la solution.

\textbf{Cas 2 : lame de cuivre dans une solution de \ce{Zn^2+}.} Les espèces présentes sont \ce{Cu} (réducteur faible) et \ce{Zn^2+} (oxydant faible) : le gamma ne peut pas se tracer entre elles. \textbf{Aucune réaction} : la lame reste intacte.
\end{exemplebox}

\begin{attention}[Erreur fréquente : oxydant contre oxydant]
On ne peut \textbf{jamais} faire réagir l'oxydant d'un couple avec l'oxydant d'un autre couple (par exemple \ce{Cu^2+} avec \ce{Zn^2+}), ni un réducteur avec un réducteur (\ce{Cu} avec \ce{Zn}). Une oxydoréduction exige un \textbf{donneur} d'électrons (réducteur) \emph{et} un \textbf{accepteur} d'électrons (oxydant) : sans transfert d'électrons possible, aucune réaction. Vérifie toujours que ton bilan met en jeu un oxydant \emph{et} un réducteur, et que les électrons ont disparu de l'équation finale.
\end{attention}

\courssub{Application : action des acides sur les métaux}

Le couple \ce{H+}/\ce{H2} occupe une place charnière dans la classification. Les ions \ce{H+} apportés par un acide dilué (acide chlorhydrique, acide sulfurique) sont des oxydants : ils ne peuvent oxyder que les métaux réducteurs \emph{plus forts} que \ce{H2}, c'est-à-dire ceux placés \textbf{au-dessous} de l'hydrogène dans la classification.

\begin{propriete}[Métaux et acides]
Seuls les métaux situés \textbf{au-dessous} du couple \ce{H+}/\ce{H2} dans la classification (\ce{Fe}, \ce{Zn}, \ce{Al}, \ce{Mg}, \ce{Na}\ldots) sont attaqués par les ions \ce{H+} d'un acide dilué, avec \textbf{dégagement de dihydrogène} \ce{H2} :
\[ \ce{M(s)} + n\,\ce{H+(aq)} \ce{->} \ce{M^{n+}(aq)} + \frac{n}{2}\,\ce{H2(g)} \]
Les métaux situés au-dessus (\ce{Cu}, \ce{Ag}, \ce{Au}) ne sont \textbf{pas attaqués} : c'est pourquoi on utilise le cuivre pour les tuyauteries et les casseroles.
\end{propriete}

\begin{exoresolu}[Le fer est-il attaqué par l'acide chlorhydrique ? Et le cuivre ?]
On plonge séparément un clou en fer et une lame de cuivre dans deux béchers contenant une solution diluée d'acide chlorhydrique (\ce{H+} + \ce{Cl-}) de concentration \SI{0,50}{mol.L^{-1}}. Prévoir ce qui se passe dans chaque bécher et écrire les équations-bilans.
\tcblower
\textbf{Étape 1 — Couples présents.} Bécher 1 : \ce{H+}/\ce{H2} et \ce{Fe^2+}/\ce{Fe}. Bécher 2 : \ce{H+}/\ce{H2} et \ce{Cu^2+}/\ce{Cu}. Les ions \ce{Cl-} sont spectateurs.

\textbf{Étape 2 — Règle du gamma.} Dans la classification, \ce{H+}/\ce{H2} est \emph{au-dessus} de \ce{Fe^2+}/\ce{Fe} mais \emph{en dessous} de \ce{Cu^2+}/\ce{Cu}.
\begin{itemize}
\item Bécher 1 : le gamma relie \ce{H+} (oxydant le plus fort) à \ce{Fe} (réducteur le plus fort) : \textbf{réaction}.
\item Bécher 2 : les espèces présentes sont \ce{H+} (oxydant faible ici) et \ce{Cu} (réducteur faible) : le gamma ne se trace pas, \textbf{aucune réaction}.
\end{itemize}

\textbf{Étape 3 — Demi-équations et bilan (bécher 1).}
\begin{align*}
\text{Oxydation du fer : } & \ce{Fe} \ce{->} \ce{Fe^2+} + 2\,e^-\\
\text{Réduction de \ce{H+} : } & \ce{2 H+} + 2\,e^- \ce{->} \ce{H2}\\[2pt]
\text{Bilan : } & \ce{Fe(s) + 2 H+(aq) -> Fe^2+(aq) + H2(g)}
\end{align*}
\emph{Observations :} effervescence de dihydrogène sur le clou (détectable par une petite détonation à la flamme), disparition progressive du fer, solution verdâtre d'ions \ce{Fe^2+}. La lame de cuivre, elle, reste parfaitement intacte.
\end{exoresolu}

\begin{exercice}[À toi de jouer]
En utilisant la classification électrochimique, prévoir si une réaction a lieu dans chacun des cas suivants ; si oui, écrire l'équation-bilan :
\begin{enumerate}
\item une lame d'argent dans une solution de sulfate de cuivre \ce{Cu^2+} ;
\item une lame de cuivre dans une solution de nitrate d'argent \ce{Ag+} ;
\item de la grenaille de zinc dans une solution diluée d'acide sulfurique (\ce{H+}) ;
\item une lame de plomb dans une solution contenant des ions \ce{Zn^2+}.
\end{enumerate}
\end{exercice}

\begin{corrige}[Exercice 1]
\begin{enumerate}
\item \ce{Ag} est un réducteur plus faible que \ce{Cu} (couple \ce{Ag+}/\ce{Ag} au-dessus de \ce{Cu^2+}/\ce{Cu}) : \ce{Ag} et \ce{Cu^2+} sont tous deux « faibles » ici, le gamma ne se trace pas. \textbf{Aucune réaction.}
\item Le gamma relie \ce{Ag+} (oxydant fort, en haut) à \ce{Cu} (réducteur plus fort que \ce{Ag}) : réaction. Demi-équations : $\ce{Cu} \ce{->} \ce{Cu^2+} + 2\,e^-$ et $\ce{Ag+} + e^- \ce{->} \ce{Ag}$ ($\times 2$). Bilan : $\ce{Cu(s) + 2 Ag+(aq) -> Cu^2+(aq) + 2 Ag(s)}$ (dépôt gris d'argent, solution qui bleuit).
\item \ce{Zn} est en dessous de \ce{H+}/\ce{H2} : réaction. Bilan : $\ce{Zn(s) + 2 H+(aq) -> Zn^2+(aq) + H2(g)}$ (effervescence).
\item \ce{Pb} est un réducteur plus faible que \ce{Zn} : \ce{Pb} ne peut pas réduire \ce{Zn^2+}. \textbf{Aucune réaction.}
\end{enumerate}
\end{corrige}

\begin{aretenir}[L'essentiel de la section]
\begin{itemize}
\item Un couple redox \ce{Ox}/\ce{Red} obéit à $\ce{Ox} + n\,e^- \ce{<=>} \ce{Red}$.
\item La classification électrochimique range les couples : pouvoir oxydant croissant \textbf{vers le haut}, pouvoir réducteur croissant \textbf{vers le bas}.
\item \textbf{Règle du gamma} : la réaction spontanée a lieu entre l'oxydant du couple le plus haut et le réducteur du couple le plus bas ; sinon, aucune réaction.
\item Seuls les métaux placés \textbf{sous} l'hydrogène sont attaqués par les acides à \ce{H+}, avec dégagement de \ce{H2}.
\end{itemize}
Ces outils vont maintenant nous permettre de comprendre comment construire une pile : il « suffira » de séparer physiquement les deux couples pour forcer les électrons à circuler dans un circuit extérieur. C'est l'objet des sections suivantes.
\end{aretenir}

\coursec{Demi-piles et électrode de référence}

Dans la section précédente, nous avons appris à classer les couples oxydant/réducteur et à prévoir le sens naturel d'une réaction grâce à la règle du gamma. Mais une question demeure : d'où viennent ces classements ? Comment peut-on affirmer, chiffres à l'appui, que le cuivre est un réducteur plus faible que le zinc ? Pour répondre, il faut d'abord définir l'objet de base de toute pile : la \cle{demi-pile}, puis accepter une idée subtile — on ne peut pas mesurer le potentiel d'une demi-pile seule, il faut une \cle{référence}. C'est toute l'histoire de l'électrode standard à hydrogène.

\courssub{La demi-pile métallique : définition et notation}

\textbf{Une intuition pour commencer.}\par
Plonge un clou en fer dans un verre d'eau contenant du sulfate de fer(II) : il ne se passe rien de spectaculaire. Pourtant, à la surface du métal, un dialogue silencieux s'établit entre les atomes de fer du clou et les ions \ce{Fe^2+} de la solution. Quelques atomes de fer ont tendance à céder des électrons et à passer en solution ; quelques ions ont tendance à capter des électrons et à se déposer sur le clou. Cet équilibre dynamique entre le métal et ses ions, c'est exactement ce qu'on appelle une demi-pile.

\begin{definition}[Demi-pile]
Une \cle{demi-pile} est l'ensemble constitué d'un \textbf{conducteur} (métallique, comme une lame de zinc, ou inerte, comme le platine) plongé dans une solution électrolytique, le conducteur et la solution mettant en présence les \textbf{deux espèces d'un même couple oxydant/réducteur} \ce{Ox}/\ce{Red}.

Pour une demi-pile métallique associée au couple \ce{M^{n+}}/\ce{M}, elle est formée d'une \textbf{lame du métal M} plongeant dans une \textbf{solution contenant ses ions \ce{M^{n+}}}.
\end{definition}

Le nom « demi-pile » est parlant : il faudra \emph{deux} demi-piles, convenablement reliées, pour fabriquer une pile complète (section suivante). Une demi-pile seule ne peut faire circuler aucun courant durable, car les électrons n'ont nulle part où aller.

\textbf{Exemples de demi-piles métalliques.}\par
\begin{itemize}
\item Une lame de zinc dans une solution de sulfate de zinc (\ce{ZnSO4}) : couple \ce{Zn^2+}/\ce{Zn}.
\item Une lame de cuivre dans une solution de sulfate de cuivre (\ce{CuSO4}) : couple \ce{Cu^2+}/\ce{Cu}.
\item Une lame de fer dans une solution de sulfate de fer(II) (\ce{FeSO4}) : couple \ce{Fe^2+}/\ce{Fe}.
\end{itemize}

\begin{definition}[Notation conventionnelle d'une demi-pile]
Une demi-pile métallique se note :
\[ \ce{M} \; \big| \; \ce{M^{n+}} \]
La \textbf{barre verticale} symbolise l'\cle{interface} métal/solution, c'est-à-dire la surface de contact où s'échangent les électrons. On écrit le métal (réducteur) à gauche de la barre et l'ion (oxydant) à droite.
\end{definition}

Ainsi, les trois exemples précédents se notent respectivement :
\[ \ce{Zn | Zn^2+} \qquad \ce{Cu | Cu^2+} \qquad \ce{Fe | Fe^2+} \]

\begin{attention}[La barre verticale n'est pas une barre de fraction]
Le symbole $|$ dans $\ce{Zn | Zn^2+}$ ne signifie ni « divisé par » ni « ou ». Il représente la \textbf{frontière physique} entre deux phases : le solide métallique d'un côté, la solution ionique de l'autre. C'est précisément à cette interface que naît la différence de potentiel électrique.
\end{attention}

\begin{popfigure}
\begin{tikzpicture}[scale=1.0]
% Bécher
\draw[thick, popdark] (-1.6,0) -- (-1.6,3.2);
\draw[thick, popdark] (1.6,0) -- (1.6,3.2);
\draw[thick, popdark] (-1.6,0) -- (1.6,0);
% Solution
\fill[popblueL, opacity=0.7] (-1.55,0.05) rectangle (1.55,2.2);
\draw[popblue, thick] (-1.55,2.2) -- (1.55,2.2);
% Lame de zinc
\fill[popmuted!60] (-0.35,0.6) rectangle (0.35,3.9);
\draw[popdark, thick] (-0.35,0.6) rectangle (0.35,3.9);
\node[popdark, font=\small\bfseries] at (0,4.25) {Lame de zinc \ce{Zn}};
% Ions en solution
\node[popblue, font=\footnotesize] at (-1.0,1.6) {\ce{Zn^2+}};
\node[popblue, font=\footnotesize] at (1.0,1.2) {\ce{Zn^2+}};
\node[popblue, font=\footnotesize] at (-0.9,0.6) {\ce{SO4^2-}};
\node[popblue, font=\footnotesize] at (1.05,1.9) {\ce{SO4^2-}};
\node[popblue, font=\footnotesize] at (-1.1,1.1) {\ce{Zn^2+}};
% Annotations interface
\draw[-{Stealth}, poporange, thick] (2.6,1.5) -- (0.45,1.5);
\node[poporange, font=\footnotesize, align=left, anchor=west] at (2.7,1.5) {interface\\métal/solution};
% Étiquette solution
\node[popblue!70!popdark, font=\small, align=center] at (0,-0.5) {Solution de sulfate de zinc \ce{ZnSO4}};
% Notation
\node[popdark, font=\normalsize\bfseries] at (0,-1.1) {Demi-pile : \ce{Zn | Zn^2+}};
\end{tikzpicture}
\legende{Demi-pile zinc : une lame de zinc plonge dans une solution contenant ses ions \ce{Zn^2+}. L'équilibre \ce{Zn^2+ + 2e- <=> Zn} s'établit à la surface du métal.}
\end{popfigure}

\textbf{Que se passe-t-il microscopiquement ?}\par
À la surface de la lame, la demi-équation du couple s'établit dans les deux sens :
\[ \ce{Zn^2+ + 2e- <=> Zn} \]
Cet échange permanent de charges crée, entre le métal et la solution, une différence de potentiel appelée \cle{potentiel d'électrode} (ou potentiel redox) du couple, notée $E(\ce{Zn^2+}/\ce{Zn})$. Plus le métal a tendance à céder ses électrons (réducteur fort), plus sa surface se charge négativement par rapport à la solution, et plus ce potentiel est bas.

\courssub{Pourquoi une référence est-elle indispensable ?}

Voici maintenant le point délicat, et le plus beau de cette section. Tu voudrais mesurer le potentiel de la demi-pile $\ce{Zn | Zn^2+}$ avec un voltmètre. Un voltmètre possède \textbf{deux} bornes : il mesure toujours une \emph{différence} de potentiel entre deux points. Si tu relies une borne à la lame de zinc, où branches-tu l'autre ? Dans la solution ? Mais pour y plonger un fil, tu introduis... un second métal, donc une \emph{seconde} interface, donc une seconde demi-pile !

\begin{propriete}[Impossibilité de mesurer un potentiel absolu]
Le potentiel d'une demi-pile \textbf{isolée n'est pas mesurable directement}. Toute mesure de tension met en jeu \textbf{deux} demi-piles : on ne mesure jamais que des \textbf{différences} de potentiel d'électrode.
\end{propriete}

La situation est analogue à l'altitude en géographie : on ne peut pas définir la hauteur du Mont Cameroun « dans l'absolu » ; on la mesure \emph{par rapport au niveau de la mer}, choisi comme référence arbitraire ($\SI{0}{m}$). De même, les électrochimistes ont choisi une demi-pile de référence universelle, à laquelle on attribue par convention le potentiel zéro : l'\cle{électrode standard à hydrogène}.

\courssub{L'électrode standard à hydrogène (ESH)}

\textbf{Pourquoi l'hydrogène ?}\par
Le couple \ce{H+}/\ce{H2} est simple, universel, reproductible dans n'importe quel laboratoire, et l'hydrogène est l'élément le plus abondant de l'Univers. C'est un excellent « niveau de la mer » de l'électrochimie. Un détail pratique se pose cependant : \ce{H2} est un gaz et \ce{H+} est en solution — aucun des deux n'est un métal conducteur. Il faut donc un conducteur \emph{inerte} pour recueillir les électrons : le \textbf{platine}.

\begin{definition}[Électrode standard à hydrogène (ESH)]
L'\cle{électrode standard à hydrogène} est la demi-pile de référence constituée de :
\begin{itemize}
\item un \textbf{fil de platine platiné} (recouvert de noir de platine, qui augmente la surface active et catalyse la réaction) ;
\item un courant de \textbf{dihydrogène gazeux} à la pression $p(\ce{H2}) = \SI{1}{bar}$ qui barbote autour du platine ;
\item une solution acide de concentration en ions oxonium $[\ce{H+}] = \SI{1}{mol.L^{-1}}$ ;
\item le tout maintenu à la température de $\SI{25}{\celsius}$.
\end{itemize}
Ces conditions ($\SI{1}{bar}$, $\SI{1}{mol.L^{-1}}$, $\SI{25}{\celsius}$) sont les \textbf{conditions standard}.
\end{definition}

La demi-équation qui s'établit à la surface du platine est :
\[ \ce{2H+ + 2e- <=> H2} \]
Il s'agit d'un équilibre rapide et réversible grâce à l'action catalytique du platine platiné.

\begin{definition}[Notation conventionnelle de l'ESH]
L'ESH se note :
\[ \ce{Pt} \; \big| \; \ce{H2}\,(\SI{1}{bar}) \; \big| \; \ce{H+}\,(\SI{1}{mol.L^{-1}}) \]
Chaque barre verticale symbolise une interface entre deux phases : platine solide $|$ gaz $|$ solution.
\end{definition}

\begin{attention}[Le platine ne participe pas à la réaction]
Le platine est un conducteur \textbf{inerte} : il n'apparaît \textbf{jamais} dans l'équation chimique. Son double rôle est :
\begin{enumerate}
\item \textbf{conduire} les électrons vers ou depuis le circuit extérieur ;
\item \textbf{catalyser} la réaction \ce{2H+ + 2e- <=> H2} grâce à sa surface platinée, qui adsorbe les molécules de \ce{H2} et facilite leur dissociation.
\end{enumerate}
Écrire \ce{Pt} dans une équation-bilan est une erreur classique au Baccalauréat : le platine n'est ni oxydé ni réduit.
\end{attention}

\begin{popfigure}
\begin{tikzpicture}[scale=1.0]
% Bécher
\draw[thick, popdark] (-2.0,0) -- (-2.0,3.6);
\draw[thick, popdark] (2.0,0) -- (2.0,3.6);
\draw[thick, popdark] (-2.0,0) -- (2.0,0);
% Solution acide
\fill[poppurpleL, opacity=0.6] (-1.95,0.05) rectangle (1.95,2.6);
\draw[poppurple, thick] (-1.95,2.6) -- (1.95,2.6);
\node[poppurple!70!popdark, font=\footnotesize] at (-1.2,2.2) {\ce{H+}};
\node[poppurple!70!popdark, font=\footnotesize] at (1.3,1.4) {\ce{H+}};
\node[poppurple!70!popdark, font=\footnotesize] at (-1.3,0.8) {\ce{H+}};
% Cloche à hydrogène
\draw[thick, popdark] (-0.7,4.6) -- (-0.7,1.0);
\draw[thick, popdark] (0.7,4.6) -- (0.7,1.0);
\draw[thick, popdark] (-0.7,4.6) -- (0.7,4.6);
% Fil de platine
\draw[very thick, popgold!80!popdark] (0,4.6) -- (0,0.5);
\node[popgold!70!popdark, font=\footnotesize, anchor=west] at (0.15,4.9) {Fil de platine platiné \ce{Pt}};
% Bulles de H2
\foreach \x/\y/\r in {-0.45/1.4/0.10, -0.3/1.9/0.08, 0.4/1.6/0.11, 0.25/2.1/0.07, -0.15/2.35/0.09, 0.45/2.4/0.08, -0.5/2.5/0.07}{
\draw[popblue, fill=popblueL] (\x,\y) circle (\r);}
\node[popblue, font=\footnotesize, anchor=west] at (1.0,2.9) {Bulles de \ce{H2} à \SI{1}{bar}};
\draw[-{Stealth}, popblue] (1.0,2.8) -- (0.3,2.2);
% Arrivée H2
\draw[-{Stealth}, very thick, popgreen!70!popdark] (-2.6,4.2) -- (-0.75,4.2);
\node[popgreen!60!popdark, font=\footnotesize, anchor=east] at (-2.65,4.2) {Arrivée de \ce{H2}};
% Annotations conditions
\node[popdark, font=\footnotesize, align=left, anchor=west] at (2.3,0.9) {$[\ce{H+}] = \SI{1}{mol.L^{-1}}$\\$\theta = \SI{25}{\celsius}$};
\draw[-{Stealth}, popdark] (2.3,0.9) -- (1.6,0.7);
% Étiquette
\node[popdark, font=\normalsize\bfseries, align=center] at (0,-0.6) {ESH : \ce{Pt | H2 (1 bar) | H+ (1 mol.L^{-1})}};
\end{tikzpicture}
\legende{L'électrode standard à hydrogène : le dihydrogène barbote autour du platine platiné plongé dans une solution acide. L'équilibre \ce{2H+ + 2e- <=> H2} s'établit à la surface du platine.}
\end{popfigure}

\courssub{Le zéro de l'échelle des potentiels}

\begin{propriete}[Convention fondamentale]
Par convention internationale, le potentiel standard de l'ESH est pris égal à \textbf{zéro} à toute température :
\[ E^\circ(\ce{H+}/\ce{H2}) = \SI{0,00}{V} \]
Tous les autres potentiels d'électrode sont mesurés \textbf{par rapport à cette référence}.
\end{propriete}

Cette convention donne enfin un sens précis au tableau de classification électrochimique de la section précédente : chaque couple y est repéré par son \cle{potentiel standard} $E^\circ$, mesuré contre l'ESH.

\begin{definition}[Potentiel standard d'un couple (complément Bac)]
Le \cle{potentiel standard} $E^\circ(\ce{Ox}/\ce{Red})$ d'un couple est, par définition, la \textbf{f.é.m. de la pile} formée en associant :
\begin{itemize}
\item la demi-pile standard du couple étudié (concentrations à $\SI{1}{mol.L^{-1}}$, $\SI{25}{\celsius}$) ;
\item l'électrode standard à hydrogène.
\end{itemize}
Par convention de signe : $E^\circ$ est \textbf{positif} si la demi-pile étudiée constitue le pôle \textbf{positif} de la pile (le couple y subit une réduction), et \textbf{négatif} si elle constitue le pôle \textbf{négatif} (le couple y subit une oxydation).
\end{definition}

\textbf{Lecture du signe et de la valeur.}\par
\begin{itemize}
\item Plus $E^\circ$ est \textbf{grand (positif)}, plus l'\textbf{oxydant} du couple est \textbf{fort} (il capte facilement les électrons, même face à \ce{H+}).
\item Plus $E^\circ$ est \textbf{petit (négatif)}, plus le \textbf{réducteur} du couple est \textbf{fort} (il cède facilement ses électrons, même face à \ce{H2}).
\end{itemize}

\begin{exemplebox}[Deux mesures historiques]
\textbf{Cas du cuivre.} On réalise la pile $\ce{Cu | Cu^2+}$ (1 mol/L) associée à l'ESH. Le voltmètre indique $\SI{0,34}{V}$, le cuivre étant le pôle positif : les ions \ce{Cu^2+} s'y réduisent ($\ce{Cu^2+ + 2e- -> Cu}$) pendant que \ce{H2} s'oxyde à l'ESH. Donc :
\[ E^\circ(\ce{Cu^2+}/\ce{Cu}) = \SI{+0,34}{V} \]
\textbf{Cas du zinc.} Avec la demi-pile $\ce{Zn | Zn^2+}$ (1 mol/L), le voltmètre indique $\SI{0,76}{V}$, mais c'est le \textbf{zinc} qui est le pôle négatif : le zinc s'oxyde ($\ce{Zn -> Zn^2+ + 2e-}$) pendant que \ce{H+} se réduit à l'ESH. Donc :
\[ E^\circ(\ce{Zn^2+}/\ce{Zn}) = \SI{-0,76}{V} \]
\end{exemplebox}

\begin{exemplebox}[Interpréter une valeur de potentiel standard]
$E^\circ(\ce{Zn^2+}/\ce{Zn}) = \SI{-0,76}{V} < \SI{0,00}{V}$ signifie concrètement que le zinc métallique est un \textbf{réducteur plus fort que \ce{H2}} : il cède ses électrons plus facilement que le dihydrogène. Conséquence directe (règle du gamma) : le zinc réduit spontanément les ions \ce{H+} d'un acide :
\[ \ce{Zn + 2H+ -> Zn^2+ + H2 ^} \]
C'est pourquoi le zinc s'attaque dans le vinaigre ou l'acide chlorhydrique avec dégagement de dihydrogène, alors que le cuivre ($E^\circ = \SI{+0,34}{V}$, réducteur plus faible que \ce{H2}) reste inerte dans ces mêmes acides non oxydants.
\end{exemplebox}

\begin{exoresolu}[Déterminer un potentiel standard à partir d'une mesure de f.é.m.]
\textbf{Énoncé.} On associe une demi-pile $\ce{Fe | Fe^2+}$ ($[\ce{Fe^2+}] = \SI{1}{mol.L^{-1}}$, $\SI{25}{\celsius}$) à une ESH. Le voltmètre mesure une f.é.m. $e = \SI{0,44}{V}$, et l'on constate que la lame de fer constitue le pôle négatif de la pile.
\begin{enumerate}
\item Écrire les demi-équations aux deux électrodes et l'équation-bilan de fonctionnement.
\item En déduire la valeur de $E^\circ(\ce{Fe^2+}/\ce{Fe})$.
\item Le fer est-il un réducteur plus fort ou plus faible que le dihydrogène ? Que se passe-t-il quand on plonge un clou en fer dans une solution d'acide chlorhydrique dilué ?
\end{enumerate}
\tcblower
\textbf{Solution détaillée.}
\begin{enumerate}
\item Le fer est le pôle \textbf{négatif} : il fournit des électrons au circuit, donc il s'\textbf{oxyde} :
\[ \ce{Fe -> Fe^2+ + 2e-} \quad \text{(anode, oxydation)} \]
À l'ESH, pôle positif, les ions \ce{H+} captent les électrons et se \textbf{réduisent} :
\[ \ce{2H+ + 2e- -> H2} \quad \text{(cathode, réduction)} \]
Les deux demi-équations échangent le même nombre d'électrons (2) ; l'équation-bilan s'obtient par simple addition :
\[ \ce{Fe + 2H+ -> Fe^2+ + H2 ^} \]
C'est une réaction \textbf{spontanée}, sans catalyseur, à température ambiante, avec dégagement gazeux de dihydrogène.
\item Par définition, $e = E^\circ_{\text{pôle }+} - E^\circ_{\text{pôle }-} = E^\circ(\ce{H+}/\ce{H2}) - E^\circ(\ce{Fe^2+}/\ce{Fe})$. Comme $E^\circ(\ce{H+}/\ce{H2}) = \SI{0,00}{V}$ :
\[ E^\circ(\ce{Fe^2+}/\ce{Fe}) = \SI{0,00}{V} - \SI{0,44}{V} = \boxed{\SI{-0,44}{V}} \]
Le signe négatif est cohérent avec le fait que le fer est le pôle négatif (il s'oxyde face à \ce{H+}).
\item $E^\circ(\ce{Fe^2+}/\ce{Fe}) < \SI{0,00}{V}$, donc le fer est un réducteur \textbf{plus fort} que \ce{H2}. Par la règle du gamma, il réduit spontanément les ions \ce{H+} : plongé dans l'acide chlorhydrique dilué, le clou se dissout lentement en dégageant du dihydrogène selon l'équation-bilan écrite ci-dessus. C'est exactement le phénomène qui fragilise les pièces en fer en milieu acide et humide — première étape vers la corrosion, que nous étudierons en fin de leçon.
\end{enumerate}
\end{exoresolu}

\begin{attention}[Pièges fréquents au Baccalauréat]
\begin{itemize}
\item \textbf{Ne pas confondre} la f.é.m. mesurée (toujours \emph{positive}) avec le potentiel standard (qui peut être \emph{négatif}). C'est la polarité de la demi-pile étudiée qui fixe le signe de $E^\circ$.
\item L'ESH n'est « standard » que si les \textbf{trois} conditions sont réunies : $p(\ce{H2}) = \SI{1}{bar}$, $[\ce{H+}] = \SI{1}{mol.L^{-1}}$, $\theta = \SI{25}{\celsius}$. Si l'une manque, le potentiel n'est plus $\SI{0,00}{V}$.
\item Dans la notation $\ce{Pt | H2 | H+}$, ne pas oublier le platine : sans conducteur, il n'y a pas d'électrode !
\end{itemize}
\end{attention}

\begin{savaistu}[Pourquoi ton pH-mètre n'utilise pas d'ESH]
L'ESH est un chef-d'œuvre théorique mais un cauchemar pratique : il faut une bouteille de dihydrogène (gaz inflammable !), une pression rigoureusement contrôlée, du platine platiné coûteux et fragile, et une solution acide à concentration exacte. En pratique, les laboratoires — y compris ceux de la SONARA ou des lycées camerounais bien équipés — utilisent des électrodes de référence secondaires, plus sûres et stables : l'\textbf{électrode au calomel} ($E = \SI{+0,27}{V}$ par rapport à l'ESH) et l'\textbf{électrode argent/chlorure d'argent} \ce{Ag}/\ce{AgCl} ($E \approx \SI{+0,22}{V}$). Ces électrodes sont soigneusement \textbf{étalonnées par rapport à l'ESH}, qui demeure la référence absolue de toute l'échelle des potentiels — exactement comme les altitudes locales restent rattachées au niveau de la mer.
\end{savaistu}

\begin{aretenir}[L'essentiel de la section]
\begin{itemize}
\item Une \textbf{demi-pile} associe un conducteur et les deux espèces d'un couple redox ; une demi-pile métallique se note $\ce{M | M^{n+}}$, la barre symbolisant l'interface métal/solution.
\item Le potentiel d'une demi-pile isolée est \textbf{inaccessible} : toute mesure est une différence, d'où la nécessité d'une référence.
\item L'\textbf{ESH} ($\ce{Pt | H2 (1 bar) | H+ (1 mol.L^{-1})}$ à $\SI{25}{\celsius}$) est la référence universelle : $E^\circ(\ce{H+}/\ce{H2}) = \SI{0,00}{V}$ par convention. Le platine y est inerte : conducteur et catalyseur seulement.
\item Le \textbf{potentiel standard} $E^\circ$ d'un couple est la f.é.m. de la pile « demi-pile étudiée + ESH », affectée du signe $+$ si la demi-pile est le pôle positif. Plus $E^\circ$ est élevé, plus l'oxydant est fort ; plus il est bas, plus le réducteur est fort.
\end{itemize}
\end{aretenir}

Nous disposons désormais de toutes les briques : des demi-piles bien définies et une échelle de potentiels étalonnée. Il ne reste plus qu'à assembler \emph{deux} demi-piles pour construire une véritable pile génératrice de courant : ce sera la célèbre pile Daniell, héroïne de la section suivante.

\coursec{La pile Daniell : constitution et fonctionnement}

Dans la section précédente, nous avons appris à construire une \cle{demi-pile} $M^{n+}/M$ et à mesurer son potentiel grâce à l'électrode standard à hydrogène (ESH). Nous savons donc que le couple \ce{Cu^2+}/\ce{Cu} a un potentiel standard $E^\circ = \SI{+0,34}{V}$ et le couple \ce{Zn^2+}/\ce{Zn} un potentiel $E^\circ = \SI{-0,76}{V}$. Une question se pose alors naturellement : si l'on \emph{relie} deux demi-piles différentes, que se passe-t-il ? Réponse : les électrons circulent, un courant électrique naît, et la réaction d'oxydoréduction spontanée devient une \emph{source d'énergie électrique}. C'est exactement le principe de la pile.

\courssub{Historique et principe général}

\begin{savaistu}[John Daniell et la révolution du télégraphe]
En 1836, le chimiste britannique \textbf{John Frederic Daniell} (1790--1845) met au point une pile associant le zinc et le cuivre, bien plus stable que la pile Volta de 1800. Sa pile alimenta les premières lignes de \textbf{télégraphe} électrique au XIX\textsuperscript{e} siècle : c'est elle qui a permis les premières communications rapides à longue distance ! Aujourd'hui encore, le même principe (transformer une réaction redox spontanée en courant électrique) fait fonctionner les batteries des motos-taxis de Yaoundé et de Douala, les piles de nos lampes torches et les batteries de nos téléphones portables.
\end{savaistu}

\begin{definition}[Pile électrochimique]
Une \cle{pile électrochimique} est un \textbf{générateur} qui transforme l'\textbf{énergie chimique} libérée par une réaction d'oxydoréduction \textbf{spontanée} en \textbf{énergie électrique}. Elle est constituée de \textbf{deux demi-piles} reliées électriquement.
\end{definition}

L'idée fondamentale est la suivante : la réaction entre le zinc métallique et les ions cuivre(II) est spontanée (nous l'avons prévue avec la règle du gamma) :
\[ \ce{Zn(s) + Cu^2+(aq) -> Zn^2+(aq) + Cu(s)} \]
Si l'on plonge simplement une lame de zinc dans une solution de sulfate de cuivre, cette réaction a lieu \emph{directement} au contact des réactifs : l'énergie est perdue en chaleur, aucun courant n'est récupéré. Le génie de Daniell consiste à \textbf{séparer physiquement} les deux réactifs dans deux compartiments distincts : les électrons sont alors \emph{obligés} de passer par un circuit extérieur pour aller du zinc vers les ions cuivre. Ce flux d'électrons, c'est le courant électrique que l'on peut exploiter (lampe, moteur, télégraphe...).

\courssub{Constitution de la pile Daniell}

\begin{experience}[Réalisation de la pile Daniell]
\textbf{Matériel :} deux béchers, une lame de zinc, une lame de cuivre, une solution de sulfate de zinc \ce{ZnSO4} à \SI{1}{mol.L^{-1}}, une solution de sulfate de cuivre \ce{CuSO4} à \SI{1}{mol.L^{-1}}, un tube en U rempli d'une solution gélifiée de \ce{KCl} (ou \ce{KNO3}), des fils de connexion, un voltmètre.

\textbf{Protocole :}
\begin{enumerate}
\item Dans le premier bécher, plonger la lame de zinc dans la solution de \ce{Zn^2+} : c'est la demi-pile \ce{Zn^2+}/\ce{Zn}.
\item Dans le second bécher, plonger la lame de cuivre dans la solution de \ce{Cu^2+} : c'est la demi-pile \ce{Cu^2+}/\ce{Cu}.
\item Relier les deux lames par un fil conducteur intercalé avec un voltmètre.
\item Plonger les deux extrémités du tube en U (pont salin) dans chacune des solutions.
\end{enumerate}

\textbf{Observations :} le voltmètre indique une tension d'environ \SI{1,10}{V}, la borne du voltmètre reliée au cuivre étant le pôle $+$. Au bout d'un certain temps, la lame de zinc s'amincit, un dépôt rougeâtre de cuivre apparaît sur la lame de cuivre, et la couleur bleue de la solution de \ce{CuSO4} pâlit.

\textbf{Interprétation :} un courant circule ; le zinc est consommé (oxydation) et le cuivre se dépose (réduction).
\end{experience}

La pile Daniell comprend donc \textbf{quatre éléments indispensables} :
\begin{itemize}
\item une demi-pile \ce{Zn} $\mid$ \ce{Zn^2+} (\SI{1}{mol.L^{-1}}) ;
\item une demi-pile \ce{Cu} $\mid$ \ce{Cu^2+} (\SI{1}{mol.L^{-1}}) ;
\item un \cle{pont salin} : tube en U contenant une solution ionique gélifiée (\ce{KCl} ou \ce{KNO3}) ;
\item un circuit électrique extérieur (fils, voltmètre ou récepteur).
\end{itemize}

\begin{definition}[Rôle du pont salin]
Le \cle{pont salin} a une double fonction :
\begin{enumerate}
\item \textbf{fermer le circuit} électrique en permettant le passage du courant entre les deux compartiments (par migration d'ions), \emph{sans mélanger} les réactifs \ce{Zn} et \ce{Cu^2+} ;
\item \textbf{assurer la neutralité électrique} des deux solutions : il fournit des cations \ce{K+} au compartiment qui s'appauvrit en charges positives (compartiment du cuivre) et des anions \ce{Cl-} au compartiment qui s'enrichit en charges positives (compartiment du zinc).
\end{enumerate}
Sans pont salin, les charges s'accumuleraient dans chaque bécher et la pile s'arrêterait presque instantanément.
\end{definition}

\begin{popfigure}
\begin{center}
\begin{tikzpicture}[scale=0.95, >=Stealth]
% Bécher gauche (zinc)
\draw[thick, popdark] (-4.6,0) -- (-4.6,2.6) (-2.2,0) -- (-2.2,2.6) (-4.6,0) -- (-2.2,0);
\draw[fill=popblueL, draw=none] (-4.55,0.05) rectangle (-2.25,1.7);
\draw[popblue] (-4.55,1.7) -- (-2.25,1.7);
\node[popblue, font=\small] at (-3.4,0.5) {\ce{Zn^2+} + \ce{SO4^2-}};
\node[popblue, font=\small] at (-3.4,1.95) {\SI{1}{mol.L^{-1}}};
% Lame de zinc
\draw[fill=popmuted!60, draw=popdark, thick] (-3.75,0.3) rectangle (-3.35,3.4);
\node[font=\small\bfseries] at (-3.55,3.65) {Zn};
\node[font=\small, popdark] at (-3.55,-0.45) {$\ominus$ anode};
% Bécher droit (cuivre)
\draw[thick, popdark] (2.2,0) -- (2.2,2.6) (4.6,0) -- (4.6,2.6) (2.2,0) -- (4.6,0);
\draw[fill=popblue!35, draw=none] (2.25,0.05) rectangle (4.55,1.7);
\draw[popblue] (2.25,1.7) -- (4.55,1.7);
\node[popblue, font=\small] at (3.4,0.5) {\ce{Cu^2+} + \ce{SO4^2-}};
\node[popblue, font=\small] at (3.4,1.95) {\SI{1}{mol.L^{-1}}};
% Lame de cuivre
\draw[fill=poporange!70, draw=popdark, thick] (3.35,0.3) rectangle (3.75,3.4);
\node[font=\small\bfseries] at (3.55,3.65) {Cu};
\node[font=\small, popdark] at (3.55,-0.45) {$\oplus$ cathode};
% Fils + voltmètre
\draw[thick, popdark] (-3.55,3.4) -- (-3.55,4.4) -- (-0.7,4.4);
\draw[thick, popdark] (3.55,3.4) -- (3.55,4.4) -- (0.7,4.4);
\draw[thick, popdark, fill=white] (0,4.4) circle (0.55);
\node[font=\bfseries] at (0,4.4) {V};
\node[font=\small] at (0,5.15) {$E \approx \SI{1,10}{V}$};
% Flèche électrons
\draw[->, very thick, poppink] (-2.9,4.9) .. controls (-1.5,5.6) and (1.5,5.6) .. (2.9,4.9);
\node[poppink, font=\small\bfseries] at (0,5.85) {électrons $e^-$};
% Pont salin
\draw[thick, popgreen!70!black, fill=popgreenL] (-1.5,1.2) -- (-1.5,3.0) -- (-1.1,3.0) -- (-1.1,1.2) -- (1.1,1.2) -- (1.1,3.0) -- (1.5,3.0) -- (1.5,1.2) -- cycle;
\node[popgreen!60!black, font=\small\bfseries] at (0,2.2) {pont salin};
\node[popgreen!60!black, font=\small] at (0,1.75) {(\ce{K+} + \ce{Cl-})};
% Migrations ioniques
\draw[->, thick, poporange] (0.4,1.35) -- (1.9,1.35);
\node[poporange, font=\small] at (1.15,1.6) {\ce{K+}};
\draw[->, thick, poppurple] (-0.4,0.95) -- (-1.9,0.95);
\node[poppurple, font=\small] at (-1.15,0.7) {\ce{Cl-}};
% Demi-équations
\node[font=\small, align=center, popdark] at (-5.9,2.9) {\ce{Zn -> Zn^2+ + 2e-}};
\node[font=\small, align=center, popdark] at (5.9,2.9) {\ce{Cu^2+ + 2e- -> Cu}};
\end{tikzpicture}
\end{center}
\legende{Schéma fonctionnel de la pile Daniell : le zinc s'oxyde (anode, pôle $\ominus$), les ions cuivre se réduisent (cathode, pôle $\oplus$), les électrons circulent de Zn vers Cu dans le circuit extérieur.}
\end{popfigure}

\courssub{Polarité de la pile : prévision par la classification}

\begin{methode}[Déterminer les pôles d'une pile]
Pour prévoir la polarité d'une pile à partir de la \textbf{classification électrochimique} :
\begin{enumerate}
\item Identifier les deux couples redox mis en jeu : ici \ce{Zn^2+}/\ce{Zn} ($E^\circ = \SI{-0,76}{V}$) et \ce{Cu^2+}/\ce{Cu} ($E^\circ = \SI{+0,34}{V}$).
\item Comparer leurs positions : le couple dont le \textbf{réducteur est le plus fort} (le plus haut dans la classification, potentiel le plus négatif) fournit le métal qui \textbf{s'oxyde}.
\item Ce métal constitue le \textbf{pôle négatif} de la pile (l'anode) ; l'autre électrode est le \textbf{pôle positif} (la cathode).
\end{enumerate}
Ici, le zinc est un réducteur plus fort que le cuivre : c'est donc le zinc qui s'oxyde. La lame de zinc est le pôle $\ominus$, la lame de cuivre le pôle $\oplus$.
\end{methode}

\begin{definition}[Anode et cathode]
\begin{itemize}
\item L'\cle{anode} est l'électrode où a lieu l'\textbf{oxydation} ; dans une pile, c'est le \textbf{pôle négatif} ($\ominus$).
\item La \cle{cathode} est l'électrode où a lieu la \textbf{réduction} ; dans une pile, c'est le \textbf{pôle positif} ($\oplus$).
\end{itemize}
Moyen mnémotechnique : \emph{anode} et \emph{oxydation} commencent par une voyelle ; \emph{cathode} et \emph{réduction} commencent par une consonne.
\end{definition}

\courssub{Équations aux électrodes et bilan de fonctionnement}

\textbf{À l'anode (lame de zinc, pôle $\ominus$) : oxydation du zinc.}
Le zinc, réducteur le plus fort, cède deux électrons et passe en solution sous forme d'ions \ce{Zn^2+} :
\[ \ce{Zn(s) -> Zn^2+(aq) + 2e-} \]
Conséquence observable : la lame de zinc \textbf{s'amincit} progressivement et la concentration en \ce{Zn^2+} augmente dans le bécher de gauche.

\textbf{À la cathode (lame de cuivre, pôle $\oplus$) : réduction des ions cuivre(II).}
Les ions \ce{Cu^2+} (oxydant le plus fort) captent les électrons arrivant par le circuit extérieur et se déposent en cuivre métallique :
\[ \ce{Cu^2+(aq) + 2e- -> Cu(s)} \]
Conséquence observable : un \textbf{dépôt rougeâtre} de cuivre épaissit la lame, et la solution bleue de sulfate de cuivre \textbf{se décolore} peu à peu car $[\ce{Cu^2+}]$ diminue.

\textbf{Équation-bilan de fonctionnement.}
En additionnant les deux demi-équations (les électrons s'éliminent, deux échangés de part et d'autre) :
\[ \ce{Zn(s) + Cu^2+(aq) -> Zn^2+(aq) + Cu(s)} \]
C'est exactement la réaction \textbf{spontanée} prévue par la règle du gamma entre ces deux couples : la pile ne fait que \emph{canaliser} cette réaction naturelle à travers un circuit électrique. La réaction est \textbf{totale} ($K \gg 10^4$), lente au contact direct mais rapide et utile à travers la pile.

\courssub{Circulation des porteurs de charge}

Dans une pile en fonctionnement, le courant est assuré par des porteurs de charge différents selon la région du circuit :

\begin{itemize}
\item \textbf{Dans le circuit extérieur} (fils métalliques) : les \textbf{électrons} circulent de la lame de zinc (pôle $\ominus$) vers la lame de cuivre (pôle $\oplus$). Le \textbf{courant conventionnel} $I$ circule en sens inverse : du cuivre vers le zinc.
\item \textbf{Dans les solutions et le pont salin} : le courant est assuré par la \textbf{migration des ions}. Dans le pont salin, les cations \ce{K+} se dirigent vers le compartiment du \textbf{cuivre} (pour compenser la disparition des \ce{Cu^2+}), tandis que les anions \ce{Cl-} se dirigent vers le compartiment du \textbf{zinc} (pour compenser l'apparition des \ce{Zn^2+}).
\end{itemize}

\begin{popfigure}
\begin{center}
\begin{tikzpicture}[scale=0.9, >=Stealth]
% Solutions
\draw[fill=popblueL, draw=none] (-4.4,0) rectangle (-2.4,1.6);
\draw[thick, popdark] (-4.4,0) -- (-4.4,2.4) (-2.4,0) -- (-2.4,2.4) (-4.4,0) -- (-2.4,0);
\draw[fill=popblue!35, draw=none] (2.4,0) rectangle (4.4,1.6);
\draw[thick, popdark] (2.4,0) -- (2.4,2.4) (4.4,0) -- (4.4,2.4) (2.4,0) -- (4.4,0);
% Pont salin
\draw[thick, popgreen!70!black, fill=popgreenL] (-1.4,0.9) -- (-1.4,2.6) -- (-1.0,2.6) -- (-1.0,0.9) -- (1.0,0.9) -- (1.0,2.6) -- (1.4,2.6) -- (1.4,0.9) -- cycle;
\node[popgreen!60!black, font=\small\bfseries] at (0,2.95) {pont salin \ce{KCl}};
% Ions dans le pont
\node[poporange, font=\small\bfseries] at (0.55,1.75) {\ce{K+}};
\node[poporange, font=\small\bfseries] at (-0.15,1.45) {\ce{K+}};
\node[poppurple, font=\small\bfseries] at (-0.55,1.9) {\ce{Cl-}};
\node[poppurple, font=\small\bfseries] at (0.2,1.15) {\ce{Cl-}};
% Flèches de migration
\draw[->, very thick, poporange] (1.1,1.6) -- (2.6,1.6);
\node[poporange, font=\small] at (1.85,1.9) {\ce{K+}};
\draw[->, very thick, poppurple] (-1.1,1.2) -- (-2.6,1.2);
\node[poppurple, font=\small] at (-1.85,0.9) {\ce{Cl-}};
% Légendes compartiments
\node[font=\small, align=center] at (-3.4,2.75) {compartiment Zn\\ (excès de \ce{Zn^2+})};
\node[font=\small, align=center] at (3.4,2.75) {compartiment Cu\\ (déficit de \ce{Cu^2+})};
% Ions dans les solutions
\node[popblue, font=\small] at (-3.4,0.5) {\ce{Zn^2+}\; \ce{SO4^2-}};
\node[popblue, font=\small] at (3.4,0.5) {\ce{Cu^2+}\; \ce{SO4^2-}};
% Annotations
\node[font=\small, align=center, popdark] at (-3.4,-0.6) {les \ce{Cl-} neutralisent\\ l'excès de charge $+$};
\node[font=\small, align=center, popdark] at (3.4,-0.6) {les \ce{K+} remplacent\\ les \ce{Cu^2+} consommés};
\end{tikzpicture}
\end{center}
\legende{Migrations ioniques dans le pont salin : les cations \ce{K+} migrent vers le compartiment de la cathode (cuivre), les anions \ce{Cl-} vers le compartiment de l'anode (zinc), garantissant l'électroneutralité des deux solutions.}
\end{popfigure}

\begin{attention}[Erreur fréquente : inverser le sens des électrons]
Beaucoup d'élèves font circuler les électrons du cuivre vers le zinc. C'est faux ! Retiens cette règle absolue : \textbf{les électrons partent toujours de l'anode} (le pôle $\ominus$, là où a lieu l'oxydation, car l'oxydation \emph{libère} des électrons) \textbf{et arrivent à la cathode} (pôle $\oplus$). Dans la pile Daniell : électrons de Zn vers Cu. Le courant conventionnel, lui, va de Cu vers Zn : ne confonds pas les deux sens !
\end{attention}

\courssub{Représentation conventionnelle et f.é.m de la pile}

\begin{definition}[Représentation conventionnelle d'une pile]
Par convention, on écrit une pile en plaçant le \textbf{pôle négatif à gauche} et le \textbf{pôle positif à droite} ; une simple barre $\mid$ sépare deux phases (métal $\mid$ solution) et une \textbf{double barre} $\parallel$ symbolise le \textbf{pont salin} :
\[ \ominus \; \ce{Zn} \mid \ce{Zn^2+} \; \parallel \; \ce{Cu^2+} \mid \ce{Cu} \; \oplus \]
\end{definition}

\begin{definition}[Force électromotrice (f.é.m)]
La \cle{force électromotrice} $E$ d'une pile est la tension mesurée à ses bornes \textbf{en circuit ouvert} (la pile ne débite pas), à l'aide d'un voltmètre. Elle est égale à la différence entre les potentiels des deux électrodes :
\[ E = E_{\text{cathode}} - E_{\text{anode}} = E_{(+)} - E_{(-)} \]
\end{definition}

\begin{exemplebox}[Calcul de la f.é.m de la pile Daniell]
Dans les conditions standard (concentrations de \SI{1}{mol.L^{-1}}, température de \SI{25}{\celsius}) :
\begin{align*}
E &= E^\circ(\ce{Cu^2+}/\ce{Cu}) - E^\circ(\ce{Zn^2+}/\ce{Zn})\\
E &= 0,34 - (-0,76)\\
E &= \SI{1,10}{V}
\end{align*}
La valeur lue sur le voltmètre en circuit ouvert est bien voisine de \SI{1,10}{V}, ce qui confirme la prévision théorique issue de la classification électrochimique.
\end{exemplebox}

\begin{exoresolu}[Pôles et f.é.m d'une pile zinc--cuivre]
On réalise une pile en associant une demi-pile \ce{Zn} $\mid$ \ce{Zn^2+} (\SI{0,50}{mol.L^{-1}}) et une demi-pile \ce{Cu} $\mid$ \ce{Cu^2+} (\SI{0,50}{mol.L^{-1}}), reliées par un pont salin au \ce{KNO3}. On donne $E^\circ(\ce{Cu^2+}/\ce{Cu}) = \SI{+0,34}{V}$ et $E^\circ(\ce{Zn^2+}/\ce{Zn}) = \SI{-0,76}{V}$.
\begin{enumerate}
\item Déterminer les pôles de la pile et écrire sa représentation conventionnelle.
\item Écrire les équations aux électrodes et l'équation-bilan de fonctionnement.
\item Calculer sa f.é.m et indiquer le sens de migration des ions \ce{K+} et \ce{NO3-} du pont salin.
\end{enumerate}
\tcblower
\textbf{1. Polarité.} Le zinc est un réducteur plus fort que le cuivre ($E^\circ$ plus négatif) : c'est lui qui s'oxyde. La lame de zinc est le pôle $\ominus$ (anode), la lame de cuivre le pôle $\oplus$ (cathode). Représentation conventionnelle :
\[ \ominus \; \ce{Zn} \mid \ce{Zn^2+} \; \parallel \; \ce{Cu^2+} \mid \ce{Cu} \; \oplus \]
\textbf{2. Équations.} Anode (oxydation) : $\ce{Zn -> Zn^2+ + 2e-}$ ; cathode (réduction) : $\ce{Cu^2+ + 2e- -> Cu}$. Bilan :
\[ \ce{Zn(s) + Cu^2+(aq) -> Zn^2+(aq) + Cu(s)} \]
\textbf{3. F.é.m et migrations.} $E = 0,34 - (-0,76) = \SI{1,10}{V}$. Dans le pont salin, les cations \ce{K+} migrent vers le compartiment du \textbf{cuivre} (cathode) et les anions \ce{NO3-} vers le compartiment du \textbf{zinc} (anode), afin de maintenir l'électroneutralité des solutions.
\end{exoresolu}

\courssub{Usure et mort de la pile}

Une pile n'est pas éternelle : elle fonctionne tant que la réaction d'oxydoréduction n'a pas atteint son \textbf{état d'équilibre}. Au fil du fonctionnement :
\begin{itemize}
\item la lame de zinc \textbf{s'amincit} : le métal est consommé par l'oxydation ;
\item la concentration $[\ce{Cu^2+}]$ \textbf{diminue} : la solution bleue se décolore progressivement ;
\item la lame de cuivre \textbf{grossit} par dépôt de cuivre métallique ;
\item la f.é.m \textbf{diminue} lentement au cours du temps.
\end{itemize}

\begin{aretenir}[Quand une pile est-elle morte ?]
Une pile est \textbf{usée} (ou « morte ») lorsque la réaction de fonctionnement atteint son \textbf{état d'équilibre} : le quotient de réaction devient égal à la constante d'équilibre ($Q_r = K$). La f.é.m s'annule alors ($E = \SI{0}{V}$) : plus aucune force ne pousse les électrons dans le circuit. En pratique, la pile est souvent hors d'usage avant, lorsque l'un des réactifs (la lame de zinc ou les ions \ce{Cu^2+}) est quasiment épuisé.
\end{aretenir}

\begin{attention}[Pièges classiques au Baccalauréat TI]
\begin{itemize}
\item Ne jamais écrire le pont salin comme un simple fil : il ne laisse passer \textbf{que des ions}, jamais des électrons.
\item Ne pas confondre anode/cathode avec les signes dans un \emph{électrolyseur} : dans une \textbf{pile} (générateur), anode $=$ pôle $\ominus$ ; c'est l'inverse dans un récepteur électrolytique.
\item Dans l'équation-bilan, les électrons \textbf{n'apparaissent jamais} : ils doivent s'éliminer en additionnant les demi-équations.
\item La f.é.m se mesure \textbf{circuit ouvert} ; dès que la pile débite, la tension à ses bornes chute légèrement (résistance interne).
\end{itemize}
\end{attention}

\begin{aretenir}[L'essentiel sur la pile Daniell]
\begin{itemize}
\item Une pile transforme l'énergie d'une réaction redox \textbf{spontanée} en énergie électrique en \textbf{séparant} les deux couples dans deux demi-piles.
\item Pile Daniell : $\ominus \; \ce{Zn} \mid \ce{Zn^2+} \parallel \ce{Cu^2+} \mid \ce{Cu} \; \oplus$, de f.é.m standard $E = \SI{1,10}{V}$.
\item Anode ($\ominus$) : oxydation $\ce{Zn -> Zn^2+ + 2e-}$ ; cathode ($\oplus$) : réduction $\ce{Cu^2+ + 2e- -> Cu}$.
\item Électrons : de Zn vers Cu à l'extérieur ; ions \ce{K+} vers le cuivre et \ce{Cl-} vers le zinc dans le pont salin.
\item La pile meurt lorsque l'équilibre chimique est atteint ($E \to 0$).
\end{itemize}
\end{aretenir}

Maintenant que la pile Daniell n'a plus de secret pour toi, nous allons généraliser cette étude : comment prévoir la polarité, la f.é.m et l'équation de fonctionnement de \emph{n'importe quelle} pile à partir de la classification électrochimique ? C'est l'objet de la section suivante.

\coursec{Généralisation et prévision des piles}

Dans la section précédente, nous avons étudié en détail la pile Daniell : deux demi-piles $\ce{Zn^{2+}}/\ce{Zn}$ et $\ce{Cu^{2+}}/\ce{Cu}$ reliées par un pont salin, un courant qui circule spontanément du pôle $+$ (cuivre) vers le pôle $-$ (zinc) dans le circuit extérieur, et une f.é.m. d'environ \SI{1,10}{V}. Une question naturelle se pose alors : la pile Daniell est-elle un cas isolé, un coup de chance de l'histoire des sciences ? La réponse est non, et c'est toute la puissance de la classification électrochimique : \textbf{toute réaction d'oxydoréduction spontanée peut être exploitée pour construire une pile}. Cette section te donne la méthode complète pour concevoir, représenter et prévoir n'importe quelle pile.

\courssub{Le principe général : séparer les deux couples d'une réaction spontanée}

\courssub{De la réaction spontanée à la pile}

Reprenons l'expérience fondatrice : quand on plonge une lame de zinc dans une solution de sulfate de cuivre, la réaction spontanée

\[ \ce{Zn(s) + Cu^{2+}(aq) -> Zn^{2+}(aq) + Cu(s)} \]

se produit \emph{directement} : les électrons passent du zinc aux ions cuivre par contact, et l'énergie chimique est intégralement dégradée en chaleur. La règle du gamma (section 1) nous avait déjà permis de prévoir cette spontanéité : le zinc, réducteur le plus fort, réduit l'oxydant de l'autre couple.

L'idée géniale de la pile consiste à \textbf{séparer physiquement les deux couples} dans deux demi-piles distinctes et à n'autoriser le passage des électrons que par un circuit extérieur. Au lieu d'être perdue en chaleur, l'énergie de la réaction devient alors de l'\cle{énergie électrique} utilisable : c'est le principe de toutes les piles.

\begin{propriete}[Principe de construction d'une pile]
Soient deux couples redox $\mathrm{Ox_1}/\mathrm{Red_1}$ et $\mathrm{Ox_2}/\mathrm{Red_2}$ tels qu'une réaction spontanée soit prévue par la règle du gamma. On construit une pile en :
\begin{itemize}
\item réalisant deux \cle{demi-piles}, chacune contenant un couple ;
\item les reliant par un \cle{pont salin} (ou une paroi poreuse) qui assure la neutralité électrique des solutions sans les mélanger ;
\item reliant les électrodes par un circuit extérieur : les électrons circulent alors de l'\cle{anode} (pôle $-$, siège de l'oxydation) vers la \cle{cathode} (pôle $+$, siège de la réduction).
\end{itemize}
La pile fonctionne tant que la réaction spontanée peut se poursuivre ; elle est \cle{usée} lorsque l'équilibre chimique est atteint.
\end{propriete}

\begin{methode}[Concevoir une pile à partir de deux couples]
\begin{enumerate}
\item \textbf{Repérer les deux couples} et leurs potentiels standard $E^\circ$ dans la classification électrochimique.
\item \textbf{Appliquer la règle du gamma} : le réducteur du couple de plus bas potentiel s'oxyde (c'est l'anode, pôle $-$) ; l'oxydant du couple de plus haut potentiel se réduit (c'est la cathode, pôle $+$).
\item \textbf{Écrire les demi-équations} aux électrodes, puis l'équation-bilan de fonctionnement.
\item \textbf{Calculer la f.é.m.} : $E = E^\circ_{\text{cathode}} - E^\circ_{\text{anode}} = E^\circ_{+} - E^\circ_{-}$.
\item \textbf{Dessiner le schéma} ou écrire la représentation conventionnelle.
\end{enumerate}
\end{methode}

\courssub{La représentation conventionnelle d'une pile}

Pour décrire une pile sans la dessiner, les chimistes utilisent une écriture symbolique normalisée, très souvent exigée au Baccalauréat TI.

\begin{definition}[Représentation conventionnelle]
La représentation conventionnelle d'une pile s'écrit :
\[ \underset{\text{pôle } -}{\underbrace{\mathrm{Red_1} \mid \mathrm{Ox_1}}} \;\; \Vert \;\; \underset{\text{pôle } +}{\underbrace{\mathrm{Ox_2} \mid \mathrm{Red_2}}} \]
\begin{itemize}
\item le \textbf{pôle négatif} (anode) est écrit \textbf{à gauche}, le \textbf{pôle positif} (cathode) \textbf{à droite} ;
\item la simple barre $\mid$ symbolise la séparation entre deux phases (métal / solution) ;
\item la double barre $\Vert$ symbolise le \cle{pont salin}.
\end{itemize}
\end{definition}

\begin{attention}[Sens d'écriture dans chaque demi-pile]
Dans la notation conventionnelle, l'ordre des espèces dépend du pôle :
\begin{itemize}
\item \textbf{Pôle $-$ (anode, à gauche) :} on écrit le \textbf{réducteur puis l'oxydant} : $\ce{Red} \mid \ce{Ox}$ (ex. $\ce{Zn} \mid \ce{Zn^{2+}}$).
\item \textbf{Pôle $+$ (cathode, à droite) :} on écrit l'\textbf{oxydant puis le réducteur} : $\ce{Ox} \mid \ce{Red}$ (ex. $\ce{Ag^{+}} \mid \ce{Ag}$).
\end{itemize}
Erreur classique au Bac : écrire $\ce{Zn^{2+}} \mid \ce{Zn}$ ou placer le pôle $+$ à gauche. Retiens : \emph{le moins à gauche, le plus à droite, comme sur une ligne de classification qui monte}.
\end{attention}

\courssub{Étude complète : la pile $\ce{Fe} \mid \ce{Fe^{2+}} \Vert \ce{Ag^{+}} \mid \ce{Ag}$}

Appliquons la méthode à deux nouveaux couples : $\ce{Fe^{2+}}/\ce{Fe}$ ($E^\circ = \SI{-0,44}{V}$) et $\ce{Ag^{+}}/\ce{Ag}$ ($E^\circ = \SI{+0,80}{V}$).

\textbf{Étape 1 — Polarité.} Le fer a le potentiel le plus bas : c'est le réducteur le plus fort. La règle du gamma prévoit donc que le fer s'oxyde : la lame de fer est l'\textbf{anode, pôle $-$}. L'argent est la \textbf{cathode, pôle $+$}.

\textbf{Étape 2 — Demi-équations.}
\begin{align*}
\text{Anode } (-) : \quad & \ce{Fe(s) -> Fe^{2+}(aq) + 2e-} \quad \text{(oxydation, la lame de fer s'amincit)} \\
\text{Cathode } (+) : \quad & \ce{Ag^{+}(aq) + e- -> Ag(s)} \quad \text{(réduction, dépôt gris d'argent)}
\end{align*}

\textbf{Étape 3 — Équation-bilan.} On multiplie la demi-équation de réduction par 2 pour équilibrer les électrons (qui ne doivent jamais apparaître dans le bilan) :
\[ \ce{Fe(s) + 2Ag^{+}(aq) -> Fe^{2+}(aq) + 2Ag(s)} \]

\textbf{Étape 4 — F.é.m.}
\[ E = E^\circ_{+} - E^\circ_{-} = 0{,}80 - (-0{,}44) = \SI{1,24}{V} \]

\begin{popfigure}
\begin{center}
\begin{tikzpicture}[scale=0.95, every node/.style={font=\small}]
% Bécher gauche (Fe)
\draw[thick, popdark] (-3.2,0) -- (-3.2,2.6) (-0.6,0) -- (-0.6,2.6) (-3.2,0) -- (-0.6,0);
\fill[popblue!25] (-3.15,0.05) rectangle (-0.65,1.6);
\node[popblue] at (-1.9,0.8) {$\ce{Fe^{2+}}$};
% Électrode Fe
\fill[popmuted!70] (-2.15,1.1) rectangle (-1.75,3.4);
\node[above] at (-1.95,3.4) {$\ce{Fe}$};
\node at (-1.95,3.85) {\textbf{pôle $-$}};
% Bécher droit (Ag)
\draw[thick, popdark] (0.6,0) -- (0.6,2.6) (3.2,0) -- (3.2,2.6) (0.6,0) -- (3.2,0);
\fill[popgreen!20] (0.65,0.05) rectangle (3.15,1.6);
\node[popgreen!60!black] at (1.9,0.8) {$\ce{Ag^{+}}$};
% Électrode Ag
\fill[popmuted!40] (1.65,1.1) rectangle (2.05,3.4);
\node[above] at (1.85,3.4) {$\ce{Ag}$};
\node at (1.85,3.85) {\textbf{pôle $+$}};
% Pont salin (U-tube plongeant dans les solutions)
\draw[thick, poporange, fill=poporange!25] 
    (-1.0,1.5) -- (-1.0,0.5) .. controls (-1.0,0.0) and (1.0,0.0) .. (1.0,0.5) -- (1.0,1.5)
    -- (1.15,1.5) -- (1.15,0.5) .. controls (1.15,0.0) and (-1.15,0.0) .. (-1.15,0.5) -- (-1.15,1.5) -- cycle;
\node[poporange!80!black] at (0,2.75) {pont salin};
% Circuit extérieur
\draw[thick, popdark] (-1.95,3.4) -- (-1.95,4.5) -- (1.85,4.5) -- (1.85,3.4);
\draw[thick, poppink, fill=poppink!20] (-0.35,4.35) rectangle (0.35,4.65);
\node at (0,4.5) {\textbf{V}};
% Sens du courant et des électrons
\draw[-{Stealth[length=3mm]}, thick, poppurple] (-1.6,4.9) -- (-0.5,4.9);
\node[poppurple, above] at (-1.05,4.9) {$e^-$};
\draw[-{Stealth[length=3mm]}, thick, popgold] (0.5,4.15) -- (1.6,4.15);
\node[popgold!80!black, below] at (1.05,4.15) {$I$};
% F.é.m.
\node[font=\bfseries, popink] at (0,-0.7) {$E = \SI{1,24}{V}$};
\end{tikzpicture}
\end{center}
\legende{Pile fer--argent : le fer (anode, pôle $-$) s'oxyde et libère des électrons qui circulent vers l'électrode d'argent (cathode, pôle $+$), où les ions $\ce{Ag^{+}}$ sont réduits en dépôt d'argent. Le pont salin (U-tube) plonge dans les deux solutions.}
\end{popfigure}

\begin{exemplebox}[Pile zinc--argent : calcul de la f.é.m.]
On associe les couples $\ce{Zn^{2+}}/\ce{Zn}$ ($E^\circ = \SI{-0,76}{V}$) et $\ce{Ag^{+}}/\ce{Ag}$ ($E^\circ = \SI{+0,80}{V}$). La règle du gamma impose : zinc anode ($-$), argent cathode ($+$).
\[ E = E^\circ_{+} - E^\circ_{-} = 0{,}80 - (-0{,}76) = \SI{1,56}{V} \]
Équation-bilan : $\ce{Zn(s) + 2Ag^{+}(aq) -> Zn^{2+}(aq) + 2Ag(s)}$.
Représentation : $\ce{Zn} \mid \ce{Zn^{2+}} \Vert \ce{Ag^{+}} \mid \ce{Ag}$.
\end{exemplebox}

\courssub{Piles à électrode inerte : le cas des couples sans métal}

Certains couples redox ne comportent \textbf{aucun métal solide} pouvant servir d'électrode : c'est le cas de $\ce{Fe^{3+}}/\ce{Fe^{2+}}$, $\ce{MnO4-}/\ce{Mn^{2+}}$ ou $\ce{Cr2O7^2-}/\ce{Cr^{3+}}$, où toutes les espèces sont en solution. Comment alors capter les électrons ?

On utilise une \cle{électrode inerte} : un fil ou une lame de \textbf{platine} (Pt), métal précieux qui conduit le courant \emph{sans participer à la réaction}. Il joue uniquement le rôle de « collecteur d'électrons ».

\begin{definition}[Demi-pile à électrode inerte]
Pour un couple dont les deux membres sont en solution, la demi-pile est constituée d'une électrode de platine plongeant dans une solution contenant les deux espèces du couple. Sa représentation conventionnelle est :
\[ \ce{Pt} \mid \ce{Fe^{2+}}, \ce{Fe^{3+}} \qquad \text{ou} \qquad \ce{Fe^{3+}}, \ce{Fe^{2+}} \mid \ce{Pt} \]
selon que la demi-pile constitue le pôle $-$ (à gauche) ou le pôle $+$ (à droite). Les deux espèces en solution sont séparées par une virgule.
\end{definition}

\begin{exemplebox}[Pile fer--ions fer(III)]
Associons $\ce{Zn^{2+}}/\ce{Zn}$ ($\SI{-0,76}{V}$) et $\ce{Fe^{3+}}/\ce{Fe^{2+}}$ ($\SI{+0,77}{V}$). Le zinc est l'anode ; la demi-pile $\ce{Fe^{3+}}/\ce{Fe^{2+}}$ nécessite du platine :
\[ \ce{Zn} \mid \ce{Zn^{2+}} \Vert \ce{Fe^{3+}}, \ce{Fe^{2+}} \mid \ce{Pt} \qquad E = 0{,}77 - (-0{,}76) = \SI{1,53}{V} \]
À la cathode de platine : $\ce{Fe^{3+} + e- -> Fe^{2+}}$. Le platine ne s'use pas : seules les concentrations en ions fer évoluent.
\end{exemplebox}

\courssub{F.é.m. et écart des potentiels : plus les couples sont éloignés, plus la pile est « forte »}

La formule $E = E^\circ_{+} - E^\circ_{-}$ a une lecture graphique immédiate sur la classification électrochimique : \textbf{la f.é.m. est la distance verticale entre les deux couples} sur l'axe des potentiels. Plus les couples sont éloignés, plus la f.é.m. est grande, et plus la réaction de fonctionnement est quantitative (constante d'équilibre élevée).

\begin{popfigure}
\begin{center}
\begin{tikzpicture}
\begin{axis}[
    width=0.85\linewidth, height=6.5cm,
    xlabel={Écart de potentiel standard $\Delta E^\circ = E^\circ_{+} - E^\circ_{-}$ (V)},
    ylabel={F.é.m. $E$ (V)},
    xmin=0, xmax=3.5, ymin=0, ymax=3.5,
    grid=both, grid style={popline!40},
    legend style={at={(0.03,0.97)}, anchor=north west, font=\small},
    xtick={0,0.5,...,3.5}, ytick={0,0.5,...,3.5},
]
\addplot[popcurve, domain=0:3.5, samples=50]{x};
\addlegendentry{$E = \Delta E^\circ$}
\addplot[only marks, mark=*, mark size=2.5pt, poporange] coordinates {
    (1.10,1.10) (1.24,1.24) (1.56,1.56) (1.53,1.53) (3.17,3.17)
};
\node[font=\footnotesize, poporange!80!black, anchor=west] at (axis cs:1.10,1.02) {Zn/Cu};
\node[font=\footnotesize, poporange!80!black, anchor=west] at (axis cs:1.24,1.16) {Fe/Ag};
\node[font=\footnotesize, poporange!80!black, anchor=west] at (axis cs:1.56,1.48) {Zn/Ag};
\node[font=\footnotesize, poporange!80!black, anchor=west] at (axis cs:1.53,1.45) {Zn/Fe$^{3+}$};
\node[font=\footnotesize, poporange!80!black, anchor=west] at (axis cs:3.17,3.09) {Mg/Ag};
\end{axis}
\end{tikzpicture}
\end{center}
\legende{La f.é.m. d'une pile est égale à l'écart des potentiels standard des deux couples : les points représentatifs des piles usuelles s'alignent sur la droite $E = \Delta E^\circ$.}
\end{popfigure}

\begin{aretenir}[Deux conséquences pratiques]
\begin{itemize}
\item Pour obtenir une \textbf{f.é.m. élevée}, on choisit deux couples très éloignés dans la classification : un réducteur très fort (métal alcalin, magnésium, zinc) face à un oxydant très fort. C'est la stratégie des fabricants de piles.
\item Pour obtenir une \textbf{tension élevée}, on monte plusieurs éléments \textbf{en série} : les f.é.m. s'additionnent.
\end{itemize}
\end{aretenir}

\courssub{Influence des concentrations : pourquoi une pile s'use}

Jusqu'ici, nous avons raisonné avec les \emph{potentiels standard}, mesurés lorsque toutes les concentrations valent \SI{1}{mol.L^{-1}}. En réalité, le potentiel de chaque couple \textbf{dépend des concentrations} : le potentiel d'un couple augmente quand la concentration de l'oxydant augmente, et diminue quand c'est le réducteur qui se concentre (loi de Nernst, approfondie en S).

Suivons le fonctionnement de la pile Daniell :
\begin{itemize}
\item à l'anode, le zinc s'oxyde : la concentration en $\ce{Zn^{2+}}$ \textbf{augmente}, donc $E(\ce{Zn^{2+}}/\ce{Zn})$ augmente légèrement ;
\item à la cathode, les ions $\ce{Cu^{2+}}$ sont consommés : leur concentration \textbf{diminue}, donc $E(\ce{Cu^{2+}}/\ce{Cu})$ diminue.
\end{itemize}
Les deux potentiels se rapprochent : la f.é.m. $E = E_{+} - E_{-}$ \textbf{diminue progressivement} au cours du fonctionnement. Lorsque les concentrations n'évoluent plus, c'est que l'équilibre chimique est atteint : les deux potentiels sont égaux, la f.é.m. est nulle, et la pile est \textbf{usée}.

\begin{propriete}[Pile usée et équilibre chimique (complément Bac)]
Une pile est usée lorsque le système atteint l'équilibre. Alors :
\begin{itemize}
\item la f.é.m. s'annule : $E = \SI{0}{V}$ ;
\item le quotient de réaction devient égal à la constante d'équilibre : $Q_r = K$.
\end{itemize}
Pour la pile Daniell, à l'équilibre : $\dfrac{[\ce{Zn^{2+}}]}{[\ce{Cu^{2+}}]} = K \approx 10^{37}$ : la réaction est totale, pratiquement tous les ions $\ce{Cu^{2+}}$ ont disparu. Une pile ne se « recharge » donc pas d'elle-même : il faudrait inverser la réaction en fournissant de l'énergie (c'est l'électrolyse, principe des accumulateurs).
\end{propriete}

\begin{attention}[Pièges fréquents sur la f.é.m.]
\begin{itemize}
\item La f.é.m. \textbf{ne dépend pas} de la surface des électrodes ni de la quantité de métal : une grosse pile et une petite pile construites avec les mêmes couples ont la même f.é.m. (mais pas la même capacité !).
\item Ne jamais écrire $E = E^\circ_{-} - E^\circ_{+}$ : une f.é.m. est toujours \textbf{positive}.
\item Une pile usée n'est pas « vide d'électrons » : c'est l'\emph{équilibre chimique} qui est atteint.
\end{itemize}
\end{attention}

\courssub{Applications : les piles et accumulateurs de la vie quotidienne}

Le principe général que tu viens d'apprendre se retrouve dans tous les générateurs électrochimiques qui t'entourent.

\begin{center}
\begin{tabularx}{\linewidth}{|l|X|X|}
\hline
\rowcolor{popblueL} \textbf{Générateur} & \textbf{Couples / principe} & \textbf{Caractéristiques} \\
\hline
Pile saline (Leclanché) & $\ce{Zn^{2+}}/\ce{Zn}$ et couple du dioxyde de manganèse $\ce{MnO2}$ & Environ \SI{1,5}{V} ; lampes torches, télécommandes \\
\hline
Pile alcaline & Mêmes couples qu'en milieu basique (KOH) & \SI{1,5}{V}, plus grande capacité et durée de vie \\
\hline
Accumulateur au plomb & $\ce{PbSO4}/\ce{Pb}$ et $\ce{PbO2}/\ce{PbSO4}$ dans $\ce{H2SO4}$ & \SI{2}{V} par élément ; \textbf{rechargeable} \\
\hline
Batterie lithium-ion & Intercalation du lithium entre deux électrodes & Environ \SI{3,7}{V} par élément ; téléphones, ordinateurs \\
\hline
\end{tabularx}
\end{center}

\begin{savaistu}[L'accumulateur au plomb de ta voiture]
La batterie d'une voiture — indispensable pour démarrer les véhicules qui sillonnent Douala ou Yaoundé — délivre une tension de \SI{12}{V}. Or un seul élément plomb/acide ne fournit que \SI{2}{V}. L'astuce : \textbf{six éléments de \SI{2}{V} montés en série}, dont les f.é.m. s'additionnent ($6 \times 2 = \SI{12}{V}$). Contrairement à une pile saline, l'accumulateur est \emph{rechargeable} : l'alternateur force le courant en sens inverse et régénère les réactifs par électrolyse. Quant aux batteries lithium-ion des téléphones portables, elles exploitent le lithium, métal le plus léger et réducteur très fort : d'où leur excellente f.é.m. par élément (\SI{3,7}{V}) et leur grande densité d'énergie.
\end{savaistu}

\begin{exoresolu}[Conception complète d'une pile cuivre--argent]
On dispose des couples $\ce{Cu^{2+}}/\ce{Cu}$ ($E^\circ = \SI{+0,34}{V}$) et $\ce{Ag^{+}}/\ce{Ag}$ ($E^\circ = \SI{+0,80}{V}$).
\begin{enumerate}
\item Prévoir la réaction spontanée et identifier les pôles.
\item Écrire la représentation conventionnelle de la pile.
\item Calculer sa f.é.m. et écrire l'équation-bilan de fonctionnement.
\end{enumerate}
\tcblower
\textbf{1.} Les deux potentiels sont proches, mais $E^\circ(\ce{Ag^{+}}/\ce{Ag}) > E^\circ(\ce{Cu^{2+}}/\ce{Cu})$. La règle du gamma impose : le cuivre (réducteur le plus fort) s'oxyde, les ions $\ce{Ag^{+}}$ sont réduits. Donc : \textbf{cuivre = anode, pôle $-$} ; \textbf{argent = cathode, pôle $+$}.

\textbf{2.} Pôle $-$ à gauche, pôle $+$ à droite, double barre pour le pont salin :
\[ \ce{Cu} \mid \ce{Cu^{2+}} \Vert \ce{Ag^{+}} \mid \ce{Ag} \]

\textbf{3.} F.é.m. :
\[ E = E^\circ_{+} - E^\circ_{-} = 0{,}80 - 0{,}34 = \SI{0,46}{V} \]
Demi-équations et bilan (on double la réduction de l'argent) :
\begin{align*}
(-) : \quad & \ce{Cu(s) -> Cu^{2+}(aq) + 2e-} \\
(+) : \quad & \ce{Ag^{+}(aq) + e- -> Ag(s)} \quad (\times 2) \\
\hline
& \ce{Cu(s) + 2Ag^{+}(aq) -> Cu^{2+}(aq) + 2Ag(s)}
\end{align*}
Remarque : la f.é.m. est faible (\SI{0,46}{V}) car les deux couples sont \emph{proches} dans la classification — conformément au lien entre f.é.m. et écart des potentiels.
\end{exoresolu}

\begin{exercice}[À toi de jouer]
On réalise une pile avec les couples $\ce{Mg^{2+}}/\ce{Mg}$ ($E^\circ = \SI{-2,37}{V}$) et $\ce{Pb^{2+}}/\ce{Pb}$ ($E^\circ = \SI{-0,13}{V}$).
\begin{enumerate}
\item Écris la représentation conventionnelle de la pile en précisant les pôles.
\item Écris les demi-équations aux électrodes et l'équation-bilan.
\item Calcule la f.é.m. de la pile.
\item La lame de magnésium voit-elle sa masse augmenter ou diminuer ? Justifie.
\end{enumerate}
\end{exercice}

\begin{corrige}[Exercice : pile magnésium--plomb]
\textbf{1.} $E^\circ(\ce{Mg^{2+}}/\ce{Mg}) < E^\circ(\ce{Pb^{2+}}/\ce{Pb})$ : le magnésium est le réducteur le plus fort, c'est l'anode (pôle $-$). Représentation :
\[ \ce{Mg} \mid \ce{Mg^{2+}} \Vert \ce{Pb^{2+}} \mid \ce{Pb} \]
\textbf{2.} Anode : $\ce{Mg(s) -> Mg^{2+}(aq) + 2e-}$ ; cathode : $\ce{Pb^{2+}(aq) + 2e- -> Pb(s)}$. Les électrons s'éliminent directement :
\[ \ce{Mg(s) + Pb^{2+}(aq) -> Mg^{2+}(aq) + Pb(s)} \]
\textbf{3.} $E = E^\circ_{+} - E^\circ_{-} = -0{,}13 - (-2{,}37) = \SI{2,24}{V}$. F.é.m. élevée : les couples sont très éloignés.
\textbf{4.} La lame de magnésium est le siège d'une \emph{oxydation} : elle perd des atomes qui passent en solution sous forme $\ce{Mg^{2+}}$. Sa masse \textbf{diminue} (elle se consomme), tandis que la lame de plomb grossit par dépôt métallique.
\end{corrige}

\begin{aretenir}[L'essentiel de la section]
\begin{itemize}
\item \textbf{Toute réaction redox spontanée} (règle du gamma) peut être transformée en pile en séparant les deux couples en deux demi-piles reliées par un pont salin.
\item Représentation conventionnelle : \textbf{pôle $-$ à gauche, pôle $+$ à droite}, double barre $\Vert$ pour le pont salin ; à l'anode (gauche) on écrit $\ce{Red} \mid \ce{Ox}$, à la cathode (droite) $\ce{Ox} \mid \ce{Red}$.
\item F.é.m. : $E = E^\circ_{+} - E^\circ_{-}$, toujours positive ; elle est d'autant plus grande que les couples sont éloignés dans la classification.
\item Couple sans métal $\Rightarrow$ électrode inerte de platine : $\ce{Pt} \mid \ce{Fe^{2+}}, \ce{Fe^{3+}}$ (anode) ou $\ce{Fe^{3+}}, \ce{Fe^{2+}} \mid \ce{Pt}$ (cathode).
\item La f.é.m. dépend des concentrations et \textbf{diminue en fonctionnement} ; pile usée $\Leftrightarrow$ équilibre atteint, $E = 0$ et $Q_r = K$.
\item Applications : piles salines et alcalines (\SI{1,5}{V}), accumulateur au plomb (\SI{12}{V} = 6 éléments de \SI{2}{V} en série), batteries lithium-ion (\SI{3,7}{V} par élément).
\end{itemize}
\end{aretenir}

\coursec{Corrosion des métaux : manifestations et mécanisme}

Dans la section précédente, nous avons appris à \emph{prévoir} et à \emph{construire} des piles : en associant deux couples redox bien choisis, séparés par un pont salin, nous transformons une réaction d'oxydoréduction spontanée en source de courant électrique utile. Mais la nature, elle, ne demande pas notre permission : dès qu'un métal réducteur se trouve au contact d'un oxydant de son environnement (le dioxygène de l'air, l'eau), une réaction d'oxydoréduction spontanée peut se déclencher \emph{contre notre volonté}. C'est exactement ce qui se passe quand un portail en fer rouille à Douala, quand une coque de pirogue s'use à Limbe ou quand les pipelines de l'oléoduc Tchad--Cameroun doivent être surveillés en permanence. Cette réaction non désirée porte un nom : la \cle{corrosion}. Comprendre la corrosion, c'est comprendre que des millions de micropiles de type Daniell se forment spontanément à la surface des métaux — et c'est la première étape pour s'en protéger.

\courssub{Définition et manifestations de la corrosion}

\textbf{Une réalité quotidienne}\par
Observe autour de toi : le fer d'une vieille marmite abandonnée sous la pluie se recouvre d'une couche brun-rouge friable ; les bijoux en argent noircissent avec le temps ; les toitures et statues en cuivre verdissent après des années d'exposition. Dans les trois cas, le métal \emph{brillant} a disparu au profit d'un composé \emph{terne} : le métal a été oxydé. Ce n'est pas une simple salissure que l'on peut essuyer : le métal lui-même a été \emph{consommé} par une réaction chimique.

\begin{definition}[Corrosion]
La \cle{corrosion} est la dégradation progressive d'un métal (ou d'un alliage) par une réaction d'\cle{oxydoréduction} avec les espèces présentes dans son environnement : dioxygène de l'air, eau (surtout si elle est salée ou acide), sol humide, etc. Au cours de la corrosion, le métal $M$ joue le rôle de \textbf{réducteur} et est oxydé en cations :
\[ \mathrm{M} \rightarrow \mathrm{M}^{n+} + n\,\mathrm{e}^- \]
L'oxydant le plus fréquent est le dioxygène \ce{O2}, réduit en présence d'eau selon :
\[ \ce{O2 + 2 H2O + 4 e- -> 4 HO-} \]
\end{definition}

\begin{attention}[Corrosion ou usure mécanique ?]
La corrosion est un phénomène \textbf{chimique} : le métal est transformé en un nouveau composé (oxyde, hydroxyde, sel). Une rayure, une abrasion ou une déformation ne sont \emph{pas} de la corrosion, même si elles peuvent la favoriser en exposant du métal frais. De même, la corrosion n'est pas une dissolution simple : le fer ne « fond » pas dans l'eau, il est \emph{oxydé}.
\end{attention}

\textbf{Les manifestations les plus courantes}\par
Chaque métal laisse une « signature » visible de sa corrosion :

\begin{center}
\begin{tabularx}{\linewidth}{|l|X|X|}
\hline
\rowcolor{popblueL} \textbf{Métal} & \textbf{Produit de corrosion} & \textbf{Aspect} \\
\hline
Fer, acier & Oxyde de fer(III) hydraté \ce{Fe2O3.xH2O} : la \cle{rouille} & Couche brun-rouge, poreuse et friable \\
\hline
Argent & Sulfure d'argent \ce{Ag2S} & Ternissement noir (dû aux traces de \ce{H2S} de l'air) \\
\hline
Cuivre, bronze & Carbonates et hydroxydes basiques de cuivre(II) : le \cle{vert-de-gris} & Couche verte (toitures, statues) \\
\hline
Aluminium & Oxyde d'aluminium \ce{Al2O3} (alumine) & Couche grise, fine, invisible, \emph{protectrice} \\
\hline
\end{tabularx}
\end{center}

\begin{savaistu}[Le paradoxe de l'aluminium]
D'après la classification électrochimique, l'aluminium ($E^\circ(\ce{Al^{3+}}/\ce{Al}) \approx \SI{-1,66}{V}$) est un réducteur \emph{bien plus fort} que le fer ($E^\circ(\ce{Fe^{2+}}/\ce{Fe}) \approx \SI{-0,44}{V}$). Il devrait donc se corroder plus vite que le fer ! Pourtant, les marmites et les tôles en aluminium durent des années. L'explication : dès qu'il est exposé à l'air, l'aluminium se recouvre instantanément d'une couche d'alumine \ce{Al2O3} d'environ \SI{5}{nm} d'épaisseur, \textbf{imperméable et adhérente}, qui isole le métal de son environnement. C'est le phénomène de \cle{passivation}. La rouille, elle, ne protège pas : nous verrons pourquoi.
\end{savaistu}

\courssub{Conditions nécessaires à la corrosion du fer}

\textbf{Une expérience classique et décisive}\par
Le fer rouille-t-il dans l'eau seule ? Dans l'air sec seul ? L'expérience des trois clous répond clairement à ces questions.

\begin{experience}[Les trois clous : conditions de la corrosion du fer]
\textbf{Dispositif et protocole} : on place trois clous en fer identiques dans trois tubes :
\begin{itemize}
\item \textbf{Tube 1} : le clou est partiellement immergé dans de l'eau du robinet, le tube reste ouvert à l'air ;
\item \textbf{Tube 2} : le clou est totalement immergé dans de l'eau \emph{bouillie} (pour chasser le dioxygène dissous), recouverte d'une couche d'huile qui empêche l'air de redissoudre du \ce{O2} ;
\item \textbf{Tube 3} : le clou est placé dans un tube \emph{sec}, bouché, contenant un dessiccant (chlorure de calcium anhydre) qui absorbe toute trace d'humidité.
\end{itemize}
\textbf{Observations après quelques jours} :
\begin{itemize}
\item Tube 1 : le clou est \textbf{rouillé}, surtout à la surface de contact eau/air ;
\item Tube 2 : le clou est \textbf{intact} (eau sans dioxygène) ;
\item Tube 3 : le clou est \textbf{intact} (air sans eau).
\end{itemize}
\textbf{Interprétation} : la corrosion du fer exige la présence \emph{simultanée} d'\textbf{eau} et de \textbf{dioxygène}. L'eau seule ou l'air sec seul ne suffisent pas. C'est pourquoi la rouille apparaît d'abord à la ligne de flottaison d'un piquet planté dans l'eau, et pourquoi le climat chaud et humide du littoral camerounais (Douala, Limbe, Kribi) est particulièrement agressif pour les structures métalliques.
\end{experience}

\begin{aretenir}[Conditions de la corrosion du fer]
Le fer ne se corrode que si \textbf{deux conditions sont réunies} : présence d'\cle{eau} (même un simple film d'humidité) \emph{et} présence de \cle{dioxygène}. Supprimer l'un des deux facteurs stoppe la corrosion : c'est le principe de nombreuses méthodes de protection.
\end{aretenir}

\courssub{Mécanisme de la corrosion : les micropiles}

\textbf{L'idée-clé : une pile Daniell miniature et involontaire}\par
Pourquoi l'eau est-elle indispensable ? Parce qu'elle joue le rôle d'\textbf{électrolyte} : elle conduit le courant ionique, exactement comme les solutions et le pont salin de la pile Daniell. La surface d'une pièce de fer n'est jamais parfaitement homogène : elle contient des impuretés (grains de carbone dans l'acier), des défauts, des rayures, des zones déformées. Entre ces zones, de minuscules différences de potentiel apparaissent : la surface du métal se comporte alors comme une multitude de \textbf{micropiles} en court-circuit.

\begin{definition}[Micropile de corrosion]
Une \cle{micropile} est une micro-pile électrochimique qui se forme spontanément à la surface d'un métal recouvert d'un film d'eau : certaines zones (défauts, zones tendues) jouent le rôle d'\textbf{anode} où le métal s'oxyde, tandis que d'autres zones (impuretés, zones riches en \ce{O2}) jouent le rôle de \textbf{cathode} où le dioxygène dissous se réduit. Le métal lui-même assure la conduction électronique, et le film d'eau (surtout s'il contient des sels) assure la conduction ionique : le circuit de la pile est bouclé.
\end{definition}

\textbf{Les deux demi-équations de la micropile}\par
Considérons une goutte d'eau salée déposée sur une plaque de fer.

\begin{itemize}
\item \textbf{Zones anodiques} (pôle $-$ des micropiles, souvent au centre de la goutte où \ce{O2} est moins disponible, ou près des défauts) : le fer est oxydé et passe en solution :
\[ \ce{Fe -> Fe^{2+} + 2 e-} \]
Le métal se \emph{dissout} littéralement : c'est là que le fer disparaît.
\item \textbf{Zones cathodiques} (pôle $+$, souvent au bord de la goutte, bien aéré) : le dioxygène dissous est réduit en ions hydroxyde :
\[ \ce{O2 + 2 H2O + 4 e- -> 4 HO-} \]
\end{itemize}

Les électrons circulent \emph{dans le métal} de la zone anodique vers la zone cathodique ; les ions migrent \emph{dans la goutte d'eau}. On retrouve exactement l'architecture de la pile Daniell, mais à l'échelle microscopique et sans aucune utilité : l'énergie libérée sert uniquement à détruire le métal.

\begin{popfigure}
\begin{center}
\begin{tikzpicture}[scale=1.0, every node/.style={font=\small}]
% Plaque de fer
\fill[popmuted!60] (-5.5,-1.4) rectangle (5.5,-0.6);
\draw[popdark, thick] (-5.5,-0.6) -- (5.5,-0.6);
\node[popdark] at (0,-1.05) {\textbf{Plaque de fer} (conduction électronique)};
% Goutte d'eau
\fill[popblue!25, opacity=0.85] (0,-0.6) arc[start angle=180, end angle=360, x radius=3.2, y radius=1.9] -- cycle;
\draw[popblue, thick] (0,-0.6) arc[start angle=180, end angle=360, x radius=3.2, y radius=1.9];
\node[popblue!70!popdark] at (0,0.55) {Goutte d'eau salée (\ce{Na+}, \ce{Cl-}, \ce{O2} dissous)};
% Zone anodique
\fill[poporange!50] (-1.1,-0.6) rectangle (1.1,-0.35);
\node[poporange!80!popdark, font=\footnotesize\bfseries] at (0,-0.2) {ANODE ($-$)};
\node[poporange!80!popdark, font=\footnotesize, align=center] at (0,-1.95) {$\ce{Fe -> Fe^{2+} + 2 e-}$\\ le fer se dissout (piqûre)};
\draw[-{Stealth}, poporange!80!popdark, thick] (0,-1.5) -- (0,-0.75);
% Zone cathodique gauche
\fill[popgreen!45] (-3.2,-0.6) rectangle (-2.2,-0.35);
\node[popgreen!60!popdark, font=\footnotesize\bfseries] at (-2.7,-0.15) {CATHODE ($+$)};
% Zone cathodique droite
\fill[popgreen!45] (2.2,-0.6) rectangle (3.2,-0.35);
\node[popgreen!60!popdark, font=\footnotesize\bfseries] at (2.7,-0.15) {CATHODE ($+$)};
\node[popgreen!60!popdark, font=\footnotesize, align=center] at (4.6,1.1) {$\ce{O2 + 2 H2O + 4 e- -> 4 HO-}$\\ (bord de goutte bien aéré)};
\draw[-{Stealth}, popgreen!60!popdark] (3.6,0.75) -- (2.9,0.1);
% Flux d'électrons dans le métal
\draw[-{Stealth}, poppurple, very thick] (-0.9,-0.9) .. controls (-1.8,-1.25) .. (-2.6,-0.85);
\draw[-{Stealth}, poppurple, very thick] (0.9,-0.9) .. controls (1.8,-1.25) .. (2.6,-0.85);
\node[poppurple, font=\footnotesize] at (-1.8,-1.35) {$e^-$};
\node[poppurple, font=\footnotesize] at (1.8,-1.35) {$e^-$};
% Migration ionique dans la goutte
\draw[-{Stealth}, poppink!80!popdark, thick] (-0.4,0.15) .. controls (-1.4,0.6) .. (-2.4,0.15);
\node[poppink!80!popdark, font=\footnotesize] at (-1.5,0.75) {\ce{Fe^{2+}}};
\draw[-{Stealth}, popblue!80!popdark, thick] (2.4,0.35) .. controls (1.4,0.9) .. (0.4,0.4);
\node[popblue!80!popdark, font=\footnotesize] at (1.6,1.05) {\ce{HO-}};
% Rouille
\fill[popgold!70] (-0.7,1.15) circle (0.18);
\fill[popgold!70] (0.1,1.35) circle (0.14);
\fill[popgold!70] (0.6,1.1) circle (0.11);
\node[popgold!60!popdark, font=\footnotesize, align=center] at (-2.3,1.6) {Rouille \ce{Fe2O3.xH2O}\\ (anneau entre les deux zones)};
\draw[-{Stealth}, popgold!60!popdark] (-1.4,1.55) -- (-0.5,1.3);
% O2 de l'air
\draw[-{Stealth}, popdark] (4.2,2.1) -- (3.3,1.5);
\node[popdark, font=\footnotesize] at (4.5,2.25) {\ce{O2} de l'air};
\end{tikzpicture}
\end{center}
\legende{Micropile de corrosion sous une goutte d'eau salée sur du fer : le centre de la goutte (pauvre en \ce{O2}) devient anodique et se corrode ; les bords bien aérés deviennent cathodiques ; la rouille se dépose en anneau entre les deux zones.}
\end{popfigure}

\textbf{De la dissolution à la rouille : la cascade de réactions}\par
Les ions \ce{Fe^{2+}} formés à l'anode et les ions \ce{HO-} formés à la cathode migrent dans la goutte et se rencontrent : il se forme d'abord un précipité vert pâle d'hydroxyde de fer(II) :
\[ \ce{Fe^{2+} + 2 HO- -> Fe(OH)2 (s)} \]
Ce précipité est à son tour oxydé par le dioxygène dissous en hydroxyde de fer(III) rouille-brun :
\[ \ce{4 Fe(OH)2 + O2 + 2 H2O -> 4 Fe(OH)3 (s)} \]
Enfin, \ce{Fe(OH)3} se déshydrate partiellement et se transforme en \cle{rouille}, un oxyde de fer(III) hydraté de formule approchée \ce{Fe2O3.xH2O} :
\[ \ce{2 Fe(OH)3 -> Fe2O3.xH2O + (3 - x) H2O} \]
C'est pourquoi la rouille n'apparaît pas exactement là où le fer se dissout : elle se dépose \emph{entre} les zones anodiques et cathodiques, là où les ions se rencontrent.

\begin{attention}[La rouille ne protège pas le fer !]
Contrairement à la couche d'alumine \ce{Al2O3} sur l'aluminium, la rouille est \textbf{poreuse, friable et non adhérente} : elle se craquelle et se détache, laissant l'eau et le dioxygène atteindre le métal frais situé en dessous. La corrosion reprend alors \emph{en profondeur}, couche après couche, jusqu'à la perforation complète de la pièce. Ne jamais croire qu'une couche de rouille « stabilise » un ouvrage : c'est un cancer qui progresse. Seuls certains aciers spéciaux (aciers « Corten ») forment une rouille compacte et protectrice.
\end{attention}

\begin{exoresolu}[Bilan électronique d'une micropile]
\textbf{Énoncé.} Une goutte d'eau salée sur une plaque de fer constitue une micropile. (1) Écrire les demi-équations aux deux zones et l'équation-bilan de la réaction de corrosion. (2) Quelle masse de fer est consommée lorsque \SI{0,32}{g} de dioxygène a été réduit ? On donne $M(\ce{Fe}) = \SI{56}{g.mol^{-1}}$ et $M(\ce{O2}) = \SI{32}{g.mol^{-1}}$.
\tcblower
\textbf{Solution.} (1) Anode (oxydation du fer) : $\ce{Fe -> Fe^{2+} + 2 e-}$. Cathode (réduction du dioxygène) : $\ce{O2 + 2 H2O + 4 e- -> 4 HO-}$. Pour combiner, on multiplie la première par 2 afin d'égaliser les électrons :
\[ \ce{2 Fe + O2 + 2 H2O -> 2 Fe^{2+} + 4 HO-} \]
soit, en regroupant les ions : $\ce{2 Fe + O2 + 2 H2O -> 2 Fe(OH)2}$.
(2) Quantité de dioxygène réduit : $n(\ce{O2}) = \dfrac{0,32}{32} = \SI{0,010}{mol}$. D'après les coefficients stœchiométriques, $n(\ce{Fe}) = 2\,n(\ce{O2}) = \SI{0,020}{mol}$. Masse de fer consommée :
\[ m(\ce{Fe}) = 0,020 \times 56 = \SI{1,12}{g} \]
Vérification par les électrons : \SI{0,010}{mol} de \ce{O2} consomme \SI{0,040}{mol} d'électrons ; \SI{0,020}{mol} de Fe en libère $2 \times 0,020 = \SI{0,040}{mol}$. Le bilan électronique est cohérent : la corrosion est bien une pile en court-circuit.
\end{exoresolu}

\courssub{Facteurs favorisants la corrosion}

La vitesse de corrosion n'est pas une fatalité constante : elle dépend fortement de l'environnement et de l'état du métal.

\begin{itemize}
\item \textbf{Humidité} : elle fournit l'électrolyte. C'est le facteur n°1 : plus l'air est humide, plus la corrosion est rapide (climat équatorial du Sud-Cameroun).
\item \textbf{Les sels dissous} (\ce{NaCl} notamment) : ils augmentent la \emph{conductivité} de l'eau, donc la circulation ionique des micropiles. L'eau de mer et les routes salées en hiver sont des accélérateurs redoutables.
\item \textbf{L'acidité} : les ions \ce{H3O+} (pluies acides, atmosphères industrielles) offrent un oxydant supplémentaire qui se réduit en dégageant \ce{H2} : $\ce{2 H3O+ + 2 e- -> H2 + 2 H2O}$.
\item \textbf{Les impuretés et l'hétérogénéité du métal} : chaque inclusion, chaque soudure, chaque zone déformée crée un couple anode/cathode supplémentaire. Un acier pur et homogène se corrode moins vite qu'un acier brut.
\item \textbf{Les contraintes mécaniques} (pliures, rivetages, vibrations) : elles fragilisent les couches superficielles et créent des zones anodiques.
\item \textbf{Le contact entre deux métaux différents} : c'est le \cle{couple galvanique}, le plus redoutable de tous.
\end{itemize}

\begin{exemplebox}[Les pirogues de Limbe]
Les pêcheurs de Limbe et d'Idenau constatent que les coques et pirogues en tôle se percent en quelques années en mer, alors que les mêmes embarcations durent beaucoup plus longtemps sur le lac Tchad ou en eau douce. L'explication est électrochimique : l'eau de mer contient environ \SI{35}{g.L^{-1}} de sels dissous (surtout \ce{NaCl}), contre moins de \SI{1}{g.L^{-1}} pour l'eau douce. La conductivité de l'électrolyte est donc des dizaines de fois plus grande : le « pont salin » des micropiles fonctionne à plein régime, et la corrosion s'accélère d'autant. C'est pourquoi les coques des navires sont inspectées et repeintes si souvent.
\end{exemplebox}

\textbf{Le couple galvanique : quand deux métaux se touchent}\par
Si deux métaux de potentiels différents sont en contact électrique dans un électrolyte, ils constituent une véritable pile : le métal le \emph{plus réducteur} (potentiel le plus bas) devient l'anode et se corrode \emph{préférentiellement}, tandis que l'autre est protégé. C'est la règle du gamma appliquée malgré nous !

\begin{popfigure}
\begin{center}
\begin{tikzpicture}[scale=1.0, every node/.style={font=\small}]
% Eau
\fill[popblue!15] (-6,-2.2) rectangle (6,0.4);
\draw[popblue, thick] (-6,0.4) -- (6,0.4);
\node[popblue!70!popdark, font=\footnotesize] at (-4.6,0.15) {eau de mer (électrolyte)};
% Plaque de fer
\fill[popmuted!70] (-4.5,-1.8) rectangle (0.5,-0.2);
\node[popdark, font=\footnotesize\bfseries] at (-2,-1.0) {FER};
\node[poporange!80!popdark, font=\footnotesize\bfseries] at (-2,-1.5) {ANODE : se corrode};
% Plaque de cuivre
\fill[popgold!60] (0.5,-1.8) rectangle (4.5,-0.2);
\node[popdark, font=\footnotesize\bfseries] at (2.5,-1.0) {CUIVRE};
\node[popgreen!60!popdark, font=\footnotesize\bfseries] at (2.5,-1.5) {CATHODE : protégé};
% Rivet / contact
\fill[popdark] (0.5,-1.0) circle (0.12);
\node[popdark, font=\footnotesize] at (0.5,-2.05) {contact (rivet)};
% Flux électrons
\draw[-{Stealth}, poppurple, very thick] (-1.5,-0.45) -- (1.5,-0.45);
\node[poppurple, font=\footnotesize] at (0,-0.2) {$e^-$};
% Equations
\node[poporange!80!popdark, font=\footnotesize, align=center] at (-3.6,-2.9) {$\ce{Fe -> Fe^{2+} + 2 e-}$\\ ($E^\circ = \SI{-0,44}{V}$)};
\node[popgreen!60!popdark, font=\footnotesize, align=center] at (3.6,-2.9) {$\ce{O2 + 2 H2O + 4 e- -> 4 HO-}$\\ (le cuivre ne s'oxyde pas)};
% Rouille sur le fer
\fill[popgold!80] (-3.4,-0.55) circle (0.13);
\fill[popgold!80] (-2.8,-0.5) circle (0.10);
\fill[popgold!80] (-3.0,-0.42) circle (0.08);
\node[popgold!60!popdark, font=\footnotesize] at (-4.3,-0.05) {rouille};
\end{tikzpicture}
\end{center}
\legende{Couple galvanique fer/cuivre en milieu humide : le fer, plus réducteur ($E^\circ = \SI{-0,44}{V} < E^\circ(\ce{Cu^{2+}}/\ce{Cu}) = \SI{+0,34}{V}$), devient l'anode et se corrode à la place du cuivre. Règle pratique : dans un assemblage de deux métaux, c'est toujours le plus réducteur qui « sacrifie ».}
\end{popfigure}

\begin{attention}[Assembler deux métaux : un piège classique]
River ou boulonner directement du cuivre sur du fer (ou de l'acier inoxydable sur de l'acier ordinaire) en milieu humide revient à \emph{fabriquer une pile de corrosion} : le métal le plus réducteur se détruit de façon accélérée au voisinage du joint. Inversement — et c'est toute la ruse de la protection cathodique que nous étudierons dans la section suivante — on peut \emph{volontairement} fixer un métal très réducteur (zinc, magnésium) sur une coque en acier pour qu'il se corrode à sa place : c'est l'\emph{anode sacrificielle}.
\end{attention}

\courssub{Le coût de la corrosion : un enjeu économique majeur}

La corrosion n'est pas qu'un problème de chimiste : c'est un problème d'ingénieur et d'économiste. On estime que la corrosion détruit chaque année l'équivalent d'\textbf{environ un quart de la production mondiale d'acier}, et que son coût global (remplacements, réparations, arrêts de production, accidents) atteint plusieurs pourcents du produit intérieur brut des pays industrialisés.

Pour le Cameroun, l'enjeu est concret et quotidien :
\begin{itemize}
\item les \textbf{ponts} métalliques (pont sur le Wouri à Douala) exposés à un air chaud, humide et chargé de sel marin ;
\item les installations du \textbf{port de Douala} et du \textbf{port en eau profonde de Kribi}, où quais, grues et coques subissent l'agression permanente de l'eau de mer ;
\item l'\textbf{oléoduc Tchad--Cameroun}, plus de \SI{1000}{km} de canalisation d'acier enfouie dans des sols humides, dont la surveillance anticorrosion conditionne la sécurité de l'exportation du pétrole ;
\item les \textbf{châteaux d'eau}, pylônes électriques et charpentes des marchés, dont l'entretien pèse sur les budgets communaux.
\end{itemize}

\begin{aretenir}[L'essentiel sur le mécanisme de la corrosion]
\begin{itemize}
\item La corrosion est une \textbf{oxydoréduction} entre un métal (réducteur) et son environnement (oxydant : \ce{O2}, parfois \ce{H3O+}).
\item Le fer ne se corrode qu'en présence \textbf{simultanée} d'eau et de dioxygène.
\item Le mécanisme est celui de \textbf{micropiles} : anode $\ce{Fe -> Fe^{2+} + 2 e-}$, cathode $\ce{O2 + 2 H2O + 4 e- -> 4 HO-}$, puis formation de \ce{Fe(OH)2}, oxydation en \ce{Fe(OH)3} et déshydratation en rouille \ce{Fe2O3.xH2O}.
\item La rouille est \textbf{poreuse} : elle ne protège pas le fer, contrairement à l'alumine sur l'aluminium.
\item Sels, acidité, impuretés, contraintes et contacts entre métaux différents (couple galvanique) \textbf{accélèrent} la corrosion.
\end{itemize}
\end{aretenir}

Puisque la corrosion est une pile involontaire, la combattre revient à \emph{casser le circuit} de cette pile ou à \emph{détourner} son fonctionnement : isoler le métal de l'eau et de l'air, ou lui imposer de devenir cathode. C'est précisément l'objet de la section suivante, consacrée aux méthodes de protection : revêtements, galvanisation et protection cathodique.

\coursec{Protection contre la corrosion}

Dans la section précédente, nous avons vu que la corrosion du fer résulte de la formation de \cle{micropiles} à la surface du métal : le fer s'oxyde en \ce{Fe^2+} pendant que le dioxygène dissous se réduit. Nous savons donc maintenant \emph{pourquoi} le fer rouille. La question pratique qui se pose alors à l'ingénieur, au charpentier métallique de Douala comme au constructeur naval, est simple : \textbf{comment empêcher cette micropile de fonctionner ?} C'est toute la science de la protection contre la corrosion, et tu vas voir qu'elle découle directement de la classification électrochimique étudiée au début de cette leçon.

\courssub{Principe général de la protection}

Puisque la corrosion est une réaction d'oxydoréduction qui nécessite le contact simultané du métal avec le dioxygène et l'eau, deux grandes stratégies s'offrent à nous :

\begin{itemize}
\item \textbf{Stratégie 1 — la barrière physique} : isoler le métal de son environnement (\ce{O2} et \ce{H2O}) par un revêtement imperméable. Si l'air et l'eau ne touchent plus le fer, la micropile ne peut plus fonctionner.
\item \textbf{Stratégie 2 — la protection électrochimique} : imposer au métal à protéger le rôle de \cle{cathode} (pôle où se produit une réduction). Or un métal ne se corrode que s'il s'oxyde, c'est-à-dire s'il joue le rôle d'anode. En forçant le fer à être une cathode, on bloque sa dissolution : c'est la \cle{protection cathodique}.
\end{itemize}

\begin{aretenir}[Les deux grandes stratégies]
Protéger un métal contre la corrosion, c'est soit \textbf{empêcher le contact} métal/environnement (revêtements), soit \textbf{inverser la polarité} du métal pour qu'il devienne cathode (protection cathodique par anode sacrificielle ou par courant imposé).
\end{aretenir}

\courssub{Protections par revêtements non métalliques}

L'idée la plus ancienne et la plus simple : recouvrir le métal d'une couche étanche.

\begin{itemize}
\item \textbf{La peinture} : les châssis métalliques, les portails, les poteaux électriques sont peints. La pellicule de peinture empêche \ce{O2} et \ce{H2O} d'atteindre le fer. Dès que la peinture s'écaille, la corrosion reprend : d'où la nécessité d'un entretien régulier.
\item \textbf{Le vernis} : utilisé sur les boîtes de conserve (intérieur verni) et les objets décoratifs.
\item \textbf{Le graissage} : une couche de graisse ou d'huile protège les pièces mécaniques (outils, chaînes, câbles). Protection temporaire mais facile à renouveler.
\item \textbf{La plastification} : le métal est enrobé d'une couche de plastique (grillages de clôture verts, fils électriques, étagères de réfrigérateur).
\end{itemize}

\begin{savaistu}[La Tour Eiffel, un combat permanent contre la rouille]
La Tour Eiffel, 7300 tonnes de fer puddlé, est \textbf{repeinte intégralement tous les 7 ans} environ : il faut près de \SI{60}{t} de peinture et des mois de travail à des peintres alpinistes. Sans ce soin, la dame de fer se serait littéralement dissoute dans l'air parisien ! La peinture est ici une barrière physique : elle coûte cher, mais bien moins que la perte du monument.
\end{savaistu}

\begin{attention}[Limite des revêtements non métalliques]
Un revêtement de peinture, de vernis ou de graisse n'est qu'une \textbf{barrière passive} : dès qu'il est rayé, écaillé ou poreux, la corrosion reprend \emph{localement}, souvent de façon plus intense (piqûres). Ces protections exigent donc une surveillance et un entretien réguliers.
\end{attention}

\courssub{Revêtements métalliques : étamage, chromage, nickelage}

On peut aussi recouvrir le fer d'un \textbf{autre métal} déposé en couche mince, le plus souvent par \cle{électrolyse} :

\begin{itemize}
\item \textbf{L'étamage} : dépôt d'étain \ce{Sn} sur la tôle d'acier. C'est le principe du \emph{fer-blanc} des \textbf{boîtes de conserve} (haricots, tomates, sardines). L'étain est peu réducteur et résiste bien à l'air humide.
\item \textbf{Le chromage} : dépôt de chrome \ce{Cr}, dur et brillant (robinetterie, pare-chocs, guidons de moto).
\item \textbf{Le nickelage} : dépôt de nickel \ce{Ni} (pièces de monnaie, ustensiles).
\end{itemize}

\begin{attention}[Le piège de l'étamage : la rayure assassine]
L'étain est placé \emph{en dessous} du fer dans la classification électrochimique : il est \textbf{moins réducteur} que le fer. Tant que le revêtement est intact, tout va bien. Mais si la boîte de conserve est \textbf{rayée}, le fer mis à nu devient l'\textbf{anode} d'une micropile \ce{Fe}/\ce{Sn} et s'oxyde \emph{plus vite} que s'il n'avait pas été étamé :
\[ \ce{Fe -> Fe^2+ + 2e-} \qquad \text{(le fer se sacrifie, l'étain est protégé !)} \]
C'est l'\textbf{inverse exact de la galvanisation}. Moralité : une boîte de conserve cabossée ou rayée doit être consommée rapidement.
\end{attention}

\courssub{La galvanisation : le zinc, gardien du fer}

\begin{definition}[Galvanisation]
La \cle{galvanisation} est le recouvrement du fer (ou de l'acier) par une couche de \textbf{zinc} \ce{Zn}, obtenue par trempage dans du zinc fondu (vers \SI{450}{\celsius}) ou par électrolyse. Le zinc protège le fer \textbf{deux fois} : mécaniquement (barrière physique) et électrochimiquement (protection cathodique).
\end{definition}

Pourquoi le zinc est-il si spécial ? Regarde la classification électrochimique : le couple \ce{Zn^2+}/\ce{Zn} est placé \emph{au-dessus} du couple \ce{Fe^2+}/\ce{Fe}. Le zinc est donc \textbf{plus réducteur} que le fer.

\begin{itemize}
\item \textbf{Tant que la couche de zinc est intacte}, elle joue le rôle de barrière : ni \ce{O2} ni \ce{H2O} n'atteignent le fer. De plus, le zinc se recouvre d'une couche protectrice d'oxyde et de carbonate de zinc (la \emph{patine} grise mate) qui le rend lui-même très résistant.
\item \textbf{Si la tôle est rayée} et le fer mis à nu, une micropile \ce{Zn}/\ce{Fe} se forme en présence d'eau humide. D'après la règle du gamma, c'est le zinc, plus réducteur, qui s'oxyde :
\[ \ce{Zn -> Zn^2+ + 2e-} \qquad \text{(anode : le zinc se sacrifie)} \]
tandis qu'à la surface du fer (cathode) le dioxygène se réduit :
\[ \ce{O2 + 2H2O + 4e- -> 4OH-} \]
Le fer, devenu cathode, \textbf{ne s'oxyde pas} : il est protégé même à découvert !
\end{itemize}

\begin{popfigure}
\begin{center}
\begin{tikzpicture}[scale=1.0, every node/.style={font=\small}]
% Couche de zinc (parties gauche et droite)
\fill[poppurple!45] (0,1.2) rectangle (4.2,2.0);
\fill[poppurple!45] (5.8,1.2) rectangle (10,2.0);
% Base fer
\fill[popblue!50] (0,0) rectangle (10,1.2);
% Rayure : fer apparent jusqu'en haut
\fill[popblue!50] (4.2,0) rectangle (5.8,2.0);
% Contours
\draw[thick, popdark] (0,0) rectangle (10,2.0);
\draw[thick, popdark] (0,1.2) -- (4.2,1.2) (5.8,1.2) -- (10,1.2);
\draw[thick, popdark] (4.2,1.2) -- (4.2,2.0) (5.8,1.2) -- (5.8,2.0);
% Labels
\node at (2.2,0.6) {\textbf{Fer} \ce{Fe}};
\node at (8,0.6) {\textbf{Fer} \ce{Fe}};
\node at (2.2,1.6) {\textbf{Zinc} \ce{Zn}};
\node at (8,1.6) {\textbf{Zinc} \ce{Zn}};
\node[popdark] at (5,1.6) {\footnotesize rayure};
% Goutte d'eau au-dessus de la rayure
\draw[popblue, thick, fill=popblue!20] (5,2.6) circle (0.28);
\node[popblue, font=\footnotesize] at (5,3.15) {eau + \ce{O2}};
% Flèches électrons : du zinc vers le fer
\draw[-{Stealth}, very thick, poporange] (4.0,1.6) .. controls (4.4,2.5) and (4.7,2.6) .. (4.85,1.75);
\draw[-{Stealth}, very thick, poporange] (6.0,1.6) .. controls (5.6,2.5) and (5.3,2.6) .. (5.15,1.75);
\node[poporange, font=\footnotesize] at (3.4,2.55) {$e^-$};
\node[poporange, font=\footnotesize] at (6.6,2.55) {$e^-$};
% Rôles
\node[popgreen!60!black, font=\footnotesize, align=center] at (2.2,2.45) {anode : \ce{Zn -> Zn^2+ + 2e-}};
\node[popink, font=\footnotesize, align=center] at (8.2,2.45) {cathode (fer protégé)};
\end{tikzpicture}
\end{center}
\legende{Coupe d'une tôle galvanisée rayée : le zinc, plus réducteur que le fer, s'oxyde à sa place. Les électrons circulent du zinc vers le fer, qui devient cathode et reste intact.}
\end{popfigure}

\begin{exemplebox}[Les tôles ondulées de nos toitures]
Les \textbf{tôles ondulées galvanisées} qui couvrent la plupart de nos maisons, de Yaoundé à Maroua, résistent \textbf{des décennies} sous la pluie et le soleil, alors qu'une tôle de fer nue rouillerait en quelques saisons des pluies. C'est la double protection du zinc : barrière physique, puis protection cathodique aux endroits rayés ou percés (trous de clous !). C'est pourquoi on choisit des clous \emph{galvanisés} eux aussi : un clou de fer ordinaire s'oxyderait rapidement au contact de la tôle.
\end{exemplebox}

\courssub{Protection cathodique par anode sacrificielle}

La galvanisation n'est pas toujours possible : comment galvaniser la coque d'un navire de 100 mètres, ou un pipeline enterré ? On utilise alors le même principe électrochimique, mais de manière localisée.

\begin{definition}[Anode sacrificielle]
Une \cle{anode sacrificielle} est un bloc de métal \textbf{plus réducteur} que le métal à protéger (zinc, magnésium, aluminium), fixé métalliquement sur la structure. Il s'oxyde \emph{à la place} du fer — il se « sacrifie » — et doit être remplacé périodiquement.
\end{definition}

Sur la coque en acier d'un navire, on boulonne des blocs de magnésium ou de zinc. La coque et le bloc forment une pile dans l'eau de mer (excellent électrolyte) : le magnésium est l'anode, la coque devient la cathode.

\[ \ce{Mg -> Mg^2+ + 2e-} \qquad \text{(anode sacrificielle, consommée lentement)} \]
\[ \ce{O2 + 2H2O + 4e- -> 4OH-} \qquad \text{(à la surface de la coque, protégée)} \]

\begin{popfigure}
\begin{center}
\begin{tikzpicture}[scale=1.0, every node/.style={font=\small}]
% Eau de mer
\fill[popblue!12] (0,0) rectangle (11,3.2);
\node[popblue!70, font=\footnotesize] at (9.6,0.4) {eau de mer};
% Coque du navire
\fill[popblue!55] (1,3.2) -- (1,4.6) -- (9,4.6) -- (10,3.6) -- (10,3.2) -- cycle;
\draw[thick, popdark] (1,3.2) -- (1,4.6) -- (9,4.6) -- (10,3.6) -- (10,3.2);
\node[white, font=\small\bfseries] at (5,4.15) {Coque en acier (cathode, protégée)};
% Anode sacrificielle
\fill[popgold!70] (4.3,2.0) rectangle (5.7,2.9);
\draw[thick, popdark] (4.3,2.0) rectangle (5.7,2.9);
\node[popdark, font=\footnotesize\bfseries, align=center] at (5,2.45) {Bloc de Mg\\ (anode)};
% Liaison métallique
\draw[line width=2.5pt, popmuted] (5,3.2) -- (5,2.9);
% Flèches électrons Mg -> coque
\draw[-{Stealth}, very thick, poporange] (4.6,2.9) -- (4.6,3.55);
\draw[-{Stealth}, very thick, poporange] (5.4,2.9) -- (5.4,3.55);
\node[poporange, font=\footnotesize] at (4.15,3.25) {$e^-$};
\node[poporange, font=\footnotesize] at (5.85,3.25) {$e^-$};
% Ions Mg2+ qui partent dans l'eau
\draw[-{Stealth}, thick, poppurple] (4.3,2.3) -- (3.2,1.5);
\draw[-{Stealth}, thick, poppurple] (5.7,2.3) -- (6.8,1.5);
\node[poppurple, font=\footnotesize] at (2.9,1.25) {\ce{Mg^2+}};
\node[poppurple, font=\footnotesize] at (7.1,1.25) {\ce{Mg^2+}};
\end{tikzpicture}
\end{center}
\legende{Anode sacrificielle en magnésium boulonnée sur une coque de navire : les électrons circulent du magnésium vers l'acier, qui devient cathode. Le bloc de Mg se dissout lentement et se remplace.}
\end{popfigure}

\begin{methode}[Choisir le métal d'une anode sacrificielle]
\begin{enumerate}
\item Repérer le métal à protéger (en général le fer, couple \ce{Fe^2+}/\ce{Fe}).
\item Consulter la classification électrochimique : le métal protecteur doit être placé \textbf{au-dessus} du fer, donc être \textbf{plus réducteur}.
\item Candidats usuels : \ce{Mg} (très efficace, eau de mer), \ce{Zn} (coques, pipelines), \ce{Al} (réservoirs, ballons d'eau chaude).
\item Vérifier la règle du gamma : la réaction entre le métal protecteur et l'oxydant de l'environnement (\ce{O2} dissous) doit être spontanée.
\item Prévoir le \textbf{remplacement périodique} du bloc : une anode sacrificielle est un consommable !
\end{enumerate}
\end{methode}

\begin{exoresolu}[Choisir entre cuivre et magnésium pour protéger une canalisation]
Une canalisation enterrée en fer doit être protégée par une anode sacrificielle. Un technicien hésite entre un bloc de cuivre et un bloc de magnésium. Classification (extraite) : \ce{Mg^2+}/\ce{Mg} au-dessus de \ce{Fe^2+}/\ce{Fe}, lui-même au-dessus de \ce{Cu^2+}/\ce{Cu}. Lequel choisir ? Justifier par les équations.
\tcblower
\textbf{Analyse.} L'anode sacrificielle doit être \emph{plus réductrice} que le fer pour s'oxyder à sa place (règle du gamma).
\begin{itemize}
\item \textbf{Cuivre} : placé \emph{en dessous} du fer, il est moins réducteur. La micropile \ce{Fe}/\ce{Cu} aurait le fer pour anode : $\ce{Fe -> Fe^2+ + 2e-}$. Le cuivre \textbf{accélérerait} la corrosion de la canalisation. Choix catastrophique !
\item \textbf{Magnésium} : placé \emph{au-dessus} du fer, il est plus réducteur. C'est lui qui s'oxyde :
\[ \ce{Mg -> Mg^2+ + 2e-} \]
et le fer, devenu cathode, est protégé : $\ce{O2 + 2H2O + 4e- -> 4OH-}$.
\end{itemize}
\textbf{Conclusion.} On choisit le \textbf{magnésium}. Bilan global de la protection : $\ce{2Mg + O2 + 2H2O -> 2Mg^2+ + 4OH-}$, réaction naturelle conforme à la règle du gamma. Le bloc de Mg sera remplacé quand il sera consommé.
\end{exoresolu}

\courssub{Protection cathodique par courant imposé}

Pour de très grandes structures (pipelines, oléoducs de plusieurs centaines de kilomètres, tanks de stockage), multiplier les anodes sacrificielles devient coûteux. On utilise alors une source extérieure : un \textbf{générateur de courant continu}.

\begin{propriete}[Principe du courant imposé]
La structure métallique à protéger est reliée au \textbf{pôle négatif} ($-$) du générateur : elle reçoit des électrons et devient ainsi une \textbf{cathode} forcée. Le pôle positif ($+$) est relié à une électrode auxiliaire enterrée (anode de mise à la terre, en ferraille ou en graphite), qui s'oxyde à la place de la structure.
\end{propriete}

Tant que le courant circule, le fer ne peut pas s'oxyder : les électrons qu'il devrait fournir lui sont imposés par le générateur. Cette méthode protège par exemple les pipelines qui transportent le pétrole de la SONARA ou les oléoducs traversant des zones humides.

\begin{attention}[Ne pas inverser les bornes !]
Si la structure est branchée par erreur sur le pôle $+$ du générateur, elle devient \textbf{anode} et se corrode de façon \emph{fulgurante}, bien plus vite qu'à l'air libre. La protection cathodique par courant imposé exige un contrôle rigoureux des polarités et une alimentation électrique permanente.
\end{attention}

\courssub{Choisir la méthode adaptée au contexte}

Il n'existe pas de protection universelle : l'ingénieur choisit selon le \textbf{coût}, la \textbf{durée} souhaitée, l'\textbf{accessibilité} de l'ouvrage et la nature du \textbf{milieu} (marin, terrestre, atmosphérique).

\begin{center}
\begin{tabularx}{\linewidth}{|l|X|X|}
\hline
\rowcolor{popblueL} \textbf{Méthode} & \textbf{Avantages} & \textbf{Limites / contexte d'emploi} \\
\hline
Peinture, vernis, graisse & Peu coûteuse, simple & Entretien fréquent ; objets accessibles (portails, ponts, machines) \\
\hline
Étamage, chromage, nickelage & Aspect décoratif, hygiène (conserve) & Dangereux si rayé ; objets de petite taille \\
\hline
Galvanisation & Double protection, durable (décennies) & Pièces de taille modérée trempées en usine (tôles, poteaux, grillages) \\
\hline
Anode sacrificielle & Autonome, sans énergie & Blocs à remplacer ; milieu conducteur (mer, sol humide) : coques, ballons d'eau chaude \\
\hline
Courant imposé & Protège de très grandes structures & Nécessite une alimentation et une surveillance ; pipelines, oléoducs \\
\hline
\end{tabularx}
\end{center}

\courssub{Bilan de la leçon}

Cette leçon t'a mené sur un chemin complet, du laboratoire au chantier :

\begin{itemize}
\item La \textbf{classification électrochimique} range les couples oxydant/réducteur et permet, grâce à la \textbf{règle du gamma}, de prévoir toute réaction d'oxydoréduction naturelle.
\item La \textbf{pile Daniell} et ses sœurs transforment une réaction redox spontanée en courant électrique : le métal le plus réducteur constitue le pôle $-$ (anode, oxydation), l'autre le pôle $+$ (cathode, réduction), et la f.é.m. se mesure au voltmètre.
\item La \textbf{corrosion} n'est autre qu'une pile involontaire : les micropiles à la surface du fer le font s'oxyder en présence d'eau et de dioxygène.
\item La \textbf{protection} inverse la logique de la pile : barrière physique contre \ce{O2} et \ce{H2O}, ou protection cathodique qui force le fer à rester cathode, grâce à un métal plus réducteur (galvanisation, anode sacrificielle) ou à un générateur (courant imposé).
\end{itemize}

\begin{aretenir}[L'idée directrice]
Toute cette leçon tient en une phrase : \textbf{dans un couple de métaux en contact, c'est toujours le plus réducteur qui s'oxyde}. La pile exploite cette loi pour produire de l'électricité ; la corrosion en est la conséquence subie ; la protection cathodique la détourne au profit de l'ouvrage. Qui maîtrise la classification électrochimique maîtrise les trois.
\end{aretenir}

\begin{exercice}[Pour vérifier ta compréhension]
\begin{enumerate}
\item Pourquoi un clou en cuivre ne doit-il jamais être utilisé pour fixer une tôle galvanisée ?
\item Un ballon d'eau chaude en acier contient une « résistance » entourée d'un fourreau de magnésium. Quel est le rôle de ce fourreau ? Écris les équations mises en jeu.
\item Classe ces métaux du plus au moins adapté pour protéger le fer par anode sacrificielle : étain, zinc, plomb, magnésium.
\end{enumerate}
\end{exercice}

\begin{corrige}[Exercice : protection des métaux]
\begin{enumerate}
\item Le cuivre est \emph{moins} réducteur que le fer : au contact de la tôle, la micropile \ce{Fe}/\ce{Cu} fait du fer l'anode ($\ce{Fe -> Fe^2+ + 2e-}$). Le clou de cuivre accélère la corrosion de la tôle au point de contact.
\item Le fourreau de magnésium est une \textbf{anode sacrificielle} : plus réducteur que le fer, il s'oxyde à sa place ($\ce{Mg -> Mg^2+ + 2e-}$) pendant que l'eau et le dioxygène se réduisent à la paroi du ballon, devenue cathode ($\ce{O2 + 2H2O + 4e- -> 4OH-}$). Il doit être changé périodiquement.
\item Du plus adapté au moins adapté : \textbf{magnésium} $>$ \textbf{zinc} (tous deux plus réducteurs que le fer) ; le plomb et l'étain, moins réducteurs que le fer, sont inutilisables (ils accéléreraient la corrosion).
\end{enumerate}
\end{corrige}

\newpage

\coursec{Travaux pratiques}

\courssub{Objectifs du TP}

À l'issue de cette séance de travaux pratiques, tu seras capable de :

\begin{itemize}
  \item construire une \cle{pile Daniell} à partir de deux demi-piles $\ce{Zn^2+}/\ce{Zn}$ et $\ce{Cu^2+}/\ce{Cu}$ ;
  \item déterminer expérimentalement la \cle{polarité} de la pile à l'aide d'un voltmètre ;
  \item mesurer la \cle{force électromotrice} (f.é.m.) de la pile et la comparer à la valeur théorique $E = \SI{1,10}{V}$ ;
  \item vérifier que la pile débite un courant en alimentant une DEL ou un petit moteur ;
  \item écrire les demi-équations aux électrodes, l'équation-bilan de fonctionnement et la représentation conventionnelle de la pile ;
  \item justifier la polarité observée à l'aide de la \cle{règle du gamma} et de la classification électrochimique.
\end{itemize}

\courssub{Matériel et produits}

\begin{center}
\rowcolors{2}{popblueL}{white}
\begin{tabularx}{\linewidth}{|l|X|}
\hline
\rowcolor{popblue!40} \textbf{Matériel} & \textbf{Produits} \\
\hline
2 béchers de \SI{100}{mL} & Solution de sulfate de zinc \ce{ZnSO4} à \SI{0,1}{mol.L^{-1}} (ou \SI{1}{mol.L^{-1}}) \\
Lame de zinc et lame de cuivre & Solution de sulfate de cuivre \ce{CuSO4} à \SI{0,1}{mol.L^{-1}} (ou \SI{1}{mol.L^{-1}}) \\
Toile émeri (ou papier de verre fin) & Solution de nitrate de potassium \ce{KNO3} (ou \ce{KCl}) \\
Tube en U + agar-agar (pont salin) & Eau distillée \\
Voltmètre (multimètre) + 2 fils à pinces & \\
DEL ou petit moteur électrique & \\
Chronomètre, gants, lunettes & \\
\hline
\end{tabularx}
\end{center}

\begin{attention}[Solutions de rechange avec du matériel local]
Si ton laboratoire ne dispose pas de tube en U, fabrique un pont salin avec une \textbf{bande de papier filtre} (ou un coton torsadé) imbibée de solution concentrée de \ce{KNO3} ; à défaut de nitrate de potassium, une solution concentrée de \textbf{sel de cuisine} (\ce{NaCl}) convient. Les lames de cuivre peuvent provenir d'un vieux câble électrique dénudé, et le zinc d'une enveloppe de pile usagée de type « saline » (à nettoyer soigneusement). Un multimètre de dépannage automobile remplace avantageusement le voltmètre de laboratoire.
\end{attention}

\courssub{Sécurité}

\begin{attention}[Consignes de sécurité]
\begin{itemize}
  \item Le sulfate de cuivre \ce{CuSO4} est \textbf{nocif} et \textbf{irritant} (pictogramme « point d'exclamation » : danger modéré) et \textbf{dangereux pour l'environnement aquatique} (pictogramme « poisson et arbre morts ») : porte des \textbf{gants} et des \textbf{lunettes de protection}, ne jette jamais les solutions dans l'évier ou dans la nature ; verse-les dans le bidon de récupération prévu par le professeur.
  \item Ne porte jamais les solutions à la bouche ; lave-toi les mains en fin de séance.
  \item En cas de contact avec la peau ou les yeux, rince abondamment à l'eau claire pendant plusieurs minutes et préviens immédiatement l'enseignant.
  \item La toile émeri peut couper : décape les lames en les posant à plat sur la paillasse, jamais dans le creux de la main.
  \item Ne branche la DEL qu'après avoir identifié les pôles : une DEL branchée à l'envers ne s'allume pas et peut être détériorée si la tension est trop élevée.
\end{itemize}
\end{attention}

\courssub{Schéma du dispositif}

\begin{popfigure}
\begin{center}
\begin{tikzpicture}[scale=1.0, every node/.style={font=\small}]
  % Bécher gauche (Zn)
  \draw[thick, popdark] (-3.2,0) -- (-3.2,2.6) (-1.2,0) -- (-1.2,2.6) (-3.2,0) -- (-1.2,0);
  \fill[popblue!25] (-3.15,0.05) rectangle (-1.25,1.5);
  \draw[popblue!60] (-3.15,1.5) -- (-1.25,1.5);
  \node[popblue!80!black] at (-2.2,0.7) {$\ce{ZnSO4}$};
  % Lame Zn
  \fill[popmuted!60] (-2.45,0.4) rectangle (-2.05,3.4);
  \node[above, popdark] at (-2.25,3.4) {Lame de Zn};
  % Bécher droit (Cu)
  \draw[thick, popdark] (1.2,0) -- (1.2,2.6) (3.2,0) -- (3.2,2.6) (1.2,0) -- (3.2,0);
  \fill[popblue!50] (1.25,0.05) rectangle (3.15,1.5);
  \draw[popblue!80] (1.25,1.5) -- (3.15,1.5);
  \node[white] at (2.2,0.7) {$\ce{CuSO4}$};
  % Lame Cu
  \fill[poporange!70] (2.05,0.4) rectangle (2.45,3.4);
  \node[above, popdark] at (2.25,3.4) {Lame de Cu};
  % Pont salin
  \draw[line width=4pt, popgreen!70] (-1.9,1.3) .. controls (-1.4,3.6) and (1.4,3.6) .. (1.9,1.3);
  \node[popgreen!60!black] at (0,3.6) {Pont salin (\ce{KNO3})};
  % Fils et voltmètre
  \draw[thick, popdark] (-2.25,3.4) -- (-2.25,4.6) -- (-0.7,4.6);
  \draw[thick, popdark] (2.25,3.4) -- (2.25,4.6) -- (0.7,4.6);
  \draw[thick, poppurple, fill=poppurpleL] (0,4.6) circle (0.7);
  \node[poppurple!80!black] at (0,4.6) {\textbf{V}};
  \node[poppurple!80!black, below] at (0,3.85) {voltmètre};
  % Polarités
  \node[popdark, font=\large\bfseries] at (-2.9,4.6) {$\ominus$};
  \node[popdark, font=\large\bfseries] at (2.9,4.6) {$\oplus$};
  % Sens du courant
  \draw[-{Stealth[length=3mm]}, thick, poppink] (1.6,5.5) -- (-1.6,5.5);
  \node[poppink!80!black, above] at (0,5.5) {sens du courant $I$};
  \draw[-{Stealth[length=3mm]}, thick, popgold!80!black] (-1.6,5.9) -- (1.6,5.9);
  \node[popgold!70!black, above] at (0,5.9) {sens des électrons $e^-$};
\end{tikzpicture}
\end{center}
\legende{Schéma fonctionnel de la pile Daniell : la lame de zinc constitue le pôle négatif (anode, siège de l'oxydation), la lame de cuivre le pôle positif (cathode, siège de la réduction). Le pont salin assure la circulation des ions et ferme le circuit.}
\end{popfigure}

\courssub{Protocole}

\begin{methode}[Réalisation de la pile Daniell]
\begin{enumerate}
  \item \textbf{Décaper} les lames de zinc et de cuivre à la toile émeri jusqu'à ce qu'elles soient brillantes, puis les rincer à l'eau distillée et les sécher.
  \item Verser environ \SI{50}{mL} de solution de \ce{ZnSO4} dans le premier bécher et \SI{50}{mL} de solution de \ce{CuSO4} dans le second.
  \item Plonger la lame de zinc dans la solution de \ce{ZnSO4} et la lame de cuivre dans la solution de \ce{CuSO4} : tu obtiens deux \cle{demi-piles}.
  \item Relier les deux béchers par le \cle{pont salin} (tube en U rempli de \ce{KNO3} gélifié, ou bande de papier filtre imbibée), en veillant à ce que ses extrémités trempent dans les deux solutions.
  \item Brancher le voltmètre (calibre \SI{2}{V} continu) : borne \textbf{COM} sur la lame de zinc, borne \textbf{V} sur la lame de cuivre. Relever la tension affichée et son signe.
  \item \textbf{Inverser les branchements} et relever à nouveau : vérifier que l'affichage devient négatif.
  \item Débrancher le voltmètre et connecter la DEL (grande patte = anode, côté cuivre) ou le petit moteur aux bornes de la pile. Observer.
  \item Laisser la pile débiter pendant \SI{15}{min} à \SI{20}{min}, puis observer attentivement la lame de cuivre et la couleur de la solution de \ce{CuSO4}.
  \item Relever une dernière fois la tension en fin de séance.
\end{enumerate}
\end{methode}

\courssub{Observations et résultats attendus}

\begin{center}
\rowcolors{2}{popgreenL}{white}
\begin{tabularx}{\linewidth}{|X|l|}
\hline
\rowcolor{popgreen!50} \textbf{Étape} & \textbf{Observation attendue} \\
\hline
Voltmètre : COM sur Zn, V sur Cu & $U \approx +\SI{1,05}{V}$ (affichage positif) \\
Voltmètre : branchements inversés & $U \approx -\SI{1,05}{V}$ (affichage négatif) \\
Branchement de la DEL & La DEL s'allume (faiblement) ; le moteur tourne \\
Après 15 à \SI{20}{min} de fonctionnement & Dépôt \textbf{rougeâtre} de cuivre sur la lame de cuivre ; la solution bleue de \ce{CuSO4} se \textbf{décolore} progressivement ; la lame de zinc s'amincit \\
Tension en fin de séance & Légèrement inférieure à la valeur initiale \\
\hline
\end{tabularx}
\end{center}

La tension mesurée est généralement un peu inférieure à la valeur théorique $E = \SI{1,10}{V}$ (valable pour des solutions à \SI{1}{mol.L^{-1}} à \SI{25}{\celsius}). Les écarts s'expliquent par : des concentrations différentes de \SI{1}{mol.L^{-1}}, des lames mal décapées (couche d'oxyde), la résistance du pont salin et les potentiels de jonction.

\courssub{Exploitation}

\begin{exoresolu}[Exploitation des résultats]
\textbf{Énoncé.} À partir de tes observations : (1) identifier les pôles de la pile ; (2) écrire les demi-équations aux électrodes ; (3) écrire l'équation-bilan de fonctionnement ; (4) donner la représentation conventionnelle de la pile ; (5) justifier la polarité par la règle du gamma ($E^\circ(\ce{Cu^2+}/\ce{Cu}) = \SI{+0,34}{V}$ ; $E^\circ(\ce{Zn^2+}/\ce{Zn}) = \SI{-0,76}{V}$).
\tcblower
\textbf{Solution détaillée.}

\textbf{(1) Polarité.} L'affichage du voltmètre est positif lorsque la borne COM est reliée au zinc : la lame de \textbf{zinc} est donc le \cle{pôle négatif} ($\ominus$) et la lame de \textbf{cuivre} le \cle{pôle positif} ($\oplus$).

\textbf{(2) Demi-équations.} Le zinc s'amincit : il s'\textbf{oxyde} (anode, pôle $-$) :
\[ \ce{Zn -> Zn^2+ + 2e-} \]
Un dépôt rougeâtre de cuivre apparaît et la solution bleue se décolore : les ions \ce{Cu^2+} sont \textbf{réduits} (cathode, pôle $+$) :
\[ \ce{Cu^2+ + 2e- -> Cu} \]

\textbf{(3) Équation-bilan de fonctionnement :}
\[ \ce{Zn + Cu^2+ -> Zn^2+ + Cu} \]

\textbf{(4) Représentation conventionnelle :}
\[ \ominus \; \ce{Zn} \mid \ce{Zn^2+} \parallel \ce{Cu^2+} \mid \ce{Cu} \; \oplus \]

\textbf{(5) Règle du gamma.} On place les deux couples sur un axe vertical de potentiels croissants :
\begin{itemize}
  \item $\ce{Cu^2+}/\ce{Cu}$ : $E^\circ = \SI{+0,34}{V}$ (couple le plus haut) ;
  \item $\ce{Zn^2+}/\ce{Zn}$ : $E^\circ = \SI{-0,76}{V}$ (couple le plus bas).
\end{itemize}
La règle du gamma prévoit que l'oxydant le plus fort (\ce{Cu^2+}) réagit avec le réducteur le plus fort (\ce{Zn}) : la réaction $\ce{Zn + Cu^2+ -> Zn^2+ + Cu}$ est bien \textbf{naturelle}, ce qui confirme que le zinc est le pôle négatif. La f.é.m. théorique vaut :
\[ E = E^\circ(\ce{Cu^2+}/\ce{Cu}) - E^\circ(\ce{Zn^2+}/\ce{Zn}) = 0{,}34 - (-0{,}76) = \SI{1,10}{V} \]
La valeur mesurée ($\approx \SI{1,05}{V}$) est en bon accord avec la théorie, l'écart provenant des concentrations et de l'état des surfaces.
\end{exoresolu}

\begin{exercice}[Pour aller plus loin]
On remplace la demi-pile au zinc par une demi-pile $\ce{Fe^2+}/\ce{Fe}$ ($E^\circ = \SI{-0,44}{V}$). Prévois la polarité de la nouvelle pile, écris son équation-bilan de fonctionnement et calcule sa f.é.m. théorique.
\end{exercice}

\begin{corrige}[Exercice]
Le fer est plus réducteur que le cuivre ($-0{,}44 < +0{,}34$) : la lame de fer est le pôle $\ominus$, la lame de cuivre le pôle $\oplus$. Bilan : $\ce{Fe + Cu^2+ -> Fe^2+ + Cu}$. F.é.m. : $E = 0{,}34 - (-0{,}44) = \SI{0,78}{V}$.
\end{corrige}

\courssub{Conclusion}

Nous avons construit une pile Daniell en associant deux demi-piles $\ce{Zn^2+}/\ce{Zn}$ et $\ce{Cu^2+}/\ce{Cu}$ reliées par un pont salin. Le voltmètre a permis d'identifier le zinc comme pôle négatif et le cuivre comme pôle positif, et de mesurer une f.é.m. voisine de \SI{1,10}{V}, conforme à la classification électrochimique. La DEL qui s'allume prouve que la pile transforme l'\textbf{énergie chimique} d'une réaction d'oxydoréduction spontanée en \textbf{énergie électrique}. Ce principe est celui de toutes les piles que tu utilises chaque jour, de la pile de ta torche à la batterie de ton téléphone portable.

\newpage

\coursec{Méthodes et exercices résolus}

\courssub{Prévoir une réaction d'oxydoréduction naturelle : la règle du gamma}

\begin{methode}[Appliquer la règle du gamma]
\begin{enumerate}
\item Repérer les deux couples redox mis en présence et écrire chacun sous la forme $\ce{Ox + n e- <=> Red}$.
\item Classer les couples sur un axe vertical des potentiels : le couple de \cle{potentiel standard le plus élevé} en haut, celui de potentiel le plus bas en bas (oxydants à gauche, réducteurs à droite).
\item Tracer le « gamma » ($\gamma$) : partir de l'oxydant du couple supérieur, descendre en diagonale vers le réducteur du couple inférieur.
\item Conclure : la réaction \cle{naturelle} (spontanée) a lieu entre l'oxydant le plus fort (en haut à gauche) et le réducteur le plus fort (en bas à droite). Les deux autres espèces ne réagissent pas entre elles.
\item Écrire les deux demi-équations, équilibrer les électrons, puis additionner pour obtenir l'équation-bilan (les électrons ne doivent plus apparaître).
\item Vérifier éventuellement avec $E^{\circ} = E^{\circ}_{\text{ox}} - E^{\circ}_{\text{red}} > 0$ : la réaction est spontanée.
\end{enumerate}
\textbf{Réflexe :} si on te donne les $E^{\circ}$, le couple au plus grand $E^{\circ}$ fournit l'\cle{oxydant} qui réagit ; l'autre couple fournit le réducteur.
\end{methode}

\begin{exoresolu}[Plongée d'une lame de zinc dans une solution de cuivre (II)]
On plonge une lame de zinc dans une solution bleue de sulfate de cuivre (II). On donne : $E^{\circ}(\ce{Cu^2+}/\ce{Cu}) = \SI{+0,34}{V}$ et $E^{\circ}(\ce{Zn^2+}/\ce{Zn}) = \SI{-0,76}{V}$.
\begin{enumerate}
\item Prévoir, par la règle du gamma, la réaction qui se produit.
\item Écrire l'équation-bilan et interpréter les observations.
\end{enumerate}
\tcblower
\textbf{1. Classement des couples.} Le couple de plus haut potentiel est \ce{Cu^2+}/\ce{Cu} ($\SI{+0,34}{V}$) ; le couple de plus bas potentiel est \ce{Zn^2+}/\ce{Zn} ($\SI{-0,76}{V}$). On dispose les espèces présentes :
\begin{popfigure}
\begin{tikzpicture}[scale=1.1]
\draw[->,thick,popdark] (0,0) -- (0,3.5) node[above] {$E^{\circ}$ (V)};
% Couple Cu en haut
\node[popblue,font=\bfseries] at (-2,2.8) {\ce{Cu^2+}};
\node[popmuted] at (2,2.8) {\ce{Cu}};
\draw[popblue, thick] (-2.5,2.8) -- (-1.5,2.8);
\draw[popmuted, thick] (1.5,2.8) -- (2.5,2.8);
% Couple Zn en bas
\node[popmuted] at (-2,0.8) {\ce{Zn^2+}};
\node[poporange,font=\bfseries] at (2,0.8) {\ce{Zn}};
\draw[popmuted, thick] (-2.5,0.8) -- (-1.5,0.8);
\draw[poporange, thick] (1.5,0.8) -- (2.5,0.8);
% Gamma
\draw[ultra thick,popgreen,->] (-1.7,2.7) -- (1.7,0.9);
\node[popgreen, font=\bfseries] at (0.5,2.0) {$\gamma$};
% Etiquettes Ox/Red
\node[popdark, font=\small] at (-3.2, 2.8) {Oxydants};
\node[popdark, font=\small] at (3.2, 2.8) {Réducteurs};
\end{tikzpicture}
\legende{Classement des couples \ce{Cu^2+}/\ce{Cu} et \ce{Zn^2+}/\ce{Zn} et tracé du $\gamma$.}
\end{popfigure}
Le gamma relie \ce{Cu^2+} (oxydant le plus fort) à \ce{Zn} (réducteur le plus fort) : la réaction naturelle a lieu entre \ce{Cu^2+} et \ce{Zn}.

\textbf{2. Demi-équations et bilan.}
\begin{align*}
\text{Réduction : } & \ce{Cu^2+ + 2e- -> Cu}\\
\text{Oxydation : } & \ce{Zn -> Zn^2+ + 2e-}
\end{align*}
Les électrons s'éliminent (2 échangés de part et d'autre) :
\[ \ce{Cu^2+ + Zn -> Cu + Zn^2+} \]
Vérification : $E^{\circ} = 0,34 - (-0,76) = \SI{+1,10}{V} > 0$ : réaction spontanée confirmée.

\textbf{Observations :} il se forme un dépôt rougeâtre de cuivre métallique sur la lame de zinc, la lame s'use, et la solution bleue se décolore progressivement (les ions \ce{Cu^2+} bleus disparaissent, remplacés par les ions \ce{Zn^2+} incolores).

\textbf{Conclusion :} le zinc réduit spontanément les ions cuivre (II) selon $\ce{Cu^2+ + Zn -> Cu + Zn^2+}$. C'est exactement cette réaction qui fournira l'énergie de la pile Daniell.
\end{exoresolu}

\courssub{Exploiter la classification électrochimique : polarité et f.é.m d'une pile}

\begin{definition}[Électrode standard à hydrogène (ESH) — Référence universelle]
L'\cle{électrode standard à hydrogène (ESH)} est la demi-pile de référence pour mesurer les potentiels standards $E^{\circ}$. Elle est constituée :
\begin{itemize}
\item d'un électrode de platine inerte platinisée (surface rugueuse pour adsorber $\ce{H2}$) ;
\item baignée dans une solution acide à $[\ce{H3O+}] = \SI{1}{mol.L^{-1}}$ (ou $[\ce{H+}] = \SI{1}{mol.L^{-1}}$) ;
\item traversée par un courant de dihydrogène gazeux $\ce{H2}$ à la pression standard $P = \SI{1}{bar}$ (soit $\SI{1e5}{Pa}$) ;
\item à la température de \SI{25}{\celsius} (\SI{298}{K}).
\end{itemize}
Sa demi-réaction de réduction est : $\ce{2H+ + 2e- <=> H2}$.
Par convention, son potentiel standard est fixé à \textbf{$E^{\circ}(\ce{H+}/\ce{H2}) = \SI{0,00}{V}$}.
Tous les potentiels standards $E^{\circ}$ du tableau de classification sont mesurés relativement à l'ESH (pile \ce{M^{n+}/M} $\parallel$ ESH).
\emph{Notation conventionnelle de l'ESH :} $\ce{Pt} \mid \ce{H2}(g,\SI{1}{bar}) \mid \ce{H+}(aq,\SI{1}{mol.L^{-1}})$.
\end{definition}

\begin{methode}[Déterminer les pôles et la f.é.m d'une pile à partir des $E^{\circ}$]
\begin{enumerate}
\item Identifier les deux couples redox constituant les demi-piles et relever leurs potentiels standard $E^{\circ}$.
\item L'électrode du couple au \cle{plus haut potentiel} est la \cle{cathode} (pôle $+$, siège d'une réduction) ; celle du couple au plus bas potentiel est l'\cle{anode} (pôle $-$, siège d'une oxydation).
\item Calculer la force électromotrice standard :
\[ E^{\circ}_{\text{pile}} = E^{\circ}_{\text{cathode}} - E^{\circ}_{\text{anode}} = E^{\circ}_{(+)} - E^{\circ}_{(-)} \quad (>0) \]
\item Écrire les demi-équations aux électrodes (réduction à la cathode, oxydation à l'anode), puis l'équation-bilan de fonctionnement en égalisant les électrons échangés.
\item Écrire la représentation conventionnelle : $\ominus$ anode $\mid$ anode ionisée $\parallel$ cathode ionisée $\mid$ cathode $\oplus$.
\end{enumerate}
\textbf{Réflexe :} la f.é.m est toujours \emph{positive} ; si ton calcul donne une valeur négative, tu as inversé les pôles.
\end{methode}

\begin{exoresolu}[Pile fer--argent]
On réalise une pile en associant une demi-pile \ce{Fe^2+}/\ce{Fe} ($[\ce{Fe^2+}] = \SI{0,10}{mol.L^{-1}}$) et une demi-pile \ce{Ag+}/\ce{Ag} ($[\ce{Ag+}] = \SI{0,10}{mol.L^{-1}}$), reliées par un pont salin. On donne $E^{\circ}(\ce{Fe^2+}/\ce{Fe}) = \SI{-0,44}{V}$ et $E^{\circ}(\ce{Ag+}/\ce{Ag}) = \SI{+0,80}{V}$.
\begin{enumerate}
\item Déterminer la polarité de la pile et sa f.é.m standard.
\item Écrire les équations aux électrodes et l'équation-bilan de fonctionnement.
\item Donner la représentation conventionnelle de la pile.
\end{enumerate}
\tcblower
\textbf{1. Polarité et f.é.m.} Comparaison des potentiels : $E^{\circ}(\ce{Ag+}/\ce{Ag}) = \SI{+0,80}{V} > E^{\circ}(\ce{Fe^2+}/\ce{Fe}) = \SI{-0,44}{V}$.
\begin{itemize}
\item Électrode d'argent : couple de plus haut potentiel $\Rightarrow$ \cle{cathode}, pôle $\oplus$ (réduction de \ce{Ag+}).
\item Électrode de fer : couple de plus bas potentiel $\Rightarrow$ \cle{anode}, pôle $\ominus$ (oxydation du fer).
\end{itemize}
\[ E^{\circ}_{\text{pile}} = E^{\circ}_{(+)} - E^{\circ}_{(-)} = 0,80 - (-0,44) = \SI{+1,24}{V} \]

\textbf{2. Équations.}
\begin{align*}
\text{Cathode (Ag), réduction : } & \ce{Ag+ + e- -> Ag} \quad (\times 2)\\
\text{Anode (Fe), oxydation : } & \ce{Fe -> Fe^2+ + 2e-}
\end{align*}
On multiplie la demi-équation de l'argent par 2 pour égaliser les électrons, puis on additionne :
\[ \ce{2Ag+ + Fe -> 2Ag + Fe^2+} \]

\textbf{3. Représentation conventionnelle} (anode à gauche, pôle négatif) :
\[ \ominus\ \ce{Fe} \mid \ce{Fe^2+} \parallel \ce{Ag+} \mid \ce{Ag}\ \oplus \]

\textbf{Conclusion :} le fer constitue le pôle négatif, l'argent le pôle positif ; la f.é.m standard vaut \SI{1,24}{V} et la pile fonctionne selon $\ce{2Ag+ + Fe -> 2Ag + Fe^2+}$. Pendant le fonctionnement, l'électrode de fer s'amincit et celle d'argent grossit.
\end{exoresolu}

\courssub{Réaliser et étudier expérimentalement une pile Daniell}

\begin{methode}[Monter une pile Daniell et mesurer ses caractéristiques]
\begin{enumerate}
\item Préparer deux demi-piles : une lame de \cle{zinc} plongeant dans une solution de sulfate de zinc ($\ce{Zn^2+}$, environ \SI{1}{mol.L^{-1}}) et une lame de \cle{cuivre} plongeant dans une solution de sulfate de cuivre ($\ce{Cu^2+}$, environ \SI{1}{mol.L^{-1}}). Décaper les lames à la toile émeri avant usage.
\item Relier les deux solutions par un \cle{pont salin} (tube en U rempli de gel de \ce{KCl} ou de \ce{KNO3}) : il assure le passage du courant ionique et la neutralité électrique des solutions.
\item Brancher un voltmètre aux bornes des lames : la borne COM du voltmètre doit être reliée au pôle $\ominus$ pour lire une valeur positive. Sinon, inverser : la borne reliée à la COM est le pôle négatif.
\item Mesurer la f.é.m : $E \approx \SI{1,10}{V}$ pour la pile Daniell dans les conditions standard.
\item Identifier les pôles : la lame de \cle{zinc} est le pôle $\ominus$ (anode, elle se consomme), la lame de \cle{cuivre} est le pôle $\oplus$ (cathode, du cuivre s'y dépose).
\item Conclure sur le sens du courant : les électrons circulent, dans le circuit extérieur, du zinc vers le cuivre ; le courant conventionnel circule du cuivre vers le zinc.
\end{enumerate}
\end{methode}

\begin{popfigure}
\begin{tikzpicture}[scale=0.85, every node/.style={font=\small}]
% Becher gauche (Zn)
\draw[thick, popdark] (0,0) -- (0,4) -- (2,4) -- (2,0);
\draw[thick, popdark] (-0.2,0) -- (2.2,0); % base
% Solution ZnSO4 (incolore)
\fill[popblueL!30] (0.1,0.1) rectangle (1.9,3.0);
% Electrode Zn
\draw[poporange, very thick] (1,3.5) -- (1,0.5);
\node[poporange, above] at (1,3.7) {\textbf{Zn}};
\node[popdark, left] at (0,2) {$\ce{Zn^2+}$ (\SI{1}{M})};
% Becher droit (Cu)
\draw[thick, popdark] (5,0) -- (5,4) -- (7,4) -- (7,0);
\draw[thick, popdark] (4.8,0) -- (7.2,0);
% Solution CuSO4 (bleue)
\fill[popblue!40] (5.1,0.1) rectangle (6.9,3.0);
% Electrode Cu
\draw[poporange, very thick] (6,3.5) -- (6,0.5);
\node[poporange, above] at (6,3.7) {\textbf{Cu}};
\node[popdark, right] at (7,2) {$\ce{Cu^2+}$ (\SI{1}{M})};
% Pont salin
\draw[poppurple, very thick] (2,1.5) .. controls (3,1.5) and (4,1.5) .. (5,1.5);
\node[poppurple, above] at (3.5,1.8) {Pont salin (\ce{KNO3})};
% Voltmètre
\draw[popdark, thick] (1,4.2) rectangle (2.5,5.0);
\node[popdark] at (1.75,4.6) {V};
% Connexions
\draw[->, thick, popred] (1,3.5) -- (1,4.2) node[left, pos=0.5] {$e^-$};
\draw[thick, popdark] (6,3.5) -- (6,4.2);
\draw[thick, popdark] (1,4.2) -- (1.75,4.2) -- (1.75,5.2) -- (5.25,5.2) -- (5.25,4.2) -- (6,4.2);
% Pôles
\node[popred, font=\bfseries\large] at (0.5,4.5) {$\ominus$};
\node[popblue, font=\bfseries\large] at (6.5,4.5) {$\oplus$};
% Flèches ions dans pont
\draw[->, popblue] (2.2,1.3) -- (3,1.3) node[midway, above] {$\ce{K+}$};
\draw[->, poporange] (4.8,1.7) -- (4,1.7) node[midway, above] {$\ce{NO3-}$};
\end{tikzpicture}
\legende{Schéma fonctionnel de la pile Daniell : anode Zn ($\ominus$), cathode Cu ($\oplus$), pont salin, voltmètre.}
\end{popfigure}

\begin{exoresolu}[Séance de TP : la pile Daniell au laboratoire du lycée]
Au laboratoire, un groupe d'élèves réalise une pile Daniell avec des solutions à \SI{1,0}{mol.L^{-1}}. Le voltmètre, branché avec sa borne COM sur la lame de zinc, affiche $U = \SI{+1,09}{V}$.
\begin{enumerate}
\item Quel est le pôle positif de la pile ? Justifier.
\item Écrire les demi-équations aux électrodes et l'équation-bilan de fonctionnement.
\item Donner la représentation conventionnelle de la pile.
\item Après plusieurs heures de fonctionnement, la masse de la lame de zinc a diminué de \SI{0,65}{g}. Calculer la quantité d'électricité débitée. On donne $M(\ce{Zn}) = \SI{65,4}{g.mol^{-1}}$ et $1\,\mathcal{F} = \SI{96500}{C.mol^{-1}}$.
\end{enumerate}
\tcblower
\textbf{1. Polarité.} Le voltmètre affiche une valeur \emph{positive} alors que sa borne COM est sur le zinc : la borne COM est donc reliée au pôle $\ominus$. Le \cle{zinc est le pôle négatif} et le \cle{cuivre est le pôle positif}. C'est cohérent avec les potentiels standard : $E^{\circ}(\ce{Zn^2+}/\ce{Zn}) = \SI{-0,76}{V} < E^{\circ}(\ce{Cu^2+}/\ce{Cu}) = \SI{+0,34}{V}$.

\textbf{2. Équations de fonctionnement.}
\begin{align*}
\text{Anode (Zn, pôle } \ominus \text{), oxydation : } & \ce{Zn -> Zn^2+ + 2e-}\\
\text{Cathode (Cu, pôle } \oplus \text{), réduction : } & \ce{Cu^2+ + 2e- -> Cu}
\end{align*}
Équation-bilan :
\[ \ce{Zn + Cu^2+ -> Zn^2+ + Cu} \]

\textbf{3. Représentation conventionnelle :}
\[ \ominus\ \ce{Zn} \mid \ce{Zn^2+} \parallel \ce{Cu^2+} \mid \ce{Cu}\ \oplus \]

\textbf{4. Quantité d'électricité débitée.}
Quantité de zinc consommé :
\[ n(\ce{Zn}) = \frac{m}{M} = \frac{0,65}{65,4} = \SI{9,94e-3}{mol} \approx \SI{9,9e-3}{mol} \]
D'après la demi-équation $\ce{Zn -> Zn^2+ + 2e-}$, chaque mole de zinc libère 2 moles d'électrons :
\[ n(\text{e}^-) = 2\,n(\ce{Zn}) = 2 \times \SI{9,94e-3}{} = \SI{1,99e-2}{mol} \]
Quantité d'électricité :
\[ Q = n(\text{e}^-) \times \mathcal{F} = \SI{1,99e-2}{} \times 96\,500 \approx \SI{1,9e3}{C} \]

\textbf{Conclusion :} la pile débite environ \SI{1,9e3}{C} (soit environ \SI{0,53}{A.h}) pendant que \SI{0,65}{g} de zinc sont consommés ; la f.é.m mesurée (\SI{1,09}{V}) est proche de la valeur standard \SI{1,10}{V}.
\end{exoresolu}

\courssub{Écrire les équations aux électrodes et exploiter quantitativement le fonctionnement d'une pile}

\begin{methode}[Bilan de matière et d'électricité dans une pile]
\begin{enumerate}
\item Écrire les demi-équations : \cle{oxydation à l'anode} (pôle $-$), \cle{réduction à la cathode} (pôle $+$). Mnémotechnique : \emph{« AnOx »} (anode--oxydation) et \emph{« RedCat »} (réduction--cathode).
\item Égaliser les nombres d'électrons échangés avant d'additionner pour obtenir l'équation-bilan.
\item Dresser un tableau d'avancement si l'on demande les quantités de matière consommées ou formées.
\item Utiliser les relations quantitatives :
\[ n(\text{e}^-) = \nu \times n(\text{métal consommé}), \qquad Q = n(\text{e}^-)\,\mathcal{F} = I \times \Delta t \]
où $\nu$ est le nombre d'électrons par atome de métal oxydé, $I$ l'intensité (A) et $\Delta t$ la durée (s).
\item Prévoir l'évolution des masses : l'anode \emph{perd} de la masse, la cathode en \emph{gagne} (dépôt métallique).
\end{enumerate}
\end{methode}

\begin{exoresolu}[Autonomie d'une pile zinc--cuivre]
Une pile Daniell débite un courant d'intensité constante $I = \SI{50}{mA}$ pendant $\Delta t = \SI{10}{h}$. La lame de cuivre est très épaisse ; la lame de zinc a une masse initiale $m_0 = \SI{6,54}{g}$. On donne $M(\ce{Zn}) = \SI{65,4}{g.mol^{-1}}$, $M(\ce{Cu}) = \SI{63,5}{g.mol^{-1}}$, $1\,\mathcal{F} = \SI{96500}{C.mol^{-1}}$.
\begin{enumerate}
\item Calculer la quantité d'électricité débitée et la quantité d'électrons échangée.
\item Déterminer la masse de zinc consommée et la masse de cuivre déposée. La lame de zinc est-elle totalement usée ?
\end{enumerate}
\tcblower
\textbf{1. Quantité d'électricité.} Conversion de la durée : $\Delta t = 10 \times 3600 = \SI{3,6e4}{s}$.
\[ Q = I \times \Delta t = \SI{0,050}{A} \times \SI{3,6e4}{s} = \SI{1,8e3}{C} \]
Quantité d'électrons :
\[ n(\text{e}^-) = \frac{Q}{\mathcal{F}} = \frac{1,8\times 10^{3}}{96\,500} = \SI{1,87e-2}{mol} \approx \SI{1,9e-2}{mol} \]

\textbf{2. Masses.} Tableau d'avancement (en mol) pour $\ce{Zn + Cu^2+ -> Zn^2+ + Cu}$ :

\begin{center}
\begin{tabularx}{\linewidth}{|l|X|X|X|X|}
\hline
\rowcolor{popblueL}
État & \ce{Zn} & \ce{Cu^2+} & \ce{Zn^2+} & \ce{Cu} \\
\hline
Initial & $n_0$ & excès & -- & -- \\
\hline
Final & $n_0 - x$ & excès & $x$ & $x$ \\
\hline
\end{tabularx}
\end{center}

D'après $\ce{Zn -> Zn^2+ + 2e-}$, on a $n(\text{e}^-) = 2\,x$, donc :
\[ x = \frac{n(\text{e}^-)}{2} = \frac{\SI{1,87e-2}{}}{2} = \SI{9,33e-3}{mol} \]
Masse de zinc consommée :
\[ m(\ce{Zn}) = x \times M(\ce{Zn}) = \SI{9,33e-3}{} \times 65,4 = \SI{0,61}{g} \]
Masse de cuivre déposée :
\[ m(\ce{Cu}) = x \times M(\ce{Cu}) = \SI{9,33e-3}{} \times 63,5 = \SI{0,59}{g} \]
Masse de zinc restante : $6,54 - 0,61 = \SI{5,93}{g} > 0$.

\textbf{Conclusion :} la pile a consommé \SI{0,61}{g} de zinc et déposé \SI{0,59}{g} de cuivre ; la lame de zinc n'est pas épuisée (il reste \SI{5,93}{g}), la pile peut continuer à fonctionner.
\end{exoresolu}

\courssub{Expliquer la corrosion par la formation de micropiles}

\begin{methode}[Analyser un phénomène de corrosion]
\begin{enumerate}
\item Identifier le métal attaqué et l'environnement : eau (humidité) + dioxygène dissous sont les deux ingrédients indispensables de la corrosion du fer.
\item Repérer les \cle{micropiles} : à la surface du métal (impuretés, zones de contrainte, contact avec un autre métal), de petites zones jouent le rôle d'anodes et de cathodes locales, l'eau salée ou humide servant d'électrolyte.
\item Écrire les demi-équations :
\begin{itemize}
\item Zones anodiques : oxydation du métal, $\ce{Fe -> Fe^2+ + 2e-}$ ;
\item Zones cathodiques : réduction du dioxygène dissous, $\ce{O2 + 2H2O + 4e- -> 4OH-}$.
\end{itemize}
\item Décrire la formation de la rouille : les ions \ce{Fe^2+} et \ce{OH-} migrent, se rencontrent et donnent un hydroxyde qui s'oxyde et se déshydrate en rouille (oxyde de fer (III) hydraté, de formule approchée \ce{Fe2O3, xH2O}).
\item Citer les facteurs aggravants : humidité, sels (bord de mer, ville côtière comme Douala), acidité, contact avec un métal plus noble, rayures.
\end{enumerate}
\textbf{Réflexe :} la corrosion est une oxydation \emph{électrochimique} : elle exige un circuit (anode + cathode + électrolyte). Supprimer l'un des éléments du circuit, c'est protéger le métal.
\end{methode}

\begin{popfigure}
\begin{tikzpicture}[scale=0.9, every node/.style={font=\small}]
% Surface métal
\fill[poporange!30] (0,0) rectangle (8,-0.5);
\draw[thick, popdark] (0,0) -- (8,0);
% Goutte d'eau
\draw[popblue, thick] (2,0) .. controls (2,1.5) and (5,1.5) .. (5,0);
\fill[popblueL!30] (2,0) .. controls (2,1.5) and (5,1.5) .. (5,0) -- cycle;
% Zone anodique (gauche)
\draw[popred, thick, ->] (2.5,-0.2) -- (2.5,0.8);
\node[popred, left] at (2.5,0.8) {Anode : $\ce{Fe -> Fe^2+ + 2e-}$};
\draw[popred, thick] (2.5,0.8) -- (3.5,1.2);
\node[popred] at (3.5,1.4) {$\ce{Fe^2+}$};
% Zone cathodique (droite)
\draw[popblue, thick, ->] (4.5,-0.2) -- (4.5,0.8);
\node[popblue, right] at (4.5,0.8) {Cathode : $\ce{O2 + 2H2O + 4e- -> 4OH-}$};
\draw[popblue, thick] (4.5,0.8) -- (3.5,1.2);
\node[popblue] at (3.5,1.0) {$\ce{OH-}$};
% Électrons dans le métal
\draw[->, thick, popdark] (2.5,-0.3) -- (4.5,-0.3) node[midway, below] {$e^-$ (dans le métal)};
% Formation rouille
\node[poppurple, font=\bfseries] at (3.5, -1.2) {Formation de $\ce{Fe(OH)2}$ puis rouille $\ce{Fe2O3, xH2O}$};
\draw[poppurple, ->] (3.5,-0.8) -- (3.5,-1.1);
% O2 dissous
\node[popdark] at (6.5, 1.0) {$\ce{O2}$ dissous};
\draw[->, popdark] (6,0.8) -- (5.2,0.8);
\end{tikzpicture}
\legende{Mécanisme de corrosion du fer : micropiles à la surface, migration des ions et formation de rouille.}
\end{popfigure}

\begin{exoresolu}[Pourquoi le portail en fer rouille-t-il plus vite à Douala ?]
Un portail en fer, peint, est installé à Douala (climat chaud et humide, air marin chargé de sel). Au bout de quelques mois, des cloques de rouille apparaissent aux endroits où la peinture est écaillée.
\begin{enumerate}
\item Expliquer, à l'aide du modèle des micropiles, pourquoi la rouille se forme aux endroits écaillés.
\item Écrire les demi-équations mises en jeu.
\item Justifier le rôle aggravant du sel marin.
\end{enumerate}
\tcblower
\textbf{1. Mécanisme.} Là où la peinture est intacte, le fer est isolé de l'air et de l'eau : aucune micropile ne peut fonctionner. Aux endroits écaillés, le fer est en contact direct avec l'eau de pluie ou l'humidité de l'air contenant du dioxygène dissous. Des zones de la surface (impuretés, zones déformées) deviennent des \cle{anodes} locales où le fer s'oxyde, tandis que d'autres zones deviennent des \cle{cathodes} où le dioxygène se réduit : de véritables micropiles se mettent en fonctionnement, l'eau humide jouant le rôle d'électrolyte. Le fer est alors consommé localement : c'est la corrosion.

\textbf{2. Demi-équations.}
\begin{align*}
\text{Zones anodiques (oxydation du fer) : } & \ce{Fe -> Fe^2+ + 2e-}\\
\text{Zones cathodiques (réduction de } \ce{O2}\text{) : } & \ce{O2 + 2H2O + 4e- -> 4OH-}
\end{align*}
Les ions \ce{Fe^2+} et \ce{OH-} se rencontrent dans la goutte d'eau ; après oxydation par le dioxygène et déshydratation, on obtient la rouille, oxyde de fer (III) hydraté \ce{Fe2O3, xH2O}, friable et non protecteur.

\textbf{3. Rôle du sel.} Le chlorure de sodium de l'air marin se dissout dans le film d'eau : les ions \ce{Na+} et \ce{Cl-} augmentent fortement la \cle{conductivité} de l'électrolyte. Les micropiles débitent alors un courant plus intense : la corrosion est accélérée. La chaleur et l'humidité permanente de Douala entretiennent en outre le film d'eau à la surface du métal.

\textbf{Conclusion :} la rouille résulte du fonctionnement de micropiles électrochimiques à la surface du fer exposé ; l'air salin et humide de Douala, en rendant l'électrolyte plus conducteur et permanent, accélère fortement le phénomène.
\end{exoresolu}

\courssub{Choisir et décrire une méthode de protection contre la corrosion}

\begin{methode}[Choisir une protection adaptée]
\begin{enumerate}
\item \cle{Protection par revêtement} (barrière) : peinture, vernis, graisse, plastique, émail. Le métal est isolé de l'eau et du dioxygène. Efficace tant que le revêtement est intact ; à renouveler.
\item \cle{Protection par métallisation} : recouvrir le fer d'un métal plus résistant (chromage, étamage des boîtes de conserve) ou d'un métal \emph{plus réducteur} (galvanisation : revêtement de zinc).
\item \cle{Galvanisation} : le zinc, plus réducteur que le fer ($E^{\circ}(\ce{Zn^2+}/\ce{Zn}) < E^{\circ}(\ce{Fe^2+}/\ce{Fe})$), s'oxyde à la place du fer même si le revêtement est rayé : le zinc joue le rôle d'\cle{anode sacrificielle}.
\item \cle{Protection cathodique} : fixer sur la structure à protéger un bloc de métal plus réducteur (magnésium, zinc) qui se corrode à sa place (anode sacrificielle) — utilisée pour les coques de navires, les pipelines, les réservoirs enterrés. Variante : imposer un courant électrique qui maintient la structure au potentiel de cathode (protection par courant imposé).
\item Justifier le choix par la classification électrochimique : le métal protecteur doit avoir un potentiel standard \emph{inférieur} à celui du métal protégé pour être oxydé à sa place.
\end{enumerate}
\end{methode}

\begin{popfigure}
\begin{tikzpicture}[scale=0.85, every node/.style={font=\small}]
% Pipeline (cathode)
\draw[popblue, very thick] (0,0) -- (8,0);
\node[popblue, above] at (4,0.3) {Pipeline en Fe (Cathode $\oplus$)};
% Anode sacrificielle Mg
\draw[poporange, thick] (4,-2) rectangle (5,-3.5);
\node[poporange, below] at (4.5,-3.8) {Bloc Mg (Anode $\ominus$)};
% Fil de connexion
\draw[popred, thick, ->] (4.5,-0.2) -- (4.5,-2) node[midway, right] {$e^-$};
% Sol (électrolyte)
\fill[popgreenL!30] (0,-0.5) rectangle (8,-4);
\draw[thick, popdark] (0,-0.5) -- (8,-0.5);
% Réactions
\node[popred, left] at (-0.5,-1.5) {Anode (Mg) : $\ce{Mg -> Mg^2+ + 2e-}$};
\node[popblue, right] at (8.5,-1.5) {Cathode (Fe) : $\ce{O2 + 2H2O + 4e- -> 4OH-}$};
% Ions
\draw[->, poporange] (4.5,-2.5) -- (3.5,-2.5) node[midway, above] {$\ce{Mg^2+}$};
\draw[->, popblue] (4.5,-1.5) -- (5.5,-1.5) node[midway, above] {$\ce{OH-}$};
\end{tikzpicture}
\legende{Protection cathodique d'un pipeline par anode sacrificielle en magnésium.}
\end{popfigure}

\begin{exoresolu}[Protéger une canalisation enterrée]
Une canalisation en fer transportant de l'eau est enterrée dans un sol humide. Pour la protéger, un technicien propose deux solutions : (a) la recouvrir d'un revêtement de cuivre ; (b) la relier à des blocs de magnésium enterrés à intervalles réguliers. On donne : $E^{\circ}(\ce{Mg^2+}/\ce{Mg}) = \SI{-2,37}{V}$ ; $E^{\circ}(\ce{Fe^2+}/\ce{Fe}) = \SI{-0,44}{V}$ ; $E^{\circ}(\ce{Cu^2+}/\ce{Cu}) = \SI{+0,34}{V}$.
\begin{enumerate}
\item Montrer que la solution (a) est dangereuse si le revêtement est rayé.
\item Justifier l'efficacité de la solution (b) et écrire les demi-équations mises en jeu.
\item Nommer cette méthode de protection.
\end{enumerate}
\tcblower
\textbf{1. Revêtement de cuivre.} On compare les potentiels : $E^{\circ}(\ce{Cu^2+}/\ce{Cu}) = \SI{+0,34}{V} > E^{\circ}(\ce{Fe^2+}/\ce{Fe}) = \SI{-0,44}{V}$. Si le revêtement de cuivre est rayé, le fer et le cuivre sont en contact dans le sol humide (électrolyte) : une micropile se forme où le \cle{fer}, métal le plus réducteur, devient l'\cle{anode} et s'oxyde :
\[ \ce{Fe -> Fe^2+ + 2e-} \]
tandis que le cuivre sert de cathode. La corrosion du fer est alors \emph{accélérée} au niveau de la rayure : la solution (a) est à rejeter.

\textbf{2. Blocs de magnésium.} On a $E^{\circ}(\ce{Mg^2+}/\ce{Mg}) = \SI{-2,37}{V} < E^{\circ}(\ce{Fe^2+}/\ce{Fe}) = \SI{-0,44}{V}$ : le magnésium est \emph{plus réducteur} que le fer. Dans la micropile formée avec le sol humide, c'est le magnésium qui devient l'anode et qui s'oxyde :
\begin{align*}
\text{Anode (bloc de Mg) : } & \ce{Mg -> Mg^2+ + 2e-}\\
\text{Cathode (canalisation en Fe) : } & \ce{O2 + 2H2O + 4e- -> 4OH-}
\end{align*}
Le fer, devenu cathode, n'est plus oxydé : il est protégé. Le magnésium se consomme lentement à sa place et doit être remplacé périodiquement.

\textbf{3. Nom de la méthode.} Il s'agit de la \cle{protection cathodique} par \cle{anode sacrificielle} (ou anode soluble).

\textbf{Conclusion :} seule la solution (b) protège durablement la canalisation : le magnésium, plus réducteur que le fer, se corrode à sa place ; un revêtement de cuivre rayé aggraverait au contraire la corrosion.
\end{exoresolu}

\begin{attention}[Les 5 erreurs qui coûtent des points au Bac]
\begin{enumerate}
\item \textbf{Inverser anode et cathode.} L'\cle{anode} est le siège de l'\cle{oxydation} (pôle $\ominus$ de la pile) ; la \cle{cathode} est le siège de la \cle{réduction} (pôle $\oplus$). Retiens : \emph{AnOx} et \emph{RedCat}. Attention : dans une pile, l'anode est le pôle négatif ; ne confonds pas avec l'électrolyse.
\item \textbf{Calculer une f.é.m négative.} La f.é.m d'une pile est toujours positive : $E = E^{\circ}_{(+)} - E^{\circ}_{(-)} = E^{\circ}_{\text{cathode}} - E^{\circ}_{\text{anode}}$. Une valeur négative signifie que tu as inversé les pôles.
\item \textbf{Laisser des électrons dans l'équation-bilan.} Avant d'additionner les demi-équations, multiplie-les pour égaliser les électrons échangés (ex. $\ce{Ag+ + e- -> Ag}$ à multiplier par 2 face à $\ce{Fe -> Fe^2+ + 2e-}$). Les électrons ne doivent jamais apparaître dans le bilan final.
\item \textbf{Oublier le rôle du pont salin.} Le pont salin ne transporte pas d'électrons : il assure la circulation des \emph{ions} et maintient l'électroneutralité des solutions. Sans lui, la pile ne débite pas. Et les électrons circulent toujours dans le circuit extérieur, jamais dans les solutions.
\item \textbf{Croire que tout revêtement métallique protège.} Un revêtement protège durablement le fer seulement s'il est \emph{intact} (peinture, chrome) ou constitué d'un métal \emph{plus réducteur} que le fer (zinc, galvanisation). Un métal plus noble que le fer (cuivre, étain rayé) \emph{accélère} la corrosion dès que le revêtement est défaillant. Justifie toujours ta réponse par la comparaison des potentiels standard.
\end{enumerate}
\end{attention}

\newpage

\coursec{Exercices}

\courssub{Je vérifie mes connaissances}

\begin{exercice}[QCM : notions sur les piles]
Pour chaque question, choisis la (ou les) bonne(s) réponse(s).
\begin{enumerate}
\item Dans une pile électrochimique, l'anode est :
\begin{enumerate}
\item[a)] le siège de la réduction ;
\item[b)] le siège de l'oxydation ;
\item[c)] toujours le pôle positif ;
\item[d)] l'électrode qui fournit des électrons au circuit extérieur.
\end{enumerate}
\item Dans la pile Daniell, le pont salin sert à :
\begin{enumerate}
\item[a)] laisser passer les électrons ;
\item[b)] assurer la neutralité électrique des solutions ;
\item[c)] fermer le circuit électrique (conduction ionique) ;
\item[d)] augmenter la f.é.m. de la pile.
\end{enumerate}
\item On donne $E^\circ(\ce{Cu^2+}/\ce{Cu}) = \SI{+0,34}{V}$ et $E^\circ(\ce{Zn^2+}/\ce{Zn}) = \SI{-0,76}{V}$. La f.é.m. standard de la pile Daniell vaut :
\begin{enumerate}
\item[a)] \SI{0,42}{V} ; \item[b)] \SI{1,10}{V} ; \item[c)] \SI{-1,10}{V} ; \item[d)] \SI{2,20}{V}.
\end{enumerate}
\item La corrosion du fer est :
\begin{enumerate}
\item[a)] une réduction du fer ;
\item[b)] une oxydation du fer ;
\item[c)] favorisée par la sécheresse ;
\item[d)] expliquée par la formation de micropiles.
\end{enumerate}
\end{enumerate}
\end{exercice}

\begin{exercice}[Vrai ou faux, avec justification]
Réponds par \textbf{Vrai} ou \textbf{Faux}, puis justifie ta réponse en une ou deux phrases.
\begin{enumerate}
\item Le réducteur d'un couple est d'autant plus fort que le potentiel standard du couple est élevé.
\item Dans une pile, les électrons circulent dans le circuit extérieur du pôle négatif vers le pôle positif.
\item Une lame d'argent plongée dans une solution de nitrate de cuivre(II) se recouvre de cuivre.
\item Le zinc protège le fer parce qu'il est plus réducteur que le fer.
\end{enumerate}
\end{exercice}

\begin{exercice}[Texte à trous]
Recopie et complète le texte suivant avec les mots ou expressions de la liste : \emph{anode ; cathode ; oxydation ; réduction ; pont salin ; f.é.m. ; référence ; galvanisation}.

Dans la pile Daniell, la lame de zinc constitue l'\ldots\ldots\ldots, siège de l'\ldots\ldots\ldots : c'est le pôle \ldots\ldots\ldots de la pile. La lame de cuivre constitue la \ldots\ldots\ldots, siège de la \ldots\ldots\ldots. Les deux demi-piles sont reliées par un \ldots\ldots\ldots qui assure le passage du courant dans les solutions. La différence de potentiel entre les deux pôles, mesurée en circuit ouvert, est appelée \ldots\ldots\ldots de la pile. Les potentiels standard des couples sont mesurés par rapport à l'électrode à hydrogène, qui sert d'électrode de \ldots\ldots\ldots. Le procédé qui consiste à recouvrir le fer d'une couche de zinc s'appelle la \ldots\ldots\ldots.
\end{exercice}

\begin{exercice}[Définitions]
Définis précisément les termes suivants, en donnant pour chacun un exemple tiré du cours :
\begin{enumerate}
\item demi-pile métallique \ce{M^{n+}}/\ce{M} ;
\item pile électrochimique ;
\item corrosion d'un métal ;
\item protection cathodique.
\end{enumerate}
\end{exercice}

\courssub{Je m'entraîne}

\begin{exercice}[Règle du gamma : fer et ions cuivre(II)]
On plonge une lame de fer dans un bécher contenant \SI{200}{mL} d'une solution de sulfate de cuivre(II). On donne : $E^\circ(\ce{Fe^2+}/\ce{Fe}) = \SI{-0,44}{V}$ et $E^\circ(\ce{Cu^2+}/\ce{Cu}) = \SI{+0,34}{V}$ ; $M(\ce{Fe}) = \SI{55,8}{g.mol^{-1}}$ ; $M(\ce{Cu}) = \SI{63,5}{g.mol^{-1}}$.
\begin{enumerate}
\item À l'aide de la règle du gamma, prévoir la réaction naturelle qui se produit. Écrire les deux demi-équations et l'équation-bilan.
\item Quelles observations expérimentales permettent de vérifier cette prévision ?
\item La réaction s'arrête lorsque \SI{0,05}{mol} d'ions \ce{Cu^2+} ont disparu. Calculer la variation de masse de la lame de fer (on précisera s'il s'agit d'une augmentation ou d'une diminution).
\end{enumerate}
\end{exercice}

\begin{exercice}[Classement des métaux et action des acides]
On donne les potentiels standard des couples suivants :
\begin{center}
\begin{tabular}{lccccc}
\toprule
Couple & \ce{Al^3+}/\ce{Al} & \ce{Zn^2+}/\ce{Zn} & \ce{Fe^2+}/\ce{Fe} & \ce{Cu^2+}/\ce{Cu} & \ce{Ag+}/\ce{Ag} \\
\midrule
$E^\circ$ (V) & $-1,66$ & $-0,76$ & $-0,44$ & $+0,34$ & $+0,80$ \\
\bottomrule
\end{tabular}
\end{center}
On donne de plus $E^\circ(\ce{H3O+}/\ce{H2}) = \SI{0,00}{V}$.
\begin{enumerate}
\item Classer les métaux \ce{Al}, \ce{Zn}, \ce{Fe}, \ce{Cu}, \ce{Ag} par pouvoir réducteur croissant. Justifier.
\item Parmi ces cinq métaux, lesquels sont attaqués par une solution d'acide chlorhydrique avec dégagement de dihydrogène ? Justifier à l'aide de la règle du gamma.
\item Écrire l'équation-bilan de la réaction entre le zinc et l'acide chlorhydrique.
\end{enumerate}
\end{exercice}

\begin{exercice}[Bilan quantitatif d'une pile Daniell]
Une pile Daniell débite dans un circuit extérieur. Après une durée de fonctionnement $t$, on constate que la masse de la lame de zinc a diminué de \SI{1,31}{g}. On donne : $M(\ce{Zn}) = \SI{65,4}{g.mol^{-1}}$ ; $M(\ce{Cu}) = \SI{63,5}{g.mol^{-1}}$ ; $1\,\mathrm{F} = \SI{96500}{C.mol^{-1}}$.
\begin{enumerate}
\item Écrire les demi-équations aux deux électrodes et l'équation-bilan de fonctionnement de la pile.
\item Calculer la quantité de matière de zinc consommée, puis la quantité d'électricité $Q$ débitée par la pile.
\item Calculer la masse de cuivre déposée sur la lame de cuivre pendant la même durée.
\item En déduire la variation totale de la masse des deux électrodes.
\end{enumerate}
\end{exercice}

\begin{exercice}[Micropiles et corrosion]
Une goutte d'eau salée repose à la surface d'une plaque de fer exposée à l'air humide.
\begin{enumerate}
\item Expliquer, en t'appuyant sur le mécanisme des micropiles, pourquoi le fer se corrode au centre de la goutte alors que le dioxygène est réduit en périphérie.
\item Écrire la demi-équation d'oxydation du fer et la demi-équation de réduction du dioxygène en milieu aqueux ($\ce{O2} + \ce{2H2O} + \ce{4e-} \ce{->} \ce{4HO-}$).
\item Citer trois facteurs qui accélèrent la corrosion du fer et proposer, pour chacun, une parade pratique.
\end{enumerate}
\end{exercice}

\courssub{Je me prépare au Bac}

\begin{exercice}[Type Bac : la pile argent--plomb]
On donne les potentiels standard : $E^\circ(\ce{Ag+}/\ce{Ag}) = \SI{+0,80}{V}$ et $E^\circ(\ce{Pb^2+}/\ce{Pb}) = \SI{-0,13}{V}$. On réalise au laboratoire une pile en associant deux demi-piles : une lame d'argent plongeant dans \SI{100}{mL} d'une solution de nitrate d'argent de concentration \SI{0,10}{mol.L^{-1}}, et une lame de plomb plongeant dans \SI{100}{mL} d'une solution de nitrate de plomb(II) de concentration \SI{0,10}{mol.L^{-1}}. Les deux demi-piles sont reliées par un pont salin au nitrate de potassium. On donne : $M(\ce{Ag}) = \SI{107,9}{g.mol^{-1}}$ ; $M(\ce{Pb}) = \SI{207,2}{g.mol^{-1}}$ ; $1\,\mathrm{F} = \SI{96500}{C.mol^{-1}}$.
\begin{enumerate}
\item À l'aide de la règle du gamma, prévoir la réaction naturelle entre les deux couples. Écrire les deux demi-équations et l'équation-bilan de fonctionnement de la pile.
\item Indiquer la polarité de chaque électrode et préciser l'anode et la cathode.
\item Donner la représentation conventionnelle de la pile.
\item Calculer la f.é.m. standard de cette pile.
\item Faire un schéma fonctionnel annoté de la pile en indiquant le sens de circulation des électrons dans le circuit extérieur et le sens de migration des ions dans le pont salin.
\item La pile débite jusqu'à ce que la concentration en ions \ce{Ag+} soit divisée par dix. Calculer la variation de masse de la lame d'argent et la variation de masse de la lame de plomb.
\end{enumerate}
\end{exercice}

\begin{exercice}[Type Bac : tôle galvanisée et boîte de conserve]
Au Cameroun, les tôles de toiture sont en fer galvanisé (recouvert de zinc), tandis que les boîtes de conserve sont en fer étamé (recouvert d'étain). On donne : $E^\circ(\ce{Zn^2+}/\ce{Zn}) = \SI{-0,76}{V}$ ; $E^\circ(\ce{Fe^2+}/\ce{Fe}) = \SI{-0,44}{V}$ ; $E^\circ(\ce{Sn^2+}/\ce{Sn}) = \SI{-0,14}{V}$.
\begin{enumerate}
\item Classer les trois métaux \ce{Zn}, \ce{Fe}, \ce{Sn} par pouvoir réducteur croissant.
\item Une tôle galvanisée est rayée : le fer est mis à nu au contact du zinc, en présence d'eau de pluie. Montrer, à l'aide de la règle du gamma et du mécanisme des micropiles, que c'est le zinc qui s'oxyde et non le fer. Écrire les demi-équations en jeu. Conclure sur la protection du fer.
\item Une boîte de conserve en fer étamé est rayée de la même façon. Montrer que, cette fois, c'est le fer qui s'oxyde préférentiellement. Expliquer pourquoi la boîte rouille alors rapidement.
\item Nommer le type de protection assuré par le zinc dans le cas de la tôle galvanisée.
\item Proposer deux autres méthodes de protection du fer contre la corrosion, en précisant leur principe.
\end{enumerate}
\end{exercice}

\begin{exercice}[Type Bac : protection cathodique d'un pipeline]
Un pipeline en acier (assimilé à du fer) transportant des produits pétroliers est enterré dans un sol humide. Pour le protéger de la corrosion, on fixe à intervalles réguliers des blocs de magnésium reliés électriquement à la canalisation. On donne : $E^\circ(\ce{Mg^2+}/\ce{Mg}) = \SI{-2,37}{V}$ ; $E^\circ(\ce{Fe^2+}/\ce{Fe}) = \SI{-0,44}{V}$ ; $M(\ce{Mg}) = \SI{24,3}{g.mol^{-1}}$ ; $M(\ce{Fe}) = \SI{55,8}{g.mol^{-1}}$ ; $1\,\mathrm{F} = \SI{96500}{C.mol^{-1}}$.
\begin{enumerate}
\item Nommer la méthode de protection utilisée. Préciser le nom donné aux blocs de magnésium.
\item Le couple magnésium--fer, au contact du sol humide, constitue une micropile. À l'aide de la règle du gamma, déterminer quel métal joue le rôle d'anode et quel métal joue le rôle de cathode. Écrire les demi-équations correspondantes.
\item Justifier le choix du magnésium plutôt que du cuivre ($E^\circ(\ce{Cu^2+}/\ce{Cu}) = \SI{+0,34}{V}$) pour protéger le pipeline.
\item Expliquer pourquoi les blocs de magnésium doivent être remplacés périodiquement.
\item Un bloc de magnésium de masse $m = \SI{4,86}{kg}$ est entièrement consommé après plusieurs années de service. Calculer la quantité d'électricité totale qu'il a fournie, puis la masse de fer qui aurait été corrodée en l'absence de protection pour la même quantité d'électricité.
\end{enumerate}
\end{exercice}

\begin{savaistu}[Le saviez-vous ? : La corrosion au Cameroun]
Au Cameroun, la corrosion est un enjeu économique majeur. L'humidité tropicale ( pluviométrie élevée, hygrométrie \SI{>80}{\%}) et la proximité de la mer (aérosols salins sur la côte, Kribi, Douala, Limbé) créent un électrolyte naturel très conducteur, transformant chaque rayure, chaque soudure, chaque hétérogénéité en micropile active.
\begin{itemize}
\item \textbf{Toitures :} Les tôles galvanisées (zinc) résistent bien, mais aux coupes et trous de vis, le fer nu est protégé par le zinc environnant (anode sacrificielle). C'est pourquoi il faut des vis \textbf{galvanisées} ou en inox (sinon la vis rouille et tache la tôle).
\item \textbf{Pétrole (SONARA) :} Les pipelines enterrés et les cuves de stockage sont protégés par \textbf{courant imposé} (anodes en oxydes métalliques mixtes MMO alimentées par redresseurs) car les anodes Mg s'usent trop vite pour de grandes structures.
\item \textbf{Eau (CAMWATER) :} Les canalisations en fonte ductile sont gainées (polyéthylène) et parfois protégées cathodiquement.
\item \textbf{Pharmacopée :} Le \emph{bili-bili} (bière de mil) ou le \emph{vin de palme} sont acides (pH 3-4). Les récipients métalliques (bidons, fûts) doivent être en inox ou plastifiés, car l'acidité attaque le fer ($\ce{Fe + 2H+ -> Fe^2+ + H2}$) et dissout le zinc (goût métallique, toxicité).
\end{itemize}
\end{savaistu}

\begin{attention}[Pièges classiques au Bac TI]
\begin{itemize}
\item \textbf{Polarité anode/cathode :} Dans une \textbf{pile} (générateur) : Anode = Pôle \textbf{Négatif} (-) ; Cathode = Pôle \textbf{Positif} (+). Dans une \textbf{électrolyse} (récepteur) : c'est l'inverse ! Ne jamais confondre.
\item \textbf{Règle du gamma :} "L'oxydant du couple du haut réagit avec le réducteur du couple du bas". Vérifier toujours la position relative dans le tableau (haut/bas) et non la valeur numérique seule.
\item \textbf{Équation-bilan :} Vérifier l'équilibre des \textbf{charges} et des \textbf{atomes}. Multiplier les demi-équations pour égaler les électrons.
\item \textbf{Masse électrodes :} Anode : masse \textbf{diminue} (métal $\to$ ions en solution). Cathode : masse \textbf{augmente} (ions $\to$ métal déposé). Variation totale $\neq 0$ (différence de masses molaires).
\item \textbf{Protection cathodique :} Le métal à protéger \textbf{doit être la cathode}. L'anode sacrificielle doit avoir un $E^\circ$ \textbf{plus bas} (plus négatif).
\end{itemize}
\end{attention}

\newpage

\coursec{Corrigés des exercices}

\courssub{Je vérifie mes connaissances}

\begin{corrige}[Exercice 1 — QCM : notions sur les piles]
\begin{enumerate}
\item \textbf{Réponses b) et d).} Par définition, l'\cle{anode} est le siège de l'\cle{oxydation} (perte d'électrons) ; elle fournit donc les électrons au circuit extérieur. Dans une pile (générateur), l'anode est le pôle \textbf{négatif} : la réponse c) est fausse. La réponse a) est fausse : la réduction a lieu à la cathode.
\item \textbf{Réponses b) et c).} Le \cle{pont salin} assure la conduction ionique (il ferme le circuit électrique à l'intérieur de la pile) et maintient la neutralité électrique des deux demi-piles : les anions migrent vers l'anode, les cations vers la cathode. Il ne laisse jamais passer les électrons (a) faux) et ne modifie pas la f.é.m. (d) faux).
\item \textbf{Réponse b).} La f.é.m. standard se calcule par :
\[ E^\circ_{\text{pile}} = E^\circ_{\text{cathode}} - E^\circ_{\text{anode}} = E^\circ(\ce{Cu^2+}/\ce{Cu}) - E^\circ(\ce{Zn^2+}/\ce{Zn}) = 0,34 - (-0,76) = \SI{1,10}{V}. \]
\item \textbf{Réponses b) et d).} La corrosion est l'\cle{oxydation} du métal par le dioxygène de l'air en présence d'eau ; elle est donc favorisée par l'\textbf{humidité}. L'option c) correspond à un facteur ne favorisant pas la corrosion (ex. : air sec, milieu basique passivant). Son mécanisme s'explique par la formation de \cle{micropiles} locales à la surface du métal (d) vrai).
\end{enumerate}
\end{corrige}

\begin{corrige}[Exercice 2 — Vrai ou faux, avec justification]
\begin{enumerate}
\item \textbf{Faux.} C'est exactement l'inverse : plus le potentiel standard $E^\circ$ d'un couple $\ce{Ox}/\ce{Red}$ est \textbf{bas} (négatif), plus le réducteur $\ce{Red}$ est fort, car il a une grande tendance à céder ses électrons. Exemple : $\ce{Zn}$ ($E^\circ = \SI{-0,76}{V}$) est plus réducteur que $\ce{Cu}$ ($E^\circ = \SI{+0,34}{V}$).
\item \textbf{Vrai.} Dans une pile fonctionnant en générateur, l'anode est le pôle négatif (source d'électrons) et la cathode le pôle positif. Les électrons circulent donc dans le circuit extérieur du pôle $(-)$ vers le pôle $(+)$, c'est-à-dire de l'anode vers la cathode.
\item \textbf{Faux.} On compare : $E^\circ(\ce{Ag+}/\ce{Ag}) = \SI{+0,80}{V} > E^\circ(\ce{Cu^2+}/\ce{Cu}) = \SI{+0,34}{V}$. L'argent est un réducteur \textbf{plus faible} que le cuivre. Par la règle du gamma, l'oxydant $\ce{Cu^2+}$ se trouve \emph{au-dessous} du réducteur $\ce{Ag}$ dans la classification : la réaction $\ce{2Ag + Cu^2+ -> 2Ag+ + Cu}$ n'est \textbf{pas spontanée}. Aucun dépôt de cuivre n'apparaît.
\item \textbf{Vrai.} $E^\circ(\ce{Zn^2+}/\ce{Zn}) = \SI{-0,76}{V} < E^\circ(\ce{Fe^2+}/\ce{Fe}) = \SI{-0,44}{V}$ : le zinc est un réducteur plus fort que le fer. Au contact du fer en milieu humide, c'est le zinc qui s'oxyde préférentiellement : il joue le rôle d'\cle{anode sacrificielle} et le fer, devenu cathode, est protégé.
\end{enumerate}
\end{corrige}

\begin{corrige}[Exercice 3 — Texte à trous]
Dans la pile Daniell, la lame de zinc constitue \textbf{l'anode}, siège de \textbf{l'oxydation} : c'est le pôle \textbf{négatif} de la pile. La lame de cuivre constitue \textbf{la cathode}, siège de \textbf{la réduction}. Les deux demi-piles sont reliées par un \textbf{pont salin} qui assure le passage du courant dans les solutions. La différence de potentiel entre les deux pôles, mesurée en circuit ouvert, est appelée \textbf{f.é.m.} de la pile. Les potentiels standard des couples sont mesurés par rapport à l'électrode à hydrogène, qui sert d'électrode de \textbf{référence}. Le procédé qui consiste à recouvrir le fer d'une couche de zinc s'appelle \textbf{la galvanisation}.
\end{corrige}

\begin{attention}[Mot non placé]
Le mot \emph{négatif} ne figurait pas dans la liste proposée : il fallait le déduire de la connaissance de la pile Daniell. Tous les autres mots de la liste sont utilisés une et une seule fois.
\end{attention}

\begin{corrige}[Exercice 4 — Définitions]
\begin{enumerate}
\item \textbf{Demi-pile métallique $\ce{M^{n+}}/\ce{M}$ :} ensemble constitué d'une lame de métal $\ce{M}$ (conducteur électronique) plongeant dans une solution contenant les cations $\ce{M^{n+}}$ de ce même métal (conducteur ionique) ; elle met en jeu le couple redox $\ce{M^{n+}}/\ce{M}$. \emph{Exemple :} lame de zinc plongeant dans une solution de sulfate de zinc, demi-pile $\ce{Zn^2+}/\ce{Zn}$.
\item \textbf{Pile électrochimique :} dispositif qui convertit l'énergie chimique d'une réaction d'oxydoréduction \textbf{spontanée} en énergie électrique. Elle associe deux demi-piles reliées par un pont salin (conduction ionique interne) et par un circuit extérieur (conduction électronique). \emph{Exemple :} la pile Daniell $\ce{Zn}|\ce{Zn^2+}||\ce{Cu^2+}|\ce{Cu}$.
\item \textbf{Corrosion d'un métal :} dégradation d'un métal par \textbf{oxydation électrochimique spontanée} au contact de son environnement (air humide, eau, sol). \emph{Exemple :} formation de rouille $\ce{Fe2O3}\cdot n\,\ce{H2O}$ sur une tôle de fer exposée à la pluie à Douala.
\item \textbf{Protection cathodique :} méthode qui consiste à imposer au métal à protéger le rôle de \textbf{cathode} (siège de la réduction), soit en le reliant à un métal plus réducteur qui se consume à sa place (anode sacrificielle), soit en lui imposant un courant électrique (courant imposé). \emph{Exemple :} blocs de magnésium fixés sur un pipeline en acier enterré.
\end{enumerate}
\end{corrige}

\courssub{Je m'entraîne}

\begin{corrige}[Exercice 5 — Règle du gamma : fer et ions cuivre(II)]
\begin{enumerate}
\item \textbf{Prévision par la règle du gamma.} On place les deux couples sur l'axe des potentiels :
\[ E^\circ(\ce{Cu^2+}/\ce{Cu}) = \SI{+0,34}{V} > E^\circ(\ce{Fe^2+}/\ce{Fe}) = \SI{-0,44}{V}. \]
L'oxydant le plus fort, $\ce{Cu^2+}$ (couple de potentiel le plus élevé), réagit spontanément avec le réducteur le plus fort, $\ce{Fe}$ (couple de potentiel le plus bas) : la réaction naturelle forme un « gamma » dans le tableau de classification.
\begin{align*}
\text{Oxydation du fer :}\quad & \ce{Fe(s) -> Fe^2+(aq) + 2e-} \\
\text{Réduction des ions cuivre :}\quad & \ce{Cu^2+(aq) + 2e- -> Cu(s)}
\end{align*}
Les électrons s'éliminent (2 de part et d'autre) ; l'équation-bilan s'écrit :
\[ \ce{Fe(s) + Cu^2+(aq) -> Fe^2+(aq) + Cu(s)} \]
\item \textbf{Observations expérimentales :}
\begin{itemize}
\item la lame de fer \textbf{s'amincit} et sa surface devient rugueuse (dissolution du fer) ;
\item un dépôt \textbf{rouge-orangé} de cuivre métallique se forme sur la lame et au fond du bécher ;
\item la teinte \textbf{bleue} de la solution (ions $\ce{Cu^2+}$) pâlit progressivement, tandis qu'apparaît une coloration \textbf{vert pâle} due aux ions $\ce{Fe^2+}$.
\end{itemize}
\item \textbf{Variation de masse de la lame de fer.} D'après la stœchiométrie de l'équation-bilan, 1 mol de $\ce{Cu^2+}$ consommée correspond à 1 mol de $\ce{Fe}$ dissous :
\[ n(\ce{Fe}) = n(\ce{Cu^2+})_{\text{disparus}} = \SI{0,05}{mol}. \]
\[ m(\ce{Fe}) = n(\ce{Fe}) \times M(\ce{Fe}) = 0,05 \times 55,8 = \SI{2,79}{g}. \]
La lame de fer \textbf{diminue} de \SI{2,79}{g} (le fer passe en solution sous forme d'ions $\ce{Fe^2+}$).
\end{enumerate}
\end{corrige}

\begin{attention}[Point de vigilance]
On demande la variation de masse de la \emph{lame de fer} uniquement : ne pas y soustraire la masse du cuivre déposé (qui serait $0,05 \times 63,5 = \SI{3,18}{g}$), sauf si le cuivre reste accroché à la lame — l'énoncé considère ici la lame seule.
\end{attention}

\begin{corrige}[Exercice 6 — Classement des métaux et action des acides]
\begin{enumerate}
\item \textbf{Classement par pouvoir réducteur croissant :}
\[ \ce{Ag} < \ce{Cu} < \ce{Fe} < \ce{Zn} < \ce{Al} \]
\textbf{Justification :} le pouvoir réducteur d'un métal est d'autant plus \textbf{grand} que le potentiel standard de son couple est \textbf{bas}. On range donc les $E^\circ$ par valeurs décroissantes : $+0,80\ (\ce{Ag}) > +0,34\ (\ce{Cu}) > -0,44\ (\ce{Fe}) > -0,76\ (\ce{Zn}) > -1,66\ (\ce{Al})$, et on lit le classement des réducteurs en sens inverse.
\item \textbf{Métaux attaqués par l'acide chlorhydrique :} \ce{Al}, \ce{Zn} et \ce{Fe}.

\textbf{Justification (règle du gamma) :} l'oxydant mis en jeu est $\ce{H3O+}$, avec $E^\circ(\ce{H3O+}/\ce{H2}) = \SI{0,00}{V}$. Pour que le dégagement de dihydrogène soit spontané, le métal doit être un réducteur plus fort que $\ce{H2}$, c'est-à-dire appartenir à un couple de potentiel $E^\circ < \SI{0,00}{V}$. C'est le cas de \ce{Al} ($-1,66$ V), \ce{Zn} ($-0,76$ V) et \ce{Fe} ($-0,44$ V). Le cuivre ($+0,34$ V) et l'argent ($+0,80$ V) ont des potentiels supérieurs à zéro : ils ne sont \textbf{pas attaqués} (ce sont des métaux « nobles »).
\item \textbf{Équation-bilan zinc + acide chlorhydrique :}
\begin{align*}
\text{Oxydation :}\quad & \ce{Zn(s) -> Zn^2+(aq) + 2e-} \\
\text{Réduction :}\quad & \ce{2H3O+(aq) + 2e- -> H2(g) + 2H2O(l)}
\end{align*}
\[ \ce{Zn(s) + 2H3O+(aq) -> Zn^2+(aq) + H2(g) + 2H2O(l)} \]
\end{enumerate}
\end{corrige}

\begin{corrige}[Exercice 7 — Bilan quantitatif d'une pile Daniell]
\begin{enumerate}
\item \textbf{Équations de fonctionnement :}
\begin{align*}
\text{Anode (pôle } - \text{), oxydation :}\quad & \ce{Zn(s) -> Zn^2+(aq) + 2e-} \\
\text{Cathode (pôle } + \text{), réduction :}\quad & \ce{Cu^2+(aq) + 2e- -> Cu(s)}
\end{align*}
Équation-bilan : $\ce{Zn(s) + Cu^2+(aq) -> Zn^2+(aq) + Cu(s)}$

Tableau d'avancement (en moles) :
\begin{center}
\begin{tabularx}{\linewidth}{|l|X|X|X|X|}
\hline
\rowcolor{popblueL} Équation & \ce{Zn} & \ce{Cu^2+} & \ce{Zn^2+} & \ce{Cu} \\
\hline
Avancement $x$ & $n_0 - x$ & $n'_0 - x$ & $n''_0 + x$ & $x$ \\
\hline
Électrons échangés & \multicolumn{4}{l|}{$n(e^-) = 2x$} \\
\hline
\end{tabularx}
\end{center}

\item \textbf{Quantité de zinc consommée :}
\[ n(\ce{Zn}) = \frac{m(\ce{Zn})}{M(\ce{Zn})} = \frac{1,31}{65,4} = \SI{2,00e-2}{mol}. \]
\textbf{Quantité d'électricité débitée :} chaque atome de zinc libère 2 électrons, donc $n(e^-) = 2\,n(\ce{Zn}) = \SI{4,00e-2}{mol}$ :
\[ Q = n(e^-) \times F = 4,00\times 10^{-2} \times 96500 = \SI{3860}{C}. \]
\item \textbf{Masse de cuivre déposée :} d'après l'équation-bilan, $n(\ce{Cu}) = n(\ce{Zn}) = \SI{2,00e-2}{mol}$ :
\[ m(\ce{Cu}) = n(\ce{Cu}) \times M(\ce{Cu}) = 2,00\times 10^{-2} \times 63,5 = \SI{1,27}{g}. \]
\item \textbf{Variation totale de la masse des électrodes :}
\[ \Delta m = \underbrace{+1,27}_{\text{cathode}} \underbrace{-1,31}_{\text{anode}} = \SI{-0,04}{g}. \]
La masse totale des électrodes diminue de \SI{0,04}{g} : la différence provient des masses molaires différentes du zinc (qui part en solution) et du cuivre (qui se dépose). La matière n'a pas disparu : les ions $\ce{Zn^2+}$ sont passés dans la solution.
\end{enumerate}
\end{corrige}

\begin{corrige}[Exercice 8 — Micropiles et corrosion]
\begin{enumerate}
\item \textbf{Explication par les micropiles.} La goutte d'eau salée est un électrolyte conducteur posé sur le fer : toutes les conditions d'une pile sont réunies. Mais la concentration en dioxygène dissous n'est \textbf{pas uniforme} dans la goutte :
\begin{itemize}
\item \textbf{Au centre} de la goutte, loin de l'interface air/eau, le dioxygène est rare : le fer y joue le rôle d'\textbf{anode} et s'oxyde : $\ce{Fe -> Fe^2+ + 2e-}$. C'est donc au centre que le fer se corrode.
\item \textbf{En périphérie} de la goutte, au contact direct de l'air, le dioxygène est abondant : il s'y réduit sur le fer qui joue le rôle de \textbf{cathode} : $\ce{O2 + 2H2O + 4e- -> 4HO-}$.
\end{itemize}
Les électrons circulent dans le métal du centre vers la périphérie ; les ions $\ce{Fe^2+}$ et $\ce{HO-}$ migrent dans la goutte, se rencontrent et forment $\ce{Fe(OH)2}$, qui s'oxyde ensuite en rouille $\ce{Fe2O3}\cdot n\,\ce{H2O}$, souvent en couronne autour du centre.
\item \textbf{Demi-équations :}
\begin{align*}
\text{Oxydation (anode) :}\quad & \ce{Fe(s) -> Fe^2+(aq) + 2e-} \\
\text{Réduction (cathode) :}\quad & \ce{O2(g) + 2H2O(l) + 4e- -> 4HO-(aq)}
\end{align*}
En combinant ($2\times$ la première $+$ la seconde) : $\ce{2Fe + O2 + 2H2O -> 2Fe(OH)2}$.
\item \textbf{Trois facteurs accélérateurs et parades :}
\begin{center}
\begin{tabularx}{\linewidth}{|X|X|}
\hline
\rowcolor{popblueL} \textbf{Facteur accélérateur} & \textbf{Parade pratique} \\
\hline
Humidité persistante (climat tropical) & Séchage, ventilation, stockage à l'abri de la pluie \\
\hline
Présence de sels (air marin de la côte : Douala, Kribi) qui augmente la conductivité de l'électrolyte & Rinçage à l'eau douce, revêtement étanche (peinture, graisse) \\
\hline
Contact avec un métal moins réducteur (cuivre, étain) créant une pile galvanique & Éviter les assemblages bimétalliques ou les isoler électriquement \\
\hline
\end{tabularx}
\end{center}
\end{enumerate}
\end{corrige}

\courssub{Je me prépare au Bac}

\begin{corrige}[Exercice 9 — Type Bac : la pile argent–plomb]
\begin{enumerate}
\item \textbf{Réaction naturelle (règle du gamma).} $E^\circ(\ce{Ag+}/\ce{Ag}) = \SI{+0,80}{V} > E^\circ(\ce{Pb^2+}/\ce{Pb}) = \SI{-0,13}{V}$ : l'oxydant le plus fort est $\ce{Ag+}$, le réducteur le plus fort est $\ce{Pb}$.
\begin{align*}
\text{Oxydation :}\quad & \ce{Pb(s) -> Pb^2+(aq) + 2e-} \\
\text{Réduction :}\quad & \ce{Ag+(aq) + e- -> Ag(s)} \quad (\times 2)
\end{align*}
Équation-bilan : $\ce{Pb(s) + 2Ag+(aq) -> Pb^2+(aq) + 2Ag(s)}$
\item \textbf{Polarités :} la lame de \textbf{plomb} est l'\textbf{anode}, pôle \textbf{négatif} $(-)$ ; la lame d'\textbf{argent} est la \textbf{cathode}, pôle \textbf{positif} $(+)$.
\item \textbf{Représentation conventionnelle} (anode à gauche) :
\[ \ce{(-) Pb(s) | Pb^2+(aq) (\SI{0,10}{mol.L^{-1}}) || Ag+(aq) (\SI{0,10}{mol.L^{-1}}) | Ag(s) (+)} \]
\item \textbf{f.é.m. standard :}
\[ E^\circ_{\text{pile}} = E^\circ(\ce{Ag+}/\ce{Ag}) - E^\circ(\ce{Pb^2+}/\ce{Pb}) = 0,80 - (-0,13) = \SI{+0,93}{V}. \]
\item \textbf{Schéma fonctionnel annoté :}
\end{enumerate}
\textbf{Barème indicatif (sur 6 pts) :} règle du gamma et équation-bilan (1,5 pt) ; polarités (0,5 pt) ; représentation conventionnelle (0,5 pt) ; f.é.m. (1 pt) ; schéma annoté avec sens des porteurs de charge (1 pt) ; variations de masse (1,5 pt).
\end{corrige}

\begin{popfigure}
\begin{tikzpicture}[scale=1, every node/.style={font=\small}]
% Bécher gauche (Pb)
\draw[thick, popdark] (-4.2,0) -- (-4.2,-2.4) -- (-1.8,-2.4) -- (-1.8,0);
\draw[fill=popblue!15, draw=none] (-4.15,-0.4) rectangle (-1.85,-2.35);
\node at (-3,-1.4) {$\ce{Pb(NO3)2}$};
\node[below] at (-3,-2.4) {$\ce{Pb^2+}$ \SI{0,10}{mol.L^{-1}}};
% Lame Pb
\draw[fill=popmuted!60, draw=popdark] (-3.15,0.6) rectangle (-2.85,-1.8);
\node[above, popdark] at (-3,0.65) {\textbf{Pb} $(-)$ anode};
% Bécher droit (Ag)
\draw[thick, popdark] (1.8,0) -- (1.8,-2.4) -- (4.2,-2.4) -- (4.2,0);
\draw[fill=popgreen!12, draw=none] (1.85,-0.4) rectangle (4.15,-2.35);
\node at (3,-1.4) {$\ce{AgNO3}$};
\node[below] at (3,-2.4) {$\ce{Ag+}$ \SI{0,10}{mol.L^{-1}}};
% Lame Ag
\draw[fill=popgold!50, draw=popdark] (2.85,0.6) rectangle (3.15,-1.8);
\node[above, popdark] at (3,0.65) {\textbf{Ag} $(+)$ cathode};
% Pont salin
\draw[line width=4pt, poppurple!60] (-2.1,-0.5) .. controls (-1,-1.6) and (1,-1.6) .. (2.1,-0.5);
\node[poppurple, below] at (0,-1.7) {pont salin $\ce{K+} + \ce{NO3-}$};
\draw[->, thick, poppurple] (-0.6,-1.35) -- (-1.4,-0.95) node[midway, below left=-2pt] {\scriptsize $\ce{NO3-}$};
\draw[->, thick, poppurple] (0.6,-1.35) -- (1.4,-0.95) node[midway, below right=-2pt] {\scriptsize $\ce{K+}$};
% Circuit extérieur
\draw[thick, popdark] (-3,0.6) -- (-3,1.6) -- (3,1.6) -- (3,0.6);
\draw[->, very thick, popink] (-1.6,1.6) -- (-0.4,1.6);
\node[popink, above] at (-1,1.6) {$e^-$};
\node[popdark, above] at (1.6,1.6) {circuit extérieur};
% Réactions
\node[align=center, poporange!80!black, below] at (-3,-3.1) {$\ce{Pb -> Pb^2+ + 2e-}$};
\node[align=center, popblue!80!black, below] at (3,-3.1) {$\ce{Ag+ + e- -> Ag}$};
\end{tikzpicture}
\legende{Pile argent–plomb : électrons de Pb vers Ag dans le circuit extérieur ; dans le pont salin, les cations $\ce{K+}$ migrent vers la cathode (Ag) et les anions $\ce{NO3-}$ vers l'anode (Pb).}
\end{popfigure}

\begin{corrige}[Exercice 9 — Suite : variations de masse]
\begin{enumerate}
\setcounter{enumi}{5}
\item \textbf{Variations de masse.} Quantité initiale d'ions argent :
\[ n_i(\ce{Ag+}) = C \times V = 0,10 \times 0,100 = \SI{1,00e-2}{mol}. \]
Concentration divisée par dix $\Rightarrow$ quantité restante $n_f = \SI{1,00e-3}{mol}$, donc quantité consommée :
\[ \Delta n(\ce{Ag+}) = 1,00\times 10^{-2} - 1,00\times 10^{-3} = \SI{9,00e-3}{mol}. \]
Tableau d'avancement :
\begin{center}
\begin{tabularx}{\linewidth}{|l|X|X|X|X|}
\hline
\rowcolor{popblueL} Équation & \ce{Pb} & \ce{2Ag+} & \ce{Pb^2+} & \ce{2Ag} \\
\hline
État initial (mol) & excès & $1,00\times 10^{-2}$ & $1,00\times 10^{-2}$ & excès \\
\hline
État final (mol) & $-x$ consommé & $1,00\times 10^{-3}$ & $+x$ formé & $+2x$ formé \\
\hline
\end{tabularx}
\end{center}
Avec $2x = \SI{9,00e-3}{mol}$, soit $x = \SI{4,50e-3}{mol}$.
\begin{itemize}
\item \textbf{Lame d'argent (cathode) :} $n(\ce{Ag})_{\text{déposé}} = 2x = \SI{9,00e-3}{mol}$ :
\[ \Delta m(\ce{Ag}) = 9,00\times 10^{-3} \times 107,9 = \SI{+0,97}{g} \quad \textbf{(augmentation)}. \]
\item \textbf{Lame de plomb (anode) :} $n(\ce{Pb})_{\text{dissous}} = x = \SI{4,50e-3}{mol}$ :
\[ \Delta m(\ce{Pb}) = -\,4,50\times 10^{-3} \times 207,2 = \SI{-0,93}{g} \quad \textbf{(diminution)}. \]
\end{itemize}
\end{enumerate}
\end{corrige}

\begin{attention}[Points de vigilance]
La demi-équation de l'argent échange \textbf{un seul} électron : il faut la multiplier par 2 avant d'écrire le bilan. Dans le tableau d'avancement, ne pas oublier le coefficient 2 devant $\ce{Ag+}$ : $n(\ce{Pb})_{\text{consommé}} = \tfrac{1}{2}\,n(\ce{Ag+})_{\text{consommé}}$. Une f.é.m. est toujours \textbf{positive}.
\end{attention}

\begin{corrige}[Exercice 10 — Type Bac : tôle galvanisée et boîte de conserve]
\begin{enumerate}
\item \textbf{Classement par pouvoir réducteur croissant} (du $E^\circ$ le plus élevé au plus bas) :
\[ \ce{Sn}\ (-0,14\ \text{V}) < \ce{Fe}\ (-0,44\ \text{V}) < \ce{Zn}\ (-0,76\ \text{V}). \]
\item \textbf{Tôle galvanisée rayée.} Le zinc et le fer nus sont en contact dans l'eau de pluie : une micropile $\ce{Zn}/\ce{Fe}$ se forme. Comme $E^\circ(\ce{Zn^2+}/\ce{Zn}) < E^\circ(\ce{Fe^2+}/\ce{Fe})$, le zinc est le réducteur le plus fort : c'est lui qui s'oxyde (anode) :
\begin{align*}
\text{Anode (zinc) :}\quad & \ce{Zn(s) -> Zn^2+(aq) + 2e-} \\
\text{Cathode (fer) :}\quad & \ce{O2(g) + 2H2O(l) + 4e- -> 4HO-(aq)}
\end{align*}
Le fer reçoit les électrons et sert seulement de support à la réduction du dioxygène : il \textbf{n'est pas oxydé}. \textbf{Conclusion :} même rayée, la tôle galvanisée reste protégée tant qu'il reste du zinc au voisinage : le zinc se consume à la place du fer.
\item \textbf{Boîte de conserve étamée rayée.} Cette fois $E^\circ(\ce{Fe^2+}/\ce{Fe}) = \SI{-0,44}{V} < E^\circ(\ce{Sn^2+}/\ce{Sn}) = \SI{-0,14}{V}$ : le \textbf{fer} est le réducteur le plus fort. Il devient l'anode de la micropile :
\begin{align*}
\text{Anode (fer) :}\quad & \ce{Fe(s) -> Fe^2+(aq) + 2e-} \\
\text{Cathode (étain) :}\quad & \ce{O2(g) + 2H2O(l) + 4e- -> 4HO-(aq)}
\end{align*}
La corrosion est d'autant plus rapide que la surface anodique (le fer nu au fond de la rayure) est \textbf{très petite} devant la surface cathodique (tout l'étain environnant) : la densité de courant anodique est énorme, d'où une piqûre profonde et une perforation rapide de la boîte. L'étain ne protège le fer que tant que le revêtement est \textbf{parfaitement continu}.
\item \textbf{Nom de la protection :} \cle{protection cathodique} par \textbf{anode sacrificielle} (le zinc est l'anode qui se sacrifie).
\item \textbf{Deux autres méthodes :}
\begin{itemize}
\item \textbf{Revêtement barrière} (peinture, vernis, émail, plastique) : on isole physiquement le fer de l'eau et du dioxygène, ce qui empêche toute micropile de fonctionner. \emph{Exemple :} peinture antirouille sur les charpentes métalliques.
\item \textbf{Protection cathodique à courant imposé :} un générateur extérieur impose au fer un potentiel de cathode ; l'anode est alors une électrode inerte (graphite, oxydes mixtes). \emph{Exemple :} protection des pipelines enterrés et des coques de navires.
\end{itemize}
\end{enumerate}
\textbf{Barème indicatif (sur 5 pts) :} classement (0,5 pt) ; micropile Zn/Fe avec demi-équations et conclusion (1,5 pt) ; micropile Fe/Sn avec explication de la corrosion accélérée (1,5 pt) ; nom de la protection (0,5 pt) ; deux méthodes justifiées (1 pt).
\end{corrige}

\begin{attention}[Point de vigilance]
Ne pas écrire que « l'étain protège le fer parce qu'il l'enrobe » sans nuance : c'est vrai seulement si le revêtement est \emph{intact}. Dès la première rayure, l'étain \textbf{accélère} la corrosion du fer, contrairement au zinc. Cette différence est une question classique du Bac TI.
\end{attention}

\begin{corrige}[Exercice 11 — Type Bac : protection cathodique d'un pipeline]
\begin{enumerate}
\item \textbf{Méthode :} \cle{protection cathodique} du pipeline. Les blocs de magnésium sont appelés \textbf{anodes sacrificielles} (ou anodes galvaniques).
\item \textbf{Rôles dans la micropile (règle du gamma).} $E^\circ(\ce{Mg^2+}/\ce{Mg}) = \SI{-2,37}{V} < E^\circ(\ce{Fe^2+}/\ce{Fe}) = \SI{-0,44}{V}$ : le magnésium est un réducteur beaucoup plus fort que le fer.
\begin{align*}
\text{Anode (magnésium, pôle } - \text{) :}\quad & \ce{Mg(s) -> Mg^2+(aq) + 2e-} \\
\text{Cathode (fer du pipeline, pôle } + \text{) :}\quad & \ce{O2(g) + 2H2O(l) + 4e- -> 4HO-(aq)}
\end{align*}
Le fer, devenu cathode, ne peut plus s'oxyder : il est protégé.
\item \textbf{Pourquoi le magnésium et pas le cuivre ?} $E^\circ(\ce{Cu^2+}/\ce{Cu}) = \SI{+0,34}{V} > E^\circ(\ce{Fe^2+}/\ce{Fe})$ : le cuivre est un réducteur \textbf{plus faible} que le fer. Dans un couple Fe/Cu, c'est le \textbf{fer} qui deviendrait l'anode et s'oxyderait : le cuivre accélérerait la corrosion du pipeline au lieu de le protéger. L'anode sacrificielle doit toujours être un métal de potentiel \textbf{inférieur} à celui du métal à protéger.
\item \textbf{Remplacement périodique :} en jouant le rôle d'anode, le magnésium s'oxyde et se dissout progressivement ($\ce{Mg -> Mg^2+ + 2e-}$) : sa masse diminue. Lorsque le bloc est entièrement consommé, la protection cesse et le pipeline redevient vulnérable ; il faut donc remplacer les anodes avant leur épuisement.
\item \textbf{Calculs.} Quantité de magnésium consommée :
\[ n(\ce{Mg}) = \frac{m}{M(\ce{Mg})} = \frac{4,86\times 10^{3}}{24,3} = \SI{200}{mol}. \]
Chaque mole de magnésium fournit 2 moles d'électrons, donc $n(e^-) = \SI{400}{mol}$ :
\[ Q = n(e^-) \times F = 400 \times 96500 = \SI{3,86e7}{C}. \]
Sans protection, la même quantité d'électricité correspondrait à l'oxydation du fer ($\ce{Fe -> Fe^2+ + 2e-}$, soit 2 mol d'électrons par mole de fer) :
\[ n(\ce{Fe}) = \frac{n(e^-)}{2} = \frac{400}{2} = \SI{200}{mol} \quad\Rightarrow\quad m(\ce{Fe}) = 200 \times 55,8 = \SI{1,116e4}{g} \approx \SI{11,2}{kg}. \]
\textbf{Conclusion :} un bloc de seulement \SI{4,86}{kg} de magnésium a évité la corrosion d'environ \SI{11,2}{kg} de fer. Le sacrifice est rentable : le rapport des masses protégées suit le rapport des masses molaires $M(\ce{Fe})/M(\ce{Mg}) \approx 2,3$, car les deux métaux échangent le même nombre d'électrons (2) par atome.
\end{enumerate}
\textbf{Barème indicatif (sur 6 pts) :} méthode et nom des blocs (0,5 pt) ; rôles anode/cathode avec demi-équations (1,5 pt) ; justification Mg vs Cu (1 pt) ; usure et remplacement (0,5 pt) ; calcul de $Q$ (1,5 pt) ; masse de fer économisée (1 pt).
\end{corrige}

\begin{attention}[Points de vigilance]
Convertir la masse en grammes avant de diviser par la masse molaire ($\SI{4,86}{kg} = \SI{4860}{g}$). Ne pas oublier le facteur 2 lié aux deux électrons échangés par atome de magnésium. La quantité d'électricité $Q$ s'exprime en coulombs (C), jamais en ampères.
\end{attention}

\newpage

\coursec{Fiche bilan}

\courssub{Carte mentale de la leçon}

\begin{popfigure}
\begin{tikzpicture}[
  central/.style={circle, draw=popdark, fill=popblue, text=white, font=\bfseries\small, minimum size=3.2cm, align=center},
  branche/.style={rounded corners=4pt, draw=popdark, font=\scriptsize, align=center, minimum width=3.0cm, inner sep=3pt}]
  \node[central] (C) {Piles\\électrochimiques\\et corrosion};
  \node[branche, fill=popblueL] (B1) at (-5.4,2.6) {Classification\\électrochimique\\Règle du gamma};
  \node[branche, fill=poporangeL] (B2) at (0,3.6) {Demi-piles\\$M^{n+}/M$\\Électrode de\\référence ESH};
  \node[branche, fill=popgreenL] (B3) at (5.4,2.6) {Pile Daniell\\$\ominus$ Zn $|$ Cu $\oplus$\\$E = E_{\oplus}-E_{\ominus}$};
  \node[branche, fill=poppinkL] (B4) at (5.4,-2.6) {Fonctionnement\\Anode : oxydation\\Cathode : réduction};
  \node[branche, fill=popgoldL] (B5) at (0,-3.6) {Prévision des piles\\Polarité, f.é.m,\\équations aux électrodes};
  \node[branche, fill=poppurpleL] (B6) at (-5.4,-2.6) {Corrosion\\Micropiles, rouille\\Facteurs aggravants};
  \node[branche, fill=popcream] (B7) at (-8.2,0) {Protection\\Revêtement, galvanisation\\protection cathodique};
  \draw[popline, line width=1.2pt] (C) -- (B1);
  \draw[popline, line width=1.2pt] (C) -- (B2);
  \draw[popline, line width=1.2pt] (C) -- (B3);
  \draw[popline, line width=1.2pt] (C) -- (B4);
  \draw[popline, line width=1.2pt] (C) -- (B5);
  \draw[popline, line width=1.2pt] (C) -- (B6);
  \draw[popline, line width=1.2pt] (C) -- (B7);
\end{tikzpicture}
\legende{Carte mentale : les sept idées-forces de la leçon 10.}
\end{popfigure}

\courssub{L'essentiel à retenir}

\begin{aretenir}[Synthèse des connaissances essentielles]
\textbf{Définitions clés}
\begin{itemize}
  \item \cle{Oxydant} : espèce chimique capable de capter un ou plusieurs électrons ; \cle{réducteur} : espèce chimique capable de céder un ou plusieurs électrons. Un \cle{couple redox} s'écrit $\text{Ox}/\text{Red}$ et sa demi-équation d'oxydoréduction est $\text{Ox} + n\,e^- \rightleftharpoons \text{Red}$.
  \item \cle{Demi-pile $M^{n+}/M$} : constituée d'une lame métallique $M$ plongée dans une solution contenant ses ions $M^{n+}$ à une concentration connue.
  \item \cle{Pile électrochimique} : association de deux demi-piles reliées par un circuit externe (fils conducteurs) et un \cle{pont salin} (ou interface) ; elle transforme l'énergie chimique d'une réaction d'oxydoréduction naturelle en énergie électrique utilisable.
  \item \cle{Force électromotrice (f.é.m)} : $E = E_{\oplus} - E_{\ominus}$ (toujours positive pour une pile fonctionnant spontanément).
  \item \cle{Corrosion} : dégradation irréversible d'un métal par oxydation sous l'action de son milieu (dioxygène $\ce{O2}$, eau $\ce{H2O}$, ions $\ce{H+}$, sels), due à la formation de \cle{micropiles} locales à la surface du métal.
\end{itemize}

\textbf{Équations-types à connaître par cœur (pile Daniell $\ce{Zn}|\ce{Cu}$)}
\begin{itemize}
  \item Anode (pôle $\ominus$, zinc, oxydation) : \ce{Zn -> Zn^2+ + 2e-}
  \item Cathode (pôle $\oplus$, cuivre, réduction) : \ce{Cu^2+ + 2e- -> Cu}
  \item Bilan de fonctionnement (réaction spontanée) : \ce{Zn + Cu^2+ -> Zn^2+ + Cu}
  \item Représentation conventionnelle (pôle $\ominus$ à gauche, $\oplus$ à droite) : 
  $\ominus\ \ce{Zn}\ |\ \ce{Zn^2+}\ (aq)\ ||\ \ce{Cu^2+}\ (aq)\ |\ \ce{Cu}\ \oplus$
  \item Électrode de référence standard : \cle{Électrode Standard à Hydrogène (ESH)} : 
  \ce{Pt} $|$ \ce{H2} (g, \SI{1}{bar}) $|$ \ce{H3O+} (\SI{1}{mol.L^{-1}}), avec $E^\circ(\ce{H3O+}/\ce{H2}) = \SI{0,00}{V}$.
\end{itemize}

\textbf{Formules et données utiles}
\begin{center}
\begin{tabularx}{\linewidth}{|l|X|}
\hline
\rowcolor{popblueL} \textbf{Grandeur} & \textbf{Expression / valeur} \\
\hline
f.é.m d'une pile & $E = E_{\oplus} - E_{\ominus} > 0$ \\
\hline
f.é.m pile Daniell (conditions standards) & $E^\circ = E^\circ_{\ce{Cu^2+}/\ce{Cu}} - E^\circ_{\ce{Zn^2+}/\ce{Zn}} = 0{,}34 - (-0{,}76) = \SI{1,10}{V}$ \\
\hline
Règle du gamma (prévision réaction naturelle) & Le réducteur le plus fort (couple bas) réagit avec l'oxydant le plus fort (couple haut) \\
\hline
Quantité d'électricité & $Q = n_e \times \mathcal{F}$, $\mathcal{F} = \SI{96500}{C.mol^{-1}}$ (constante de Faraday) \\
\hline
Potentiels standards usuels ($E^\circ$ en V/ESH) & 
$\ce{Cu^2+}/\ce{Cu}$ : $+0{,}34$ ; 
$\ce{Fe^2+}/\ce{Fe}$ : $-0{,}44$ ; 
$\ce{Zn^2+}/\ce{Zn}$ : $-0{,}76$ ; 
$\ce{Al^3+}/\ce{Al}$ : $-1{,}66$ ; 
$\ce{Mg^2+}/\ce{Mg}$ : $-2{,}37$ \\
\hline
\end{tabularx}
\end{center}

\textbf{Méthodes express}
\begin{itemize}
  \item \cle{Polarité d'une pile} : le pôle $\ominus$ (anode, oxydation) correspond au métal du couple dont le potentiel standard est le \emph{plus bas} (le plus fort réducteur) ; le pôle $\oplus$ (cathode, réduction) correspond au couple au potentiel le \emph{plus haut} (le plus fort oxydant).
  \item \cle{Prévision d'une réaction d'oxydoréduction} : placer les deux couples sur l'axe des potentiels standards croissants, tracer le \cle{gamma ($\gamma$)} reliant l'oxydant du couple haut au réducteur du couple bas ; la réaction naturelle est celle du gamma.
  \item \cle{Sens des courants} : les électrons circulent dans le circuit extérieur du pôle $\ominus$ (anode) vers le pôle $\oplus$ (cathode) ; le courant conventionnel va de $\oplus$ vers $\ominus$. Dans la solution, les cations migrent vers la cathode, les anions vers l'anode.
  \item \cle{Rôle du pont salin} : assure la neutralité électrique de chaque demi-pile par migration ionique (ex. $\ce{K+}$, $\ce{NO3-}$) et ferme le circuit électrique, sans mélanger les électrolytes.
  \item \cle{Protection contre la corrosion} : 
  \begin{enumerate}
    \item \textbf{Revêtement passif} : peinture, vernis, graisse, plastique, étamage (isolent le métal de l'air et de l'eau).
    \item \textbf{Galvanisation (zincage)} : revêtement de zinc (plus réducteur que le fer) qui protège par anode sacrificielle même si rayé.
    \item \textbf{Protection cathodique} : liaison du métal à protéger (cathode) à une anode sacrificielle (\ce{Zn}, \ce{Mg}, \ce{Al}) ou courant imposé.
    \item \textbf{Alliages inoxydables} : ajout de chrome (\ce{Cr}) formant une couche passive $\ce{Cr2O3}$ imperméable.
  \end{enumerate}
\end{itemize}
\end{aretenir}

\begin{aretenir}[Checklist avant le Bac]
\begin{itemize}
  \item[$\square$] Je sais définir oxydant, réducteur, couple redox et écrire une demi-équation équilibrée (atomes et charges).
  \item[$\square$] Je sais exploiter la classification électrochimique pour classer les oxydants (du plus fort au plus faible) et les réducteurs (du plus faible au plus fort).
  \item[$\square$] Je sais appliquer la règle du gamma pour prévoir le sens naturel d'une réaction d'oxydoréduction entre deux couples.
  \item[$\square$] Je sais définir une demi-pile $M^{n+}/M$ et citer l'électrode de référence (ESH, $E^\circ = \SI{0,00}{V}$).
  \item[$\square$] Je sais schématiser la pile Daniell (électrodes, solutions, pont salin, voltmètre, sens des électrons et du courant conventionnel).
  \item[$\square$] Je sais déterminer les pôles d'une pile à partir des potentiels standards des couples mis en jeu.
  \item[$\square$] Je sais calculer la f.é.m : $E = E_{\oplus} - E_{\ominus}$ (ex. Daniell : \SI{1,10}{V}).
  \item[$\square$] Je sais écrire les équations aux électrodes (anode : oxydation, cathode : réduction) et l'équation-bilan de fonctionnement.
  \item[$\square$] Je sais écrire la représentation conventionnelle d'une pile ($\ominus$ à gauche, $\oplus$ à droite, barre simple pour changement de phase, double barre pour pont salin).
  \item[$\square$] Je sais expliquer la corrosion du fer par la formation de micropiles (zones anodiques : $\ce{Fe -> Fe^2+ + 2e-}$ ; zones cathodiques : $\ce{O2 + 2H2O + 4e- -> 4OH-}$) et citer les facteurs aggravants (humidité, sel, acidité, température, contact métaux différents).
  \item[$\square$] Je sais décrire au moins trois méthodes de protection : revêtement passif, galvanisation, protection cathodique (anode sacrificielle ou courant imposé), alliages inoxydables.
\end{itemize}
\end{aretenir}

\findecours

\end{document}
